\section{The local lattice link}\label{sec:pf:linkG}

\providecommand{\linkGV}{\mathcal V}
\providecommand{\linkGH}{\mathcal H}
\providecommand{\linkGR}{\mathcal R}
\providecommand{\linkGNs}{N^{\mathrm{src}}}
\providecommand{\linkGNt}{N^{\mathrm{tgt}}}
\providecommand{\linkGLP}{\mathsf{LP}}
\providecommand{\linkGLPat}{\mathsf{LP}_{\mathrm{at}}}
\providecommand{\linkGGR}{\mathsf{GR}}
\providecommand{\linkGLL}{\mathsf{LL}}
\providecommand{\linkGEF}{\mathsf{EF}}
\providecommand{\linkGST}{\mathsf{ST}}
\providecommand{\linkGES}{\mathsf{ES}}
\providecommand{\linkGAB}{\mathsf{AB}}
\providecommand{\linkGst}{\mathrm{st}}
\providecommand{\linkGVV}{\mathit{VV}}
\providecommand{\linkGgain}{\mathrm{gain}}
\providecommand{\linkGlen}{\ell}

This section proves \cref{thm:linkG} and \cref{thm:families:local}. The first says that the link
$\sfname{LinkG}$ of \cref{sec:families:links} is a closed property on every family class on which the group
compatibility of the extra data and the size condition $\sfname{Adm2F}$ are closed. The second is a purely local
statement about lattices in tensor products of dual Weyl modules over a $p$-adic ring of integers: from every
lattice one reaches, along ``Hecke edges'', vectors of the weight lattices which are highly divisible on all
competing edges and whose pairings with the dual weight lattice have almost full rank. The local statement is the
deepest ingredient of $\sfname{LinkG}$; the rest of the argument transports it to the members of a family and
converts it into bounds for lengths of lattices of algebraic modular forms.

\paragraph{The argument.}
The proof has a local half and a global half.
\begin{enumerate}[label=(\arabic*)]
\item \emph{The local lattice theorem} (\cref{sec:pf:linkG:local,sec:pf:linkG:edges,sec:pf:linkG:criterion}).
We state it in full (\cref{pf:linkG:localthm}): the definitions involved (unramified quadratic settings, lattices
and Hecke edges, the weight lattices $\linkGH_L$ and the edge lattices $\linkGR_e$, the property $\linkGLP$) are
self-contained. Its proof is organised as follows. A \emph{criterion} reduces the conclusion to $37$ explicit
statements (``clauses''), about representations of $\GL_n$ over discrete valuation rings, polynomial models of Weyl
modules, Gelfand--Tsetlin patterns, orbit combinatorics and $p$-adic numerics; none of them depends on the datum which
formally carries it. The last step of the criterion is a counting argument which we prove in full
(\cref{sec:pf:linkG:edges}): by homogeneity it suffices to treat the standard lattice; there are at least
$p^{n'r(m+1-r)}$ edges with a given target; each edge carries an \emph{edge family} of vectors with the required
divisibility and almost the expected number of elements; and a Smith normal form argument turns the edge families
into the rank statement. The construction of the edge families (the \emph{edge step}) and the verification of the
clauses occupy most of the formal proof; their architecture is given in \cref{sec:pf:linkG:criterion}, the inner
reductions of the criterion are proved in \cref{sec:pf:linkG:inner}, and the symbol space statement (clause 27) in
\cref{sec:pf:linkG:second}.
\item \emph{From local to global} (\cref{sec:pf:linkG:lengths}). At a place $w$ of the coefficient field above a
prime $p$, the lattice data of a member consist of spaces $M\subseteq M^{\mathrm{na}}$ of algebraic modular forms
(functions on tuples of lattices with values in the weight lattices) and Hecke operators $y_v$ given by sums over
Hecke edges. The local lattice theorem at the places $v$ above $p$ implies, by lemma J and the walk lemma
(\cref{pf:linkG:lemmaJ,pf:linkG:walkabs}), that up to an explicit exceptional set, products of powers of the $y_v$ map $M$
into a highly divisible sublattice; we prove in full the
module-theoretic lemma which converts such ``all but few'' statements into a bound for the length of
$M/N$ (\cref{pf:linkG:lengthpart2,pf:linkG:core}).
\item \emph{The family argument} (\cref{sec:pf:linkG:family,sec:pf:linkG:proof}). The local theorem, specialised to
the parameters of the members of a family, gives the condition $\sfname{Adm6F}$; together with eight further
admissibility conditions it gives, by the ``Kudla family statement'' (\cref{pf:linkG:stmtF}), the gain bounds and
the von Mangoldt inequality of $\sfname{LinkG}$; the length bounds of the previous step give the local length
inequalities. Finally ten closed properties are combined to prove \cref{thm:linkG}.
\end{enumerate}

\paragraph{What is proved here and what is cited.}
We prove in full: the reduction of the local lattice theorem to edge families and its numerical part
(\cref{sec:pf:linkG:edges}); the assembly of the local lattice theorem from the criterion and the clauses, one of the
clauses and a subspace count (\cref{sec:pf:linkG:criterion,sec:pf:linkG:verif}); the length lemmas, the walk lemma, lemma J,
the many-orbits statement and the length bound for lattices of modular forms (\cref{sec:pf:linkG:lengths}); the
derivation of the inequalities of $\sfname{LinkG}$ from the family statements, and the closedness arguments
(\cref{sec:pf:linkG:family,sec:pf:linkG:proof}). At the level of their main steps, with the principal estimates, we prove
the inner reductions of the criterion (\cref{sec:pf:linkG:inner}), the second part of the tool package and the space
statement (\cref{sec:pf:linkG:second}), the closed form of the Kudla family statement (\cref{pf:linkG:stmtF}) and its six
components (\cref{sec:pf:linkG:famcomp}), and the membership part of ``Proposition 5'' (\cref{pf:linkG:memb}); with all
estimates and counting arguments in full we prove the good path weight bound (\cref{sec:pf:linkG:goodpaths}) and the choice
part of ``Proposition 5'' (\cref{pf:linkG:pick}). The
following inputs, all proved in the formal development \cite{LeanSource2026}, are cited with their statements: the deeper
ingredients of these reductions (among them the operator words bound, the realisation of test polynomials, the generator
congruence modulo the uniformiser and the decomposition into products over segments in the membership part), the verification of
most clauses (representation theory of $\GL_n$ over discrete valuation rings, Fock-model lattice lemmas and
Gelfand--Tsetlin combinatorics, about $16\,000$ lines), and some local facts on the slot lattices. Peculiarities of the
formal statements are collected in \cref{sec:pf:linkG:remarks}.

% ---------------------------------------------------------------------------
\subsection{Families, local lattice data and the link}\label{sec:pf:linkG:link}

We use the notation of \cref{sec:families}. A family $\mathcal F\in\sfname{FamP}(\eps)$ has members
$(D,L)=(D_{x,i},L_{x,i})$ for $x>x_1$ and $i\in I_x$, where $D$ is a datum and $L$ its size data, with scale
quantities $a,b,b',x,\delta_1=x^{-\eps/8}$, $X'$, $X''(n)$, $p_0(k)$ and $w''(n)$ (functions of $x$ and $\eps$). Besides
the extra data $\mathrm{ex}(x,i)$ of \cref{def:families:closed}, a family carries, for every member, a
\emph{parameter instance} $X=G_x(x,i)$ and a \emph{lattice instance} $Y=G_y(x,i)$, and a list $G_K$ of
non-negative \emph{link constants} $\Lambda_K,\rho_K,k_{II},t_{II},D_{II},s_{II},k'_{II},\dots$, the same for
all members.

\begin{definition}[Closed properties]\label{pf:linkG:closed}
We make \cref{def:families:closed} explicit. The constants of a family $\mathcal F$ are arranged in four lists:
\emph{outer} constants $c_{\mathrm{out}}$, \emph{middle} constants $c_{\mathrm{mid}}(C_H)$ (which contain $C_H$ and
the link constant $\Lambda_K$), \emph{last} constants $c_{\mathrm{last}}$ (which contain the link constants
$\rho_K,k_{II},t_{II},D_{II},s_{II},k'_{II},D_I,k_I,C_N$ and two thresholds for $\kappa$), and
\emph{preliminary} constants. A \emph{staged threshold} is a quadruple
$K=(K_{\mathrm{pre}},K_t,K_{\mathrm{mid}},K_{\mathrm{last}})$ of a real number and three functions from lists of
reals to reals; $\mathcal F$ is \emph{large} for $(K,C_H)$ if every preliminary constant is at least
$K_{\mathrm{pre}}$, every middle constant is at least $K_{\mathrm{mid}}(c_{\mathrm{out}})$, and every last constant
is at least $K_{\mathrm{last}}(c_{\mathrm{out}}\cup c_{\mathrm{mid}}(C_H))$. A property $P(\mathcal F,C_H)$ is
\emph{closed} on a class $\mathcal C$ if there is $\eps_1>0$ such that for every $\eps\in(0,\eps_1]$ there is a
staged threshold $K$ with $P(\mathcal F,C_H)$ for all $\mathcal F\in\mathcal C\cap\sfname{FamP}(\eps)$ and all
$C_H$ for which $\mathcal F$ is large.
\formal{Formalism.MainNotions.Closed}
\end{definition}

Closed properties are stable under finite conjunctions (take the least $\eps_1$ and, stage by stage, the maximum
of the thresholds; \lref{Formalism.MainNotions.Closed.and}) and under implications valid for every family
with $\eps>0$ (\lref{Formalism.MainNotions.Closed.imp}). For a family $\mathcal F$ and a property $P(x,i)$
of its members we write $\sfname{Ev}_{\mathcal F}(P)$ (``eventually'') for: there is $x_0$ with $P(x,i)$ for all
$x\ge x_0$ and $i\in I_x$. Eventual properties are stable under finite conjunctions (take the largest $x_0$), and
$\sfname{Ev}_{\mathcal F}(x>x_1)$ holds.

\begin{definition}[Local group data]\label{pf:linkG:gatp}
Let $u$ be a word of $D$ with level prime $T_u$, and let $w$ be a finite place of the coefficient field
$\mathit{LL}$ of $D$ whose residue characteristic $p=p(w)$ is admissible for $u$ (a condition on
$p$ and $T_u$). The datum $D$ attaches to $(u,w)$ a \emph{local group datum} $\mathcal G=\mathcal G(u,w)$,
consisting of:
\begin{enumerate}[label=(\alph*)]
\item a complete discrete valuation ring $\OO$ with fraction field $K_f$ and an element $\varpi_0\in\OO$;
\item finitely many \emph{places} $v=1,\dots,n_V$, each with a Hida setting in the sense of \cref{pf:linkG:hida}
below, with rings $W_v,Q_v,\OO,K_f$, maps $\iota_v,\sigma_v,j_v$, degree $n'_v$ and the common element $\varpi_0$;
integers $m,r\ge0$ (the same for all $v$) and dominant weights $\lambda_v:\Z/n'_v\to\Z^{m+1}$;
\item a finite set of \emph{components}, each with a \emph{level}: a group $\Gamma$ with homomorphisms
$\Gamma\to\GL_{m+1}(Q_v)$ for all $v$ and a character $\chi:\Gamma\to\OO^\times$;
\item an integer $e\ge0$ (the sandwich exponent).
\end{enumerate}
A \emph{tuple of lattices} is a family $\mathbf L=(L_v)_v$ of lattices $L_v\subseteq Q_v^{m+1}$ (\cref{pf:linkG:edges}).
The weight lattices of \cref{pf:linkG:weightlat} at the places $v$ give, by tensor product over $v$, an
$\OO$-lattice $\linkGH(\mathbf L)$ (built from the dual Weyl lattices $\nabla$) and an $\OO$-lattice
$\linkGH^{W}(\mathbf L)$ (built in the same way from the Weyl lattices $\Delta$) in a common polynomial space $\mathcal X$.
For a level, $M^{\mathrm{na}}$ (resp.\ $M$) is the $\OO$-module of functions $F$ from tuples to $\mathcal X$ with
$F(\mathbf L)\in\linkGH(\mathbf L)$ (resp.\ $\linkGH^{W}(\mathbf L)$) for tuples of lattices, $F=0$ on other tuples, and
$F(\gamma\mathbf L)=\chi(\gamma)\,\gamma\cdot F(\mathbf L)$ for $\gamma\in\Gamma$, where $\Gamma$ acts on $\mathcal X$ by the
natural (determinant-twisted) action; $M^{\mathrm{na}}$ and $M$ of $\mathcal G$ are the products over the
components. We put $d=\rank_\OO M$. For a place $v$ the \emph{Hecke operator} $y_v$ is
\[
(y_vF)(\mathbf L)=p^{A_r(\lambda_v)}\sum_{L''}F(\mathbf L\text{ with }L_v\text{ replaced by }L''),
\]
the sum over the sources $L''$ of the Hecke edges $(L''\subseteq L_v)$ of type $r$ (\cref{pf:linkG:edges}), where
$A_r$ is defined in \cref{pf:linkG:weightlat}. The datum $\mathcal G$ has \emph{trivial stabilisers} if for every
level, $\gamma\mathbf L=\mathbf L$ for a tuple of lattices implies $\gamma=1$; it has \emph{finite classes} if every
level has finitely many orbits on tuples of lattices; and it satisfies the \emph{sandwich} condition if
$p^eM^{\mathrm{na}}\subseteq M\subseteq M^{\mathrm{na}}$.
\formal{Formalism.MainNotions.G.GAtP}
\end{definition}

\begin{definition}[Orbit data at inert primes]\label{pf:linkG:orbitdata}
Let $H$ be the finite abelian group of $D$ which indexes the \emph{tops}, and $TT\subseteq H$ the set of tops of
the host. For a prime $p$ the host provides a Frobenius element $h_p\in H$; the translations by $h_p$ (or $2h_p$ if
$p$ is in the \emph{signal} class) decompose $H$ into orbits, with representatives $\mathrm{reps}(p)$ chosen by
$X$. For a representative $t_0$, the set of tops on its orbit has $N_v(p,t_0)$ elements and determines a subset of
$\Z/\ord(h_p)$; $t_0$ is \emph{good} if this subset satisfies the good-orbit tile condition of
\cref{pf:linkG:orbitpar} with the parameters $r_G=\lceil8(h_G+2)/\eps\rceil+1$, $h_G=2+8\eps$, $\delta_1$, $x$ of
the scale. A prime $p$ is a \emph{candidate} if $p<X'$, $p$ lies in the inert class of the host,
$\ord(h_p)\ge x^{\eps/4}$, and $p$ avoids an explicit band of primes with a power close to $X'$ when $p\ge p_0(k)$.
With the local parameters $t_v(p,t_0)$, $D_p$ and $\kappa_v(p,t_0)$ attached to the good orbits by $X$ and the link
constants (\cref{sec:pf:linkG:family}), put $\Sigma(p)=m\sum_{t_0\text{ good}}\kappa_v(p,t_0)$. The set of \emph{served
orbits} $\linkGVV(p)$ is the set of good representatives if $p$ is a candidate and $\Sigma(p)\le x^{-\eps}$, and is
empty otherwise. Then
\begin{equation}\label{eq:pf:linkG:tauio}
\tau(p)=e_p+\sum_{t_0\in\linkGVV(p)}m\bigl(\lceil t_v(p,t_0)\rceil+D_p+1\bigr),\qquad
\iota(p)=2\,\Sigma(p)\max_{t_0\in\linkGVV(p)}S_v(p,t_0)
\end{equation}
if $\linkGVV(p)\ne\emptyset$, and $\tau(p)=\iota(p)=0$ otherwise; here $e_p$ is the sandwich exponent chosen by
$X$, and $S_v(p,t_0)$ is the ceiling of the sum over the lifts $j$ of the positive part of an explicit
\emph{orbit gain} (a difference of sums over the tops of the orbit involving the weights $a_T,b_T$ of $L$). The
\emph{gain} $\mathrm{gain}_P(p,j)$ is the sum over the served orbits of the orbit gains and over the other
representatives of their positive ``odd'' parts, and is $0$ if $p$ is in the inert class with $h_p=0$, or if
$p>X''(n)$.
\formal{Formalism.MainNotions.G.tauP}
\formal{Formalism.MainNotions.G.gainP}
\end{definition}

\paragraph{The link.}
The size data $L$ contain numbers $\tau(w),\iota(w)$, local lengths $\mathrm{len}_{MN}(u,w)$, gains
$\mathrm{gain}(j,w)$, sets of tops $TT_w$, von Mangoldt class sums $\Lambda_{j,t}(p)$ and the lattices $N_u$ of
\cref{sec:families}; ``fill'' identities of the family (see \eqref{eq:pf:linkG:fill}) relate them to the quantities just
defined. The error terms $\mathrm{lo}_G(p)=\mathrm{lo}_G(L,C,b_o,p)$ and
$\mathrm{fl}_G(p)=Cf\,N(m)^2\log_pa+Cf(\log_pa)^2$ are explicit expressions in the scale.

\begin{definition}[The local lattice link]\label{pf:linkG:linkGat}
Let $(D,L)$ be a member with parameter instance $X$, link constants $K$, and let $C,b_o\in\R$.
$\sfname{LinkG}_{\mathrm{at}}(X,K,C,b_o)$ is the conjunction of:
\begin{enumerate}[label=(\roman*)]
\item for every word $u$ and place $w$ with $p(w)\ne T_u$: the lattice $N_u$ is contained, locally at $w$, in
$\mathfrak p_w^{\lfloor e_w\tau(w)\rfloor}M_u^{\mathrm{ram}}$, where $\mathfrak p_w$ is the prime ideal of $w$, $e_w$ its
ramification index and $M_u^{\mathrm{ram}}$ the ramified lattice of $u$ (every $x\in N_u$ satisfies
$sx\in\mathfrak p_w^{\lfloor e_w\tau(w)\rfloor}M_u^{\mathrm{ram}}$ for some $s\notin\mathfrak p_w$);
$\mathrm{len}_{MN}(u,w)\le d_u(\tau(w)+\iota(w))$; and if $\linkGVV(p(w))=\emptyset$ then
$\tau(w)=\iota(w)=\mathrm{len}_{MN}(u,w)=0$;
\item for every prime $p$ and $t_0\in\linkGVV(p)$: $S_v(p,t_0)\le2m\,N_v(p,t_0)\,a/p$;
\item for every place $w$ with $p=p(w)$ in the inert or signal class and $p\le X''(n)$: for all $1\le j\le m$,
$0\le\mathrm{gain}(j,w)\le4|TT|a/p$, and
\begin{equation}\label{eq:pf:linkG:vonMangoldt}
\sum_{j=1}^m\sum_{t\in TT_w}\bigl(\Lambda_{j,t}(p)-\mathrm{lo}_G(p)\bigr)-m\,\mathrm{fl}_G(p)\ \le\
\sum_{j=1}^m\mathrm{gain}(j,w)-\tau(w);
\end{equation}
\item the served orbits contain all good representatives of every candidate $p$ with $\Sigma(p)\le x^{-\eps}$, and
the mean of $\kappa$ over the weight boxes is at most $N_v(p,t_0)x^{-12}$ at every good orbit of every candidate
(and $\kappa_v\le N_vx^{-12}$ pointwise when $p<p_0(k)$).
\end{enumerate}
$\sfname{LinkG}(\mathcal F)$ means: there are $C,x_0$ and a function $\delta$ with $\delta(x)\to0$ as $x\to\infty$
such that $\sfname{LinkG}_{\mathrm{at}}(G_x(x,i),G_K,C,b^{\delta(x)})$ for all $x\ge x_0$ and $i\in I_x$.
\formal{Formalism.MainNotions.LinkG_at}
\formal{Formalism.MainNotions.LinkG}
\end{definition}

The two hypotheses of \cref{thm:linkG} are as follows. The \emph{group compatibility} $\mathsf{GC}(\mathcal F)$ says that
eventually, for every member, the comparison data $\mathrm{ex}(x,i).\mathrm{cmp}$ admit a \emph{group compatibility
record}: a hermitian level datum, forms at all words, and a comparison record tying them to the local group data
(it contains, for every $(u,w)$ and place $v$, ring isomorphisms identifying the rings of $\mathcal G(u,w)$ with
explicit completions, and a description of $\mathcal G(u,w)$ as the datum attached to a definite unitary group).
$\sfname{Adm2F}(\mathcal F)$ says that eventually $\tfrac12n\delta_1\le|TT|\le2n\delta_1$ and
\begin{equation}\label{eq:pf:linkG:adm2}
B_{TT}:=\sum_{p}\frac{\log p}{p}\,\#\{t\in TT:\ \text{the orbit of }t\text{ at }p\text{ is not good}\}\le|TT|\,x^{-2+\eps},
\end{equation}
the sum over the primes $p<X'$ in the inert class with $\ord(h_p)\ge x^{\eps/4}$.

% ---------------------------------------------------------------------------
\subsection{The local lattice theorem}\label{sec:pf:linkG:local}

Throughout this subsection $p$ is a prime and all rings are commutative. For a discrete valuation ring $R$ with
fraction field $K$ and a $K$-vector space $V$, an $R$-submodule $A\subseteq V$ and $c\in K$, we write $cA$ for the
image of $A$ under multiplication by $c$, and $A^\vee=\{G\in\Hom_K(V,K):\ G(A)\subseteq R\}$ for the \emph{dual
lattice}.

\begin{definition}[Hida settings]\label{pf:linkG:hida}
A \emph{Hida setting} $(p,n',W,Q_q,\OO,K_f,\iota,\sigma,j,\varpi)$ consists of discrete valuation rings $W$ and
$\OO$ with fraction fields $Q_q$ and $K_f$, a field embedding $\iota:Q_q\to K_f$, a field automorphism $\sigma$ of
$Q_q$, a ring homomorphism $j:W\to\OO$ and an element $\varpi\in\OO$, such that: $p$ is a prime with $p\ge5$;
$n'\ge2$ is even; $p$ is a uniformiser of $W$, $|W/pW|=p^{n'}$ and $W$ is $p$-adically complete; for every
$x\in W$ there is $y\in W$ with $\sigma(x)=x^p+py$; $\sigma^{n'}=1$; $\iota$ restricted to $W$ is $j$ followed by
the inclusion $\OO\subseteq K_f$; $j$ is finite; $\OO$ is complete; and $\varpi^2=p$. It is an \emph{unramified
quadratic setting} if moreover $\varpi$ is a uniformiser of $\OO$, $|\OO/\varpi\OO|=p^{2n'}$, and every rational
integer is a square in $\OO$.
\formal{Formalism.MainNotions.G.HidaSetting}
\formal{Formalism.MainNotions.G.UnramifiedQuadraticSetting}
\end{definition}

Thus $W$ is the unramified extension of $\Z_p$ of degree $n'$ with its Frobenius $\sigma$. In an unramified quadratic
setting the residue field of $\OO$ has $p^{2n'}$ elements and $\varpi^2=p$ with $\varpi$ a uniformiser, so $\OO$ is a
totally ramified quadratic extension of the unramified extension of $\Z_p$ of degree $2n'$. Part~I writes $w$ for
$\varpi$.

\begin{definition}[Lattices and Hecke edges]\label{pf:linkG:edges}
Fix $k\ge1$. A \emph{lattice} $L\subseteq Q_q^k$ is the $W$-span of the columns of a matrix $g\in\GL_k(Q_q)$;
the \emph{standard lattice} is $W^k$, the span of the columns of the identity. An \emph{edge} is a pair
$e=(e_s,e_t)$ of $W$-submodules of $Q_q^k$, the \emph{source} and the \emph{target}. For $n',r\ge0$ the edge $e$
is a \emph{Hecke edge of type $r$} if $e_s$ and $e_t$ are lattices, $p\,e_t\subseteq e_s\subseteq e_t$, and
$[e_t:e_s]=p^{n'r}$.
\formal{Formalism.MainNotions.G.IsLattice}
\formal{Formalism.MainNotions.G.IsEdge}
\end{definition}

Since $|W/pW|=p^{n'}$, a Hecke edge of type $r$ with target $L$ is the same as a $W/p$-subspace of $L/pL$ of
codimension $r$, and its source is the preimage of that subspace.

\begin{definition}[Weight lattices]\label{pf:linkG:weightlat}
Let $R$ be a ring, $K$ a field containing $R$, and $k,n\ge1$.
\begin{enumerate}[label=(\alph*)]
\item In the polynomial ring $R[x]=R[x_{ab}:\ a<k,\ b<n]$ the variable $x_{ab}$ has \emph{column weight} the
$b$-th unit vector of $\N^n$. For $\nu\in\N^n$, the \emph{dual Weyl module} $\nabla_R(\nu)\subseteq R[x]$ is the
module of polynomials $F$ of column weight $\nu$ which are invariant under the column operations
$x_{\cdot j}\mapsto x_{\cdot j}+t\,x_{\cdot l}$ for all $l<j$ ($t$ an additional variable). For
$g\in M_k(K)$ let $\rho(g)$ be the substitution $x_{ab}\mapsto\sum_{c}g_{ac}x_{cb}$, and put
$\nabla_R(\nu,g)=\{F\in K[x]:\ \rho(g)F\in\nabla_R(\nu)\}$.
\item Let $k=n=m+1$ and let $\lambda\in\Z^{m+1}$ be non-increasing. For $g\in\GL_{m+1}(K)$ put
$\nabla^{\mathrm{tw}}_R(\lambda,g)=\det(g)^{-\lambda_m}\nabla_R(\lambda-\lambda_m,g)$. For an $R$-submodule
$\Lambda\subseteq K^{m+1}$, $H^0(\lambda,\Lambda)$ is the intersection of the $\nabla^{\mathrm{tw}}_R(\lambda,g)$ over
all $g\in\GL_{m+1}(K)$ whose columns span $\Lambda$ over $R$; for $\lambda$ not non-increasing,
$H^0(\lambda,\Lambda)=0$.
\item Now fix a Hida setting, $m\ge0$ and a \emph{weight} $\lambda_0:\Z/n'\to\Z^{m+1}$; $\lambda_0$ is
\emph{dominant} if every $\lambda_0(s)$ is non-increasing. Put $\lambda_s=\lambda_0(s)$ for $s$ even and
$\lambda_s=(-\lambda_0(s)_{m-i})_{0\le i\le m}$ for $s$ odd, and
\[
A_r(\lambda_0)=\sum_{s\in\Z/n'}\ \sum_{i<r}\lambda_s(i).
\]
The space $\linkGV=\linkGV(\lambda_0)$ is the tensor product over $s\in\Z/n'$ of the
$\nabla_{K_f}(\lambda_s-\lambda_s(m))$ (a factor is $0$ if $\lambda_s$ is not non-increasing), realised in
$K_f[x_{s,a,b}]$ as the span of products of polynomials in
disjoint sets of variables; $\dim\linkGV$ is its $K_f$-dimension. For $\OO$-submodules $\Lambda_s\subseteq
K_f^{m+1}$ the \emph{tensor lattice} $\bigotimes_sH^0(\lambda_s,\Lambda_s)$ is the $\OO$-span of the products of
elements of the $H^0(\lambda_s,\Lambda_s)$.
\item For a $W$-module $L\subseteq Q_q^{m+1}$ and $s\in\Z/n'$ let $L^{(s)}\subseteq K_f^{m+1}$ be the $\OO$-span of
the vectors $\iota(\sigma^s(v))$, $v\in L$. The \emph{weight lattice} of $L$ is
$\linkGH_L=\linkGV\cap\bigotimes_sH^0(\lambda_s,L^{(s)})$. For an edge $e$ put
$L_e^{(s)}=e_s^{(s)}+\varpi^{-1}(p\,e_t)^{(s)}$, $\linkGH_{e}=\linkGV\cap\bigotimes_sH^0(\lambda_s,L_e^{(s)})$ and
\[
\linkGR_e=\varpi^{-A_r(\lambda_0)}\linkGH_e\subseteq\linkGV,\qquad
\linkGR'_e=\varpi^{-A_r(\lambda_0)}\linkGH_e^\vee\subseteq\linkGV^\vee .
\]
\item The pairing $b:\linkGV\times\linkGV^\vee\to K_f$ is $b(F,G)=G(p^{A_r(\lambda_0)}F)$. For $t\in\N$ and
lattices $L'',L$ put
\[
\linkGNs(t,L'')=\varpi^{-t}\linkGH_{L''}\cap\bigcap_{e:\ e_s=L''}\linkGR_e,\qquad
\linkGNt(t,L)=\varpi^{-t}\linkGH_L^\vee\cap\bigcap_{e:\ e_t=L}\linkGR'_e,
\]
the intersections over the Hecke edges of type $r$.
\end{enumerate}
\formal{Formalism.GradedWeylDivisibility.dualWeyl}
\formal{Formalism.MainNotions.G.H0}
\formal{Formalism.MainNotions.G.HL}
\formal{Formalism.MainNotions.G.Re}
\formal{Formalism.MainNotions.G.Nsrc}
\formal{Formalism.MainNotions.G.Ntgt}
\end{definition}

The model (a) of the dual Weyl module is the classical description of the induced module as functions on matrices
invariant under right translation by upper unitriangular matrices, and $H^0(\lambda,\Lambda)$ is an integral form
attached to the lattice $\Lambda$, independent of the choice of basis by construction \cite{Jantzen2003}. Since
$p\,e_t\subseteq e_s\subseteq e_t$, the module $L_e^{(s)}=e_s^{(s)}+\varpi\,e_t^{(s)}$ lies halfway between the
two ends of the edge.

\begin{definition}[Good rank and the property $\linkGLP$]\label{pf:linkG:LP}
Let $\OO$ be a ring with fraction field $K_f$, $A$ an $\OO$-module, $\pi\in\OO$, $I$ a finite set and
$\varphi_i:A\to K_f$ ($i\in I$) maps. For $n,D\in\N$, $\linkGGR(\pi,A,\varphi,n,D)$ means: there are an
$\OO$-linear map $u:\OO^n\to A$ and a matrix $(M_{ji})\in\OO^{n\times I}$ with
$\sum_iM_{ji}\varphi_i(u(x))=\pi^Dx_j$ for all $x\in\OO^n$ and $j$. In a Hida setting, with $m,r,\lambda_0$ fixed,
let $t,D,c_*\in\N$, $\kappa\in\R$, and let $L\subseteq Q_q^{m+1}$ be a lattice.
\begin{enumerate}[label=(\alph*)]
\item $\linkGLPat(t,D,\kappa,c_*,L)$ means: there are Hecke edges $e_1,\dots,e_M$ of type $r$ with target $L$ and
vectors $F_1,\dots,F_M\in\linkGV$ such that $F_i\in\linkGNs(t,e_{i,s})$, $F_i\in\varpi^{c_*}\linkGR_{e'}$ for every
Hecke edge $e'\ne e_i$ of type $r$ with source $e_{i,s}$, and $\linkGGR\bigl(p,\linkGNt(t,L),(G\mapsto
b(F_i,G))_i,n,D\bigr)$ for some $n\in\N$ with $n\ge(1-\kappa)\dim\linkGV$.
\item $\linkGLP(t,D,\kappa,c_*)$ means: $4(D+1)\le c_*$ and $\linkGLPat(t,D,\kappa,c_*,L)$ for every lattice $L$;
$\linkGLP_1(t,D,\kappa,c_*)$ means: $4(D+1)\le c_*$ and $\linkGLPat(t,D,\kappa,c_*,L)$ for some lattice $L$.
\end{enumerate}
\formal{Formalism.MainNotions.G.GoodRankK}
\formal{Formalism.MainNotions.G.LPprimeAt}
\formal{Formalism.MainNotions.G.LPprime}
\end{definition}

$\linkGGR(\pi,A,\varphi,n,D)$ says that the functionals $\varphi_i$, restricted to $A$, have ``rank at least $n$ up
to the denominator $\pi^D$''. Two monotonicity properties are immediate: $\linkGLP(t,D,\kappa,c_*)$ implies
$\linkGLP(t,D,\kappa',c_*)$ for $\kappa\le\kappa'$, and $\linkGLP(t,D,\min(1,\kappa),c_*)$ is equivalent to
$\linkGLP(t,D,\kappa,c_*)$, because $(1-\kappa)\dim\linkGV\le0$ when $\kappa\ge1$.

\begin{definition}[Orbit parameters]\label{pf:linkG:orbitpar}
Let $r\ge1$ and $\delta,h,x\in\R$. Put $u=\lceil1/\delta\rceil$, $\ell=(6r+1)u$,
$M_0=\lceil10^{6r+1}(h+3)\log x\rceil$, $G'=(M_0+5)\ell$ and
$N_1=\lceil\max(6e\,G'\delta,(h+3)\log_2x)\rceil+1$ (we write $G'=G'(r,\delta,h,x)$ and $N_1=N_1(r,\delta,h,x)$).
Let $n\ge1$ and $T\subseteq\{0,\dots,n-1\}$. A \emph{window} of length $L$ at $i$ is
$\{(i+j)\bmod n:\ 0\le j<L\}$. For $s\ge0$ and $1\le c\le6r+1$ the \emph{cell} $(s,c)$ is the interval
$[s\ell+(c-1)u,\,s\ell+cu)$; the tile $s$ is \emph{marked} by $T$ if $T$ meets no cell $(s,c)$ with $c$ odd, and
meets each cell $(s,c)$ with $c$ even in exactly one point, which is even. For $L_2,K\in\N$, $T$ satisfies the
\emph{good-orbit tile condition} if (i) every window of length $L_2$ contains at most $K$ points of $T$; (ii) for
every tile index $s_0<n/\ell$ one of the tiles $s_0+j\bmod\lfloor n/\ell\rfloor$, $0\le j\le M_0$, is marked; and
(iii) every window of length $G'$ contains at most $N_1$ points of $T$. For a set $TT\subseteq\Z/n'$ and
$L_1,L_2\in\N$, the \emph{orbit hypothesis} $\mathrm{OH}(n',TT,r,L_1,L_2,\delta,h,x)$ means: $TT$ (viewed in
$\{0,\dots,n'-1\}$) satisfies the good-orbit tile condition with $(n,L_2,K)=(n',L_2,r-1)$, $2L_1\le L_2<u$, and
$3G'\le n'$.
\formal{Formalism.MainNotions.G.IsGoodOrbit}
\formal{Formalism.MainNotions.tableLineDegreeFive.OrbitHyp}
\end{definition}

\begin{definition}[Weight hypothesis]\label{pf:linkG:weighthyp}
Let $TT\subseteq\Z/n'$, $a_T,b_T:\Z/n'\to\Z$, $c_{pp},c_0\in\N$ and $a\in\R$. A weight $\lambda_0$ is a
\emph{base weight} if for $s\in TT$: $\lambda_0(s)_0=a_T(s)$, $\lambda_0(s)_m=-b_T(s)$ and
$|\lambda_0(s)_i|\le c_{pp}$ for $0<i<m$; and for $s\notin TT$: $|\lambda_0(s)_i|\le c_{pp}$ for all $i$. The
\emph{weight hypothesis} $\mathrm{WH}$ means: $\lambda_0$ is a dominant base weight, and for $s\in TT$:
$a+c_{pp}\le a_T(s)\le2a-c_{pp}$, $0\le b_T(s)$ and $b_T(s)+2c_{pp}\le c_0$.
\formal{Formalism.MainNotions.tableLineDegreeFive.WeightHyp}
\end{definition}

\begin{definition}[The local statement]\label{pf:linkG:LocLat}
For $C\in\R$ put $\Lambda_C(N_1,a)=CN_1^2\log a$ and, for $k=m+1$,
\[
t_{\mathrm{B}}=C\,r\,|TT|\Bigl(\frac{\Lambda_C(N_1,a)\,a}{pk}+a\,p^{-L_1}+c_0+k\,p^{G'}\Bigr)+2rn'c_{pp}.
\]
$\linkGLL(C)$ is the following statement. Let $(p,n',W,Q_q,\OO,K_f,\iota,\sigma,j,\varpi)$ be an unramified quadratic
setting with $p\ge13$; let $m\in\N$, $r\ge2$, $TT\subseteq\Z/n'$, $L_1,L_2\in\N$, $\delta,h,x,a,Y\in\R$, a weight
$\lambda_0$, $a_T,b_T:\Z/n'\to\Z$ and $c_{pp},c_0,D_p,c_*\in\N$; put $G'=G'(r,\delta,h,x)$, $N_1=N_1(r,\delta,h,x)$
and $k=m+1$. Assume
\begin{enumerate}[label=(H\arabic*)]
\item\label{pf:linkG:H1} $\mathrm{OH}(n',TT,r,L_1,L_2,\delta,h,x)$ and $N_1^2\le r|TT|$;
\item\label{pf:linkG:H2} $\mathrm{WH}(TT,a_T,b_T,c_{pp},c_0,a,\lambda_0)$;
\item\label{pf:linkG:H3} $r\le Y$, $\Lambda_C(N_1,a)\le k$ and $N_1+2r\le k-r$;
\item\label{pf:linkG:H4} $256k^2|TT|Y^2\le p^{L_1-1}$, $C|TT|k^2\log_pa\le D_p$ and $8(D_p+1)\le c_*$;
\item\label{pf:linkG:H5} $C\,k^3r^2|TT|\,(p^{G'})^5\bigl(c_*+r|TT|c_0+2rn'c_{pp}\bigr)\,Y\le a$.
\end{enumerate}
Then $\linkGLP\bigl(t,D_p,\,C|TT|r^2k/Y,\,c_*\bigr)$ holds for every integer $t\ge t_{\mathrm B}$.
\formal{Formalism.MainNotions.tableLineDegreeFive.aux₃}
\end{definition}

In the development $\linkGLL(C)$ is a proposition $\mathrm{aux}_3(D,C)$ indexed by a datum $D$; the datum does not
occur in it (\cref{sec:pf:linkG:remarks}). The proof of the local lattice theorem is conditional on two packages
of results, which are themselves proved in the development (\cref{pf:linkG:locmodel}).

\begin{definition}[The tool packages]\label{pf:linkG:hyps}
\begin{enumerate}[label=(\alph*)]
\item $(\mathrm E)$ is the \emph{extremal projector statement}: for every field $K$ of characteristic $0$, every
finite-dimensional representation $(E_{ab})_{a,b<n}$ of $\mathfrak{gl}_n$ on a $K$-vector space $V$, every
$\mu\in\Z^n$ with $\mu_a-\mu_b+(b-a)\ge1$ for $a<b$, every normal ordering of the positive roots and every
$T\ge\dim V$, an explicit operator $P_\mu=\mathrm{pp}(E,\mu,\mathrm{ord},T)$ preserves the weight space $V_\mu$,
annihilates $V_\mu\cap\mathfrak n^-V$, and fixes every vector of $V_\mu$ annihilated by the raising operators
$E_{ab}$, $a<b$.
\item $(\mathrm{L}6)$ is the conjunction of the following seventeen statements:
\begin{itemize}
\item the \emph{criterion} $\mathsf{Crit}(D)$ for every datum $D$ (\cref{pf:linkG:criterion});
\item the ``second part'' (\cref{pf:linkG:planecount}) for every $D$;
\item three counting forms of the Weyl dimension formula, the branching rule and the Pieri rule for
Gelfand--Tsetlin patterns;
\item seven lattice statements for the weight lattices (\cref{tab:pf:linkG:clauses}):
clauses 3, 4, 5 and 29--32;
\item the Witt identification (clause 28);
\item for every $D$, the four clauses numbered 12, 23, 36 and 37 in \cref{tab:pf:linkG:clauses}.
\end{itemize}
\end{enumerate}
\formal{Formalism.MainNotions.G.ExtremalProjectorTools}
\formallabel{Formalism.MainNotions.G.Loc6}
\end{definition}

\begin{theorem}[\cref{thm:families:local}, local lattice theorem]\label{pf:linkG:localthm}
Assume $(\mathrm E)$ and $(\mathrm L6)$. Then there is $C_0\ge0$ such that $\linkGLL(C_0)$ holds, in the following
form: for every datum $D$, the statement $\mathrm{aux}_3(D,C_0)$ (which is $\linkGLL(C_0)$) holds.
\formallabel{Formalism.MainNotions.G.second_theorem_absolute_constant}
\end{theorem}

The theorem is proved in \cref{sec:pf:linkG:verif}. Since $(\mathrm E)$ and $(\mathrm L6)$ hold
(\cref{pf:linkG:locmodel}), the constant $C_0$ is absolute, as stated in \cref{thm:families:local}.

% ---------------------------------------------------------------------------
\subsection{Reduction to edge families}\label{sec:pf:linkG:edges}

We fix a Hida setting $(p,n',W,Q_q,\OO,K_f,\iota,\sigma,j,\varpi)$, integers $m,r$ and a weight $\lambda_0$, and write
$k=m+1$. ``Edge'' means Hecke edge of type $r$ in $Q_q^k$. An element $x\in K_f$ \emph{divides $p^D$} if
$x\in\OO$ and $xy=p^D$ for some $y\in\OO$. A vector $c\in\OO^n$ is \emph{primitive} if some $c_i$ is a unit.

\begin{definition}[Edge families, source tests, edge step]\label{pf:linkG:edgefam}
\begin{enumerate}[label=(\alph*)]
\item Let $t_S,t_T,D,c_*,n\in\N$ and let $e$ be an edge. $\linkGEF(t_S,t_T,D,c_*,e,n)$ means: there are
$F_1,\dots,F_n\in\linkGV$ such that (i) $F_i\in\varpi^{-t_S}\linkGH_{e_s}\cap\linkGR_e$; (ii) $F_i\in\varpi^{c_*}\linkGR_{e'}$
for every edge $e'\ne e$ with $e'_s=e_s$; (iii) for every primitive $c\in\OO^n$ there is
$G\in\varpi^{-t_T}\linkGH_{e_t}^\vee\cap\linkGR'_e$ with $G\in\varpi^{c_*}\linkGR'_{e'}$ for every edge $e'\ne e$ with
$e'_t=e_t$, such that $b(\sum_ic_iF_i,G)$ divides $p^D$.
\item Let $t,D,c_*\in\N$, $\kappa\in\R$ and let $L$ be a lattice. $\linkGST(t,D,\kappa,c_*,L)$ means: there are a
finite set $E$ of edges with target $L$, integers $n_e\ge0$ and vectors $F_{e,i}\in\linkGV$ ($e\in E$, $i<n_e$) such
that $F_{e,i}\in\linkGNs(t,e_s)$; $F_{e,i}\in\varpi^{c_*}\linkGR_{e'}$ for every edge $e'\ne e$ with $e'_s=e_s$; for
every $e\in E$ and every primitive $c\in\OO^{n_e}$ there is $G\in\linkGNt(t,L)$ such that
$b(\sum_ic_iF_{e,i},G)$ divides $p^D$ and $G\in\varpi^{c_*}\linkGR'_{e'}$ for all $e'\in E\smallsetminus\{e\}$; and
$\sum_{e\in E}n_e\ge(1-\kappa)\dim\linkGV$.
\item For constants $C,C_{t},C_{D},C_r$, $\linkGES(C,C_t,C_D,C_r)$ is the statement: for every unramified quadratic
setting with $p\ge13$ and all parameters as in \cref{pf:linkG:LocLat} with $r\ge2$ satisfying
\ref{pf:linkG:H1}--\ref{pf:linkG:H5} (with the constant $C$), and every edge $e$ with target $W^k$, there are
$n,t_S,t_T,D_1\in\N$ with
\begin{gather*}
t_S\le t_0+\omega,\qquad t_T\le t_0+|TT|c_0+\omega,\qquad D_1\le C_D|TT|k^2\log_pa,\\
\Bigl(1-\frac{C_r|TT|r^2k}{Y}\Bigr)\Bigl(1-\frac{kc_0}{a}\Bigr)^{|TT|}\frac{\dim\linkGV}{p^{n'r(k-r)}}\le n,
\end{gather*}
and $\linkGEF(t_S,t_T,D_1,c_*,e,n)$, where
$t_0=C_t|TT|\bigl(\Lambda_C(N_1,a)a/(pk)+a\,p^{-L_1}+k\,p^{G'}\bigr)$ and $\omega=r|TT|c_0+2rn'c_{pp}$.
\end{enumerate}
\formallabel{Formalism.MainNotions.Pf14.EdgeFam}
\formallabel{Formalism.MainNotions.Pf14.SrcTests}
\formallabel{Formalism.MainNotions.Pf14.EdgeStep}
\end{definition}

Two of the clauses of \cref{sec:pf:linkG:criterion} enter the reduction.

\begin{definition}[Perfectness and homogeneity]\label{pf:linkG:perfect}
The \emph{perfectness statement} $(\mathrm P)$ is: for every Hida setting, $m,r$, dominant $\lambda_0$ and edge $e$,
$G\in\linkGR'_e$ if and only if $b(F,G)\in\OO$ for all $F\in\linkGR_e$, and $F\in\linkGR_e$ if and only if $b(F,G)\in\OO$
for all $G\in\linkGR'_e$. The \emph{homogeneity statement} $(\mathrm H)$ is: for every Hida setting, $m,r$, dominant
$\lambda_0$ and $t,D,\kappa,c_*$, $\linkGLP_1(t,D,\kappa,c_*)$ implies $\linkGLP(t,D,\kappa,c_*)$.
\formal{Formalism.MainNotions.G.BePerfect}
\formal{Formalism.MainNotions.G.LPOneEvery}
\end{definition}

Both are proved in the development: $(\mathrm P)$ in \lref{Formalism.L2.G.Model.bePerfect} (the first
equivalence is a direct computation from the definitions, $b(\varpi^{-A}z,\varpi^{-A}\psi)=\psi(z)$ because
$p^A\varpi^{-2A}=1$; the second uses an elementary divisor basis adapted to the edge), and $(\mathrm H)$ in
\lref{Formalism.L2.G.Model.lPOneEvery} (any two lattices differ by an element of $\GL_k(Q_q)$, and
$\linkGLPat$ is transported along this action).

\begin{lemma}[Many edges]\label{pf:linkG:manyedges}
For every Hida setting and $r\le k$ there is a set $E$ of Hecke edges of type $r$ in $Q_q^k$ with target $W^k$ and
$|E|\ge p^{n'r(k-r)}$.
\formallabel{Formalism.MainNotions.Pf14.manyEdges}
\end{lemma}

\begin{proof}
Write $k=r+K$, $\bar W=W/pW$, and choose a lift $\tilde x\in W$ of every $x\in\bar W$. For a matrix
$A\in\bar W^{r\times K}$ let $\psi_A:W^k\to\bar W^r$ be the $W$-linear map
$\psi_A(v)_i=\bigl(v_i+\sum_{c<K}\tilde A_{ic}v_{r+c}\bigr)\bmod p$, and let $e(A)$ be the edge with target $W^k$ and
source $\ker\psi_A$ (viewed in $Q_q^k$). Since $\psi_A$ vanishes on $pW^k$ we have $pW^k\subseteq\ker\psi_A\subseteq
W^k$. The map $\psi_A$ is surjective (take $v=(\tilde x,0)$), so $[W^k:\ker\psi_A]=|\bar W|^r=p^{n'r}$. As a
submodule of the free module $W^k$ over the principal ideal domain $W$, $\ker\psi_A$ is free, of rank $k$ because it
contains $pW^k$; hence it is the span of the columns of an invertible matrix, i.e.\ a lattice. So $e(A)$ is a Hecke
edge of type $r$. The map $A\mapsto e(A)$ is injective: if $\ker\psi_A=\ker\psi_B$, then for $c<K$ the vector $v$
with $v_{r+c}=1$, $v_{r+c'}=0$ for $c'\ne c$ and $v_i=-\tilde A_{ic}$ for $i<r$ lies in $\ker\psi_A$, and
$\psi_B(v)_i=B_{ic}-A_{ic}=0$ for all $i$. Hence the set of the $e(A)$ has $|\bar W|^{rK}=p^{n'rK}$ elements.
\end{proof}

For comparison, the number of all Hecke edges of type $r$ with a given target is the Gaussian binomial coefficient
$\binom{k}{r}_{p^{n'}}$, which is at least $p^{n'r(k-r)}$ by \cref{pf:linkG:gauss}.

\begin{lemma}[Smith normal form]\label{pf:linkG:smith}
Let $\OO$ be a local principal ideal domain, $I$ a finite set, $S\subseteq\OO^I$ a submodule and $a\in\OO$.
Suppose that for every $\varphi\in\OO^I$ with some unit coordinate there is $s\in S$ such that $\sum_i\varphi_is_i$
divides $a$. Then $ax\in S$ for every $x\in\OO^I$.
\formallabel{Formalism.MainNotions.Pf14.smith_mem}
\end{lemma}

\begin{proof}
By the elementary divisor theorem there are a basis $(b_j)_{j\in I}$ of $\OO^I$, an integer $n$, an injection
$f:\{1,\dots,n\}\to I$ and elements $a_1,\dots,a_n\in\OO$ such that $(a_\nu b_{f(\nu)})_\nu$ is a basis of $S$. Fix
$j\in I$ and let $b_j^*$ be the $j$-th coordinate functional of the basis $(b_j)$; write $b_j^*(y)=\sum_i\varphi_iy_i$
with $\varphi_i=b_j^*(\text{$i$-th unit vector})$. Some $\varphi_i$ is a unit: otherwise
$1=b_j^*(b_j)=\sum_i\varphi_i(b_j)_i$ would lie in the maximal ideal. By hypothesis there is $s\in S$ with
$b_j^*(s)\mid a$. Write $s=\sum_\nu c_\nu a_\nu b_{f(\nu)}$; then $b_j^*(s)=c_\nu a_\nu$ if $j=f(\nu)$, and
$b_j^*(s)=0$ if $j\notin f(\{1,\dots,n\})$. In the first case $a_\nu\mid a$, say $a=ta_\nu$, and
$ab_j=t\,(a_\nu b_{f(\nu)})\in S$. In the second case $0\mid a$ forces $a=0$ and $ab_j=0\in S$. Hence $ab_j\in S$ for
all $j$, and $ax=\sum_jb_j^*(x)\,ab_j\in S$.
\end{proof}

\begin{lemma}[Good rank from primitive tests]\label{pf:linkG:goodrank}
Let $\OO$ be a discrete valuation ring with fraction field $K_f$, $A$ a submodule of an $\OO$-module, $\pi\in\OO$,
$D\in\N$, and $\varphi_1,\dots,\varphi_M:A\to K_f$ additive maps with $\varphi_i(cG)=c\varphi_i(G)$ for $c\in\OO$. Assume
$\varphi_i(A)\subseteq\OO$ for all $i$, and that for every primitive $c\in\OO^M$ there is $G\in A$ such that
$\sum_ic_i\varphi_i(G)$ divides $\pi^D$. Then $\linkGGR(\pi,A,\varphi,M,D)$ holds, with $M_{ji}=\delta_{ji}$.
\formallabel{Formalism.MainNotions.Pf14.goodRank_of_prim}
\end{lemma}

\begin{proof}
The set $S=\{(\varphi_i(G))_i:\ G\in A\}$ is an $\OO$-submodule of $\OO^M$. For primitive $c$ choose $G$ as in the
hypothesis; then $s=(\varphi_i(G))_i\in S$ and $\sum_ic_is_i\mid\pi^D$. By \cref{pf:linkG:smith},
$\pi^D\epsilon_l\in S$ for every unit vector $\epsilon_l$, i.e.\ there are $G_l\in A$ with
$\varphi_i(G_l)=\pi^D\delta_{il}$. The map $u(x)=\sum_lx_lG_l$ is $\OO$-linear and
$\varphi_j(u(x))=\sum_lx_l\varphi_j(G_l)=\pi^Dx_j$.
\end{proof}

\begin{lemma}[Source tests give $\linkGLPat$]\label{pf:linkG:s7}
Assume $(\mathrm P)$. Let $\lambda_0$ be dominant, $t,D,c_*\in\N$ with $2(D+1)\le c_*$, $\kappa\in\R$ and $L$
a lattice. Then $\linkGST(t,D,\kappa,c_*,L)$ implies $\linkGLPat(t,D,\kappa,c_*,L)$.
\formallabel{Formalism.MainNotions.Pf14.s7}
\formallabel{Formalism.MainNotions.Pf14.dvd_of_add}
\end{lemma}

\begin{proof}
Let $(E,(n_e),(F_{e,i}))$ be as in $\linkGST$, and enumerate the pairs $(e,i)$, $e\in E$, $i<n_e$, as $1,\dots,M$
with $M=\sum_en_e\ge(1-\kappa)\dim\linkGV$. Take the edges $e_{(e,i)}=e$ and the vectors $F_{(e,i)}=F_{e,i}$; the first
two conditions of $\linkGLPat$ are those of $\linkGST$, and it remains to prove
$\linkGGR\bigl(p,\linkGNt(t,L),(G\mapsto b(F_{e,i},G))_{(e,i)},M,D\bigr)$. We apply \cref{pf:linkG:goodrank} with
$A=\linkGNt(t,L)$; the maps $G\mapsto b(F,G)$ are additive and $\OO$-linear.

\emph{Integrality.} Since $F_{e,i}\in\linkGNs(t,e_s)\subseteq\linkGR_e$ and $\linkGNt(t,L)\subseteq\linkGR'_e$ ($e$ has
target $L$), $(\mathrm P)$ gives $b(F_{e,i},G)\in\OO$ for $G\in A$.

\emph{Primitive tests.} Let $c=(c_{e,i})$ be primitive, say $c_{e_0,i_0}$ is a unit. The vector
$(c_{e_0,i})_{i<n_{e_0}}$ is primitive, so $\linkGST$ provides $G\in A$ such that $\alpha=b(\sum_ic_{e_0,i}F_{e_0,i},G)$
divides $p^D$, and $G\in\varpi^{c_*}\linkGR'_{e}$ for all $e\in E\smallsetminus\{e_0\}$. For such $e$ write
$G=\varpi^{c_*}G_0$ with $G_0\in\linkGR'_e$; by $(\mathrm P)$, $y=b(F_{e,i},G_0)\in\OO$, and since $\varpi^2=p$
and $c_*\ge2(D+1)$,
\[
b(F_{e,i},G)=\varpi^{c_*}y=p^{D+1}\,\varpi^{c_*-2(D+1)}y\in p^{D+1}\OO .
\]
Hence $\sum_{(e,i)}c_{e,i}b(F_{e,i},G)=\alpha+p^{D+1}z$ with $z\in\OO$. Write $\alpha\beta=p^D$. Then
$\alpha+p^{D+1}z=\alpha(1+pz\beta)$, and $1+pz\beta$ is a unit since $p$ lies in the maximal ideal of $\OO$; so
$\alpha+p^{D+1}z$ divides $p^D$. \Cref{pf:linkG:goodrank} gives the claim.
\end{proof}

\begin{lemma}[Assembly of edge families]\label{pf:linkG:asm}
Let $L$ be a lattice, $E$ a finite set of edges with target $L$, $n_e\in\N$ ($e\in E$), and $t,D,c_*\in\N$,
$\kappa\in\R$. Suppose that for every $e\in E$ there are $t_S,t_T\le t$ and $D_1\le D$ with
$\linkGEF(t_S,t_T,D_1,c_*,e,n_e)$, and that $\sum_en_e\ge(1-\kappa)\dim\linkGV$. Then
$\linkGST(t,D,\kappa,c_*,L)$ holds.
\formallabel{Formalism.MainNotions.Pf14.asm}
\end{lemma}

\begin{proof}
Take for $F_{e,i}$ the vectors of the edge families. We use two inclusions valid for $\OO$-submodules $A$ and integers
$a\le b$, $c\ge0$: $\varpi^{-a}A\subseteq\varpi^{-b}A$ and $\varpi^cA\subseteq A$. Then
$F_{e,i}\in\varpi^{-t_S}\linkGH_{e_s}\subseteq\varpi^{-t}\linkGH_{e_s}$, $F_{e,i}\in\linkGR_e$, and
$F_{e,i}\in\varpi^{c_*}\linkGR_{e'}\subseteq\linkGR_{e'}$ for the other edges $e'$ with source $e_s$; so
$F_{e,i}\in\linkGNs(t,e_s)$, and the second condition of $\linkGST$ is condition (ii) of $\linkGEF$. For a primitive
$c$, condition (iii) gives $G\in\varpi^{-t_T}\linkGH_L^\vee\cap\linkGR'_e\subseteq\varpi^{-t}\linkGH_L^\vee\cap\linkGR'_e$
with $G\in\varpi^{c_*}\linkGR'_{e'}\subseteq\linkGR'_{e'}$ for all other edges $e'$ with target $L$, so $G\in\linkGNt(t,L)$;
these $e'$ include the elements of $E\smallsetminus\{e\}$; and $b(\sum c_iF_{e,i},G)$ divides $p^{D_1}$, hence $p^D$.
\end{proof}

\begin{lemma}[Numerics]\label{pf:linkG:num}
Let $C,C_t,C_D,C_r$ be real numbers with $C_t,C_D,C_r>0$, $C\ge2$, $C_t\le C$, $C_D\le C$ and $C_r+1\le C$. Let the
parameters of \cref{pf:linkG:LocLat} satisfy $p\ge13$, $r\ge2$ and \ref{pf:linkG:H1}--\ref{pf:linkG:H5}, and let
$t_0,\omega$ be as in \cref{pf:linkG:edgefam}. Then $a\ge1$ and
\begin{enumerate}[label=(\alph*)]
\item $t_0+|TT|c_0+\omega\le t_{\mathrm B}$;
\item $C_D|TT|k^2\log_pa\le D_p$;
\item $1-C|TT|r^2k/Y\le\bigl(1-C_r|TT|r^2k/Y\bigr)\bigl(1-kc_0/a\bigr)^{|TT|}$.
\end{enumerate}
\formallabel{Formalism.MainNotions.Pf14.num}
\formallabel{Formalism.MainNotions.Pf14.num_size}
\formallabel{Formalism.MainNotions.Pf14.level_bound_numeric}
\formallabel{Formalism.MainNotions.Pf14.num_count}
\formallabel{Formalism.MainNotions.Pf14.numeric_product_inequality}
\end{lemma}

\begin{proof}
We have $N_1\ge1$, so $|TT|\ge1$ by $N_1^2\le r|TT|$; also $k\ge1$, $Y\ge r\ge2$, $Q=p^{G'}\ge1$ and $c_*\ge8$. Since
$c_0\le r|TT|c_0+2rn'c_{pp}$ and $k|TT|\le Ck^3r^2|TT|Q^5$, hypothesis \ref{pf:linkG:H5} gives
\begin{equation}\label{eq:pf:linkG:numsize}
k|TT|\bigl(c_*+r|TT|c_0+2rn'c_{pp}\bigr)Y\le a,\qquad\text{hence}\qquad a\ge1\quad\text{and}\quad |TT|\,k\,c_0\,Y\le a,
\end{equation}
because all factors on the left are at least $1$. In particular $\log a\ge0$, $\log_pa\ge0$ and $\Lambda_C(N_1,a)\ge0$.
(a) Put $X=\Lambda_C(N_1,a)a/(pk)+a\,p^{-L_1}+k\,p^{G'}\ge0$. Since $C_t\le C\le Cr$ and $r+1\le2r\le Cr$,
\[
t_0+|TT|c_0+\omega=C_t|TT|X+(r+1)|TT|c_0+2rn'c_{pp}\le Cr|TT|(X+c_0)+2rn'c_{pp}=t_{\mathrm B}.
\]
(b) follows from $C_D\le C$, $\log_pa\ge0$ and \ref{pf:linkG:H4}. (c) Put $\beta=|TT|r^2k/Y$ and $y=kc_0/a\ge0$. By
\eqref{eq:pf:linkG:numsize}, $|TT|y\le1/Y\le\beta$ (as $|TT|r^2k\ge1$) and $y\le1/2$. By Bernoulli's inequality
$u=(1-y)^{|TT|}$ satisfies $1\ge u\ge1-|TT|y\ge1-\beta$. With $A=C_r\beta\ge0$ we get $(1-A)u\ge1-A-\beta$: if $A\le1$
then $(1-A)u\ge(1-A)(1-\beta)\ge1-A-\beta$, and if $A>1$ then $(1-A)u\ge1-A$. Finally $1-C\beta\le1-(C_r+1)\beta=1-A-\beta$.
\end{proof}

\begin{proposition}[From the edge step to the local statement]\label{pf:linkG:edgeassembly}
Assume $(\mathrm P)$ and $(\mathrm H)$. Let $C_t,C_D,C_r,C_1>0$ be such that
$\linkGES(C,C_t,C_D,C_r)$ holds for every $C\ge C_1$, and put $C_0=\max(C_1,2,C_t,C_D,C_r+1)$. Then $C_0>0$ and
$\linkGLL(C)$ holds for every $C\ge C_0$.
\formallabel{Formalism.MainNotions.Pf14.prvE_of_gl_operators_form_representation}
\formallabel{Formalism.MainNotions.Pf14.theorem_II_of_edge_bound_constants}
\end{proposition}

\begin{proof}
Let $C\ge C_0$, fix data satisfying the hypotheses of $\linkGLL(C)$, and let $t\ge t_{\mathrm B}$; put
$\kappa=C|TT|r^2k/Y$. The weight $\lambda_0$ is dominant by \ref{pf:linkG:H2}, and $4(D_p+1)\le8(D_p+1)\le c_*$. By
$(\mathrm H)$ it suffices to prove $\linkGLPat(t,D_p,\kappa,c_*,W^k)$, and since $2(D_p+1)\le c_*$,
\cref{pf:linkG:s7} reduces this to $\linkGST(t,D_p,\kappa,c_*,W^k)$.

We have $r\le k$ by \ref{pf:linkG:H3}. Let $E$ be the set of edges of \cref{pf:linkG:manyedges}, so that every $e\in E$
has target $W^k$ and $|E|\ge Q:=p^{n'r(k-r)}$. The hypotheses of $\linkGES(C,C_t,C_D,C_r)$ hold, so for every $e\in E$
there are $n_e,t_S(e),t_T(e),D_1(e)$ as in \cref{pf:linkG:edgefam}(c). By \cref{pf:linkG:num}, whose hypotheses
hold since $C\ge C_0$,
\[
t_S(e)\le t_0+\omega\le t_0+|TT|c_0+\omega\le t_{\mathrm B}\le t,\qquad t_T(e)\le t,\qquad D_1(e)\le C_D|TT|k^2\log_pa\le D_p .
\]
By \cref{pf:linkG:asm} it remains to show $\sum_{e\in E}n_e\ge(1-\kappa)\dim\linkGV$. Put
$\rho=(1-C_r|TT|r^2k/Y)(1-kc_0/a)^{|TT|}$. By \cref{pf:linkG:num}(c), $(1-\kappa)\dim\linkGV\le\rho\dim\linkGV$. If
$\rho\dim\linkGV<0$ there is nothing to prove. Otherwise every $n_e\ge\rho\dim\linkGV/Q\ge0$, and
\[
\sum_{e\in E}n_e\ \ge\ |E|\,\frac{\rho\dim\linkGV}{Q}\ \ge\ Q\,\frac{\rho\dim\linkGV}{Q}=\rho\dim\linkGV. \qedhere
\]
\end{proof}

In the development this proposition is stated with four abstract hypotheses, which are \cref{pf:linkG:manyedges,pf:linkG:s7,pf:linkG:asm,pf:linkG:num}; in the
combined form it takes $(\mathrm P)$ and $(\mathrm H)$ from the clauses (\cref{sec:pf:linkG:criterion})
and the edge step from a hypothesis of the form described there.

% ---------------------------------------------------------------------------
\subsection{The clauses and the criterion}\label{sec:pf:linkG:criterion}

The local lattice theorem is deduced from a list of $37$ statements, the \emph{clauses}. Formally they are
propositions $\mathrm{adm}_i(D)$ indexed by a datum $D$ (and, for one of them, by four constants
$C_{\Lambda},C_k,C_{rk},C_{\deg}$), but in each of them the datum is a dummy argument: the statements do not mention
$D$. We write $\sfname{Adm}(D,C_\Lambda,C_k,C_{rk},C_{\deg})$ for their conjunction. Several clauses
carry names which do not describe their content; \cref{tab:pf:linkG:clauses} gives the content. In the table,
$\OO$ is a discrete valuation ring over $\Z_p$ with a chosen uniformiser, $K=\mathrm{Frac}\,\OO$,
``$\GL_n$-representation'' means a finite-dimensional representation $(E_{ab})$ of $\mathfrak{gl}_n$ over $K$, a
\emph{$U_\Z$-lattice} is a finitely generated $\OO$-submodule stable under the generators of the integral form
$U_\Z$ (the divided powers $E_{ab}^{(k)}$), and $N_0(n)$, $d_p(n,\ell)$
are explicit exponents depending on $n$, $p$ and a length $\ell$.

{\small
\begin{longtable}{>{\raggedright\arraybackslash}p{0.05\textwidth}>{\raggedright\arraybackslash}p{0.29\textwidth}>{\raggedright\arraybackslash}p{0.56\textwidth}}
\caption{The clauses of the criterion.}\label{tab:pf:linkG:clauses}\\
\toprule
& formal name & content\\
\midrule
\endfirsthead
\toprule
& formal name & content\\
\midrule
\endhead
1 & $\mathrm{adm}_1$ & orbit positions: for $TT$ satisfying the orbit hypothesis, the $3r$ points of a marked tile are
even and far from the other points of $TT$; consecutive marked tiles are joined by a stretch of length $<G'$
containing at most $N_1$ points; fewer than $r$ even points of $TT$ lie at forward distance $1,\dots,L_1$ from any
point; every even point has a next even point at distance at most $G'$ with at most $N_1$ points in between\\
2 & $\mathrm{adm}_2$ & orbit counts: for $0<\eps<1/52$, $h=2+8\eps$, $r=\lceil8(h+2)/\eps\rceil+1$ and $\delta=x^{-\eps/8}$,
eventually in $x$: $G'\le C\log x/\delta$, $N_1\le C\log x$, and for $n'\ge x^{\eps/4}$: $3G'\le n'$ and every $TT$
with the tile condition has $n'\delta/(C\log x)\le|TT|\le40n'\delta$\\
3 & $\mathrm{adm}_3$ & $\sfname{DualHeightInclusion}$: $\linkGH_{e_s}\subseteq\linkGR_e$ and $\linkGH_{e_t}^\vee\subseteq\linkGR'_e$ for
every edge\\
4 & \textsf{Archimedean\-Signature\-Matching} & $(\mathrm P)$ (\cref{pf:linkG:perfect})\\
5 & $\mathrm{adm}_5$ & $(\mathrm H)$ (\cref{pf:linkG:perfect})\\
6 & $\mathrm{adm}_6$ & the extremal projector statement $(\mathrm E)$\\
7 & $\mathrm{adm}_7$ & invariant projections: for a $\GL_n$-representation with a $U_\Z$-lattice $M$, the invariant projection
equals $P_0\circ\pi_0$ and $p^{N_0d_p(n,\ell(M))}$ times it preserves $M$\\
8 & $\mathrm{adm}_8$ & $\nabla$ versus $\Delta$: for an irreducible module with highest weight vector $v$ of weight $\lambda$,
$p^{N_0d_p(n,\ell(\lambda))}\nabla(\lambda,v)\subseteq\Delta(v)$, and every $U_\Z$-lattice $M$ with
$M\cap V_\lambda=\OO v$ satisfies $\Delta(v)\subseteq M\subseteq\nabla(\lambda,v)$\\
9 & $\mathrm{adm}_9$ & denominators of isotypic projectors on $U_\Z$-lattices\\
10 & $\mathrm{adm}_{10}$ & translation corollary: an isotypic $Y\in M$ all of whose integral tests are divisible by $p^c$ lies in
$p^{c-N_0d_p(n,\ell(\lambda))}M$\\
11 & $\mathrm{adm}_{11}$ & filtered version of 9 on polynomial rings with commuting families of isotypic projectors\\
12 & \textsf{Filtered\-Projection\-Criterion} & divided powers of the $\mathfrak{gl}_n$-operators on polynomial rings, their
adjointness for the Fock pairing, degree bounds for isotypic projectors and their products\\
13 & $\mathrm{adm}_{12}$ & properties of the valuations $\mathrm{val}_S$, $\lambda_{1,S}$ of polynomials (additivity, and their
description by the weights of the isotypic components)\\
14 & $\mathrm{adm}_{13}$ & the polynomials $\Psi(P)$ attached to Gelfand--Tsetlin paths are highest weight vectors and linearly
independent\\
15 & $\mathrm{adm}_{14}$ & the $\Psi(P,I)$ for standard pairs $(P,I)$ form a basis of the polynomial space\\
16--21 & $\mathrm{adm}_{15},\dots,\mathrm{adm}_{20}$ & statements on Gelfand--Tsetlin paths of $\GL_{r+K}$: a rank-one statement, the
splitting of paths into paths of $\GL_r$ and $\GL_K$ with an interlacing condition, lowest weighted-homogeneous components
of column polynomials, saturation of joint types, dependence of Witt-lifted substitutions only on the rational plane,
and properties of canonical translates\\
22 & $\mathrm{adm}_{21}$ & determinant divisibility: $\mathrm{val}_S(G)\ge\theta>\deg_{\mathrm{other}}$ implies
$\det_A^{\,\theta-\deg_{\mathrm{other}}}\mid G$\\
23 & $\mathrm{adm}_{22}$ & $p$-adic unit products: for $P=p^{t+1}$, the numbers $\prod_{i<n}\prod_{1\le y<P}(1+Pi/y)$, their
$p$-adic logarithms, valuations of $P^s\sum_{y<P}y^{-s}$, Faulhaber polynomials, decay of Mahler coefficients and
polynomial approximation of products of functions with decaying Mahler coefficients\\
24 & $\mathrm{adm}_{23}$ & a factorial unit decomposition and a degree truncation lemma\\
25 & \textsf{Frobenius\-Expansion\-Congruence} & a derivative valuation bound and digit expansions of derivatives\\
26 & $\sfname{HerrExtensionClause}$ & $\mathrm{val}_S$ and $\lambda_{1,S}$ on $\OO[x]$: additivity, values on variables, and
$\mathrm{val}_S=1$ on determinants of distinct vector letters\\
27 & $\mathrm{adm}_{26}$ & the symbol space statement with constants $(C_\Lambda,C_k,C_{rk},C_{\deg})$: for valid frames, free symbol
spaces with complements, of rank at least $(1-C_{rk}|TT|r^2k/Y)\prod_t\binom{d_t+k-1}{k-1}p^{-n'rK}$, generated in degrees
bounded in terms of $C_{\deg}$\\
28 & \textsf{Herr\-Cocycle\-Extension\-Class} & $\sfname{WittIdent}$: $W$ is isomorphic to the Witt vectors of $\F_{p^{n'}}$,
compatibly with $\sigma$\\
29 & \textsf{Fontaine\-Formal\-Properties} & $\sfname{FlatTop}$: for a flat weight, an $\OO$-isomorphism $\Theta$ from a space of
polynomials onto $\linkGV$ carrying integrality to $\linkGH_L$ and the filtration of Witt-lifted substitutions to the
scaled lattices $\varpi^i\linkGR_e$\\
30 & \textsf{Ambient\-Metric\-Lower\-Bound} & $\sfname{DualHeightBasis}$: for every edge a basis $(x_i)$ of $\linkGV$ and integers
$\pi_i$ with $\linkGH_{e_s}=\langle x_i\rangle$, $\linkGH_{e_t}=\langle p^{\pi_i}x_i\rangle$,
$\linkGH_e=\langle\varpi^{\pi_i}x_i\rangle$, $A_r-\pi_i\ge0$, and explicit bases of $\linkGR_e$, $\linkGR'_e$\\
31 & $\sfname{HostOrbitConditions}$ & $\sfname{HLSpans}$: $\linkGH_L$ is spanned by a $K_f$-basis of $\linkGV$\\
32 & $\sfname{OrderConsequences}$ & $\sfname{TestSide}$: duality of edges, and an isomorphism $\linkGV^\vee\cong\linkGV$
comparing $\linkGH_L^\vee$ with $\linkGH_{L^\vee}$ and $\linkGR'_e$ with $\linkGR_{e^\vee}$ up to $p^{\kappa_1}$\\
33 & $\sfname{TopsChoiceLemma}$ & segment rings of frames: saturation, a basis of products of path polynomials, and valuation
bounds\\
34 & $\sfname{ZeroGainAbove}$ & the tile condition gives a valid frame for $TT$ with prescribed positions in the marked tiles\\
35 & $\sfname{ChosenTopsCount}$ & $\mathrm{val}_S$, $\lambda_{1,S}$ commute with renaming and with $\OO\subseteq K$\\
36 & $\sfname{FewBadOrbits}$ & Weyl's dimension formula
$\dim\linkGV=\prod_s\prod_{i<j}\frac{\lambda_s(i)-\lambda_s(j)+j-i}{j-i}$\\
37 & $\sfname{LevelPrimesProviso}$ & products of sections: $F\in\bigotimes_sH^0(\lambda_s,\Lambda_s)$ and
$G\in\bigotimes_sH^0(\mu_s,\Lambda_s)$ imply $FG\in\bigotimes_sH^0(\lambda_s+\mu_s,\Lambda_s)$\\
\bottomrule
\end{longtable}
}

\begin{definition}[The criterion]\label{pf:linkG:criterion}
For a datum $D$, $\mathsf{Concl}(D)$ means: there is $C_0>0$ such that $\linkGLL(C)$ holds for all $C\ge C_0$.
$\mathsf{Crit}(D)$ means: for all $C_\Lambda,C_k,C_{rk},C_{\deg}$, if
$\sfname{Adm}(D,C_\Lambda,C_k,C_{rk},C_{\deg})$ holds and two auxiliary existence statements hold (a ``source
cleaning'' statement with positive constants $C'_\Lambda,C_D,C_{sz}$, and a ``plane Gram'' statement with constants
$C_H,C_E,C_6,C_{18}$), then $\mathsf{Concl}(D)$ holds. $\sfname{Loc}$ is the statement that
$\mathsf{Crit}(D)$ holds for every datum $D$.
\formal{Formalism.MainNotions.tableLineDegreeFive.ConclusionE}
\formal{Formalism.MainNotions.G.SecondTheoremReduced}
\formal{Formalism.L2.G.Loc}
\end{definition}

\begin{definition}[Symbol realisability and the cloud step]\label{pf:linkG:cloud}
Let $\varphi:W(\F_{p^{n'}})\cong W$ be a Witt identification compatible with $\sigma$ (clause 28). A \emph{frame} for
$TT$ is the combinatorial datum of clause 34: a choice, for every integer $j$, of $r$ positions in the orbit $\Z/n'$ (in
the marked tiles), with letters, distances and ``oldest'' positions satisfying the validity conditions of clause 1. A
frame determines a polynomial ring $R_F$ (in variables indexed by letters), \emph{symbol spaces} in $R_F$ (clause 27),
and, for a Witt matrix $g\in\GL_k(W(\F_{p^{n'}}))$, a substitution $\mathrm{subst}_g$ of polynomials on $\Z/n'\times k$
variables and a ``symbol map'' $\mathrm{sym}_g$ on $R_F$.
\begin{enumerate}[label=(\alph*)]
\item $\mathsf{Real}(C,C_\Lambda)$: in every unramified quadratic setting with the hypotheses of
$\linkGES$ (constant $C$), for every Witt identification, every decomposition $k=r+K$ and every valid frame, every
element of the symbol space is \emph{realisable}: for every $g$, its image under $\mathrm{sym}_g$ is
$\mathrm{subst}_g$ of a polynomial of bounded degree, with explicit bounds.
\item $\mathsf{Cloud}(C,C_t,C_D,C_\Lambda)$: under the same hypotheses, for a linearly independent family $(B_i)$
of generators of the symbol spaces whose span has a complement, every edge $e$ with target $W^k$ and every edge basis of
$e$ of the form $h\,(\varphi(g)^{-1})^{T}$, and every family of polynomials $F^{\mathrm s}_i$ realising the $B_i$, the
\emph{cloud statement} holds: there are an isomorphism $\Theta$ as in clause 29, an integer $n_2$, \emph{shift vectors}
$\mathrm{sh}_1,\dots,\mathrm{sh}_{n_2}$ of the complementary weight, integers $o_S,t_T,D_1$ with
$o_S\le\omega$, $t_T\le t_0+|TT|c_0+\omega$, $D_1\le C_D|TT|k^2\log_pa$, the count stated below, and test functionals for all primitive
combinations of the products of the realisations with the shifts. The count reads
$(1-kc_0/a)^{|TT|}\dim\linkGV\le n_2\prod_{t\in TT}\binom{d_t+k-1}{k-1}$.
\end{enumerate}
\formal{Formalism.MainNotions.Pf14.SymbolRealisableBound}
\formal{Formalism.MainNotions.Pf14.CloudStep}
\end{definition}

The edge step is obtained from these two statements.

\begin{proposition}[The edge step]\label{pf:linkG:edgestep}
Let $D$ be a datum and assume $\sfname{Adm}(D,C_\Lambda,C_k,C_{rk},C_{\deg})$. Assume that there are $C_2,C_3$, $C_t,C_D>0$
such that $\mathsf{Real}(C,C_\Lambda)$ holds for $C\ge C_2$ and $\mathsf{Cloud}(C,C_t,C_D,C_\Lambda)$
holds for $C\ge C_3$. Then there are $C'_t,C'_D,C_r,C_1>0$ such that $\linkGES(C,C'_t,C'_D,C_r)$ holds for all
$C\ge C_1$; one can take $C'_t=\max(C_t,|C_{\deg}|)$, $C'_D=C_D$ and $C_r=\max(C_{rk},1)$.
\formallabel{Formalism.MainNotions.Pf14.edgeStep_of}
\formallabel{Formalism.MainNotions.Pf14.restEdge_of}
\formallabel{Formalism.MainNotions.Pf14.norm_index_count_bound}
\end{proposition}

\begin{proof}[Proof (construction of the edge families)]
Let $C\ge C_1=\max(C_g,C_2,C_3,2,|C_\Lambda|,C_D)$, where $C_g$ is the constant of the symbol space statement
(\lref{Formalism.MainNotions.Pf14.symbol_space_at}, from clause 27). Fix data with the hypotheses of $\linkGES$ and an
edge $e$ with target $W^k$; write $k=r+K$ with $K\ge1$.
\begin{enumerate}[label=(\arabic*)]
\item By clause 28 choose a Witt identification $\varphi$. By the orbit hypothesis there is a valid frame $F$ for $TT$
(\lref{Formalism.MainNotions.Pf14.frames_of}).
\item The symbol space statement gives a free $\OO$-submodule $S_g$ of $R_F$ with a complement, of rank
\[
\rank S_g\ \ge\ \Bigl(1-\frac{C_{rk}|TT|r^2k}{Y}\Bigr)\prod_{t\in TT}\binom{d_t+k-1}{k-1}\,p^{-n'rK},
\]
generated by elements of bounded degree; choose a basis $(B_i)$ of $S_g$ consisting of generators
(\lref{Formalism.MainNotions.Pf14.basis_fin}). Here $d_t$ is the degree of the flat part of $\lambda_0$ at $t$.
\item Every edge with target $W^k$ has an edge basis of the form $h\,(\varphi(g)^{-1})^{T}$ with $g$ a Witt matrix
(\lref{Formalism.MainNotions.Pf14.edgePlane}); fix one for $e$.
\item By $\mathsf{Real}$ every $B_i$ is realised: $\mathrm{sym}_g(B_i)=\mathrm{subst}_g(\sum_{o\le X}F^{\mathrm
s}_{i,o})$ with graded pieces of degree $o\le X$, where
$X=C_{\deg}|TT|\bigl(C_\Lambda N_s^2\log a\cdot a/(pk)+a\,p^{-L_1}+k\,p^{G'}\bigr)$ and $N_s=\min(N_1,2|TT|)$.
\item $\mathsf{Cloud}$ gives $\Theta$, the shift vectors $\mathrm{sh}_l\in\linkGR_e\cap\varpi^{-o_S}\linkGH_{e_s}$ of the
complementary weight ($l\le n_2$), the integers $o_S,t_T,D_1$, the count for $n_2$, and test functionals for all
primitive combinations of the products $S_{i,l}=\Theta(\text{realisation of }B_i)\cdot\mathrm{sh}_l$.
\item Each graded piece of a realisation is lifted to a vector $F^{\mathrm v}_{i,o}\in\varpi^{-o}\linkGH_{e_s}$ of the
flat weight, whose difference with the corresponding $\Theta$-image lies in $\varpi^{c'}\linkGR_e$, and which lies in
$\varpi^{c'}\linkGR_{e'}$ for the other edges $e'$ with the same source
(\lref{Formalism.MainNotions.Pf14.flatSrc}); here $c'$ is an explicit integer with $c'\ge c_*$.
\item The edge family consists of the products $F_{i,l}=\bigl(\sum_{o\le X}F^{\mathrm
v}_{i,o}\bigr)\cdot\mathrm{sh}_l$, which lie in $\linkGV(\lambda_0)$ because $\lambda_0$ is the sum of its flat part and the
complementary weight; $F_{i,l}-S_{i,l}\in\varpi^{c'}\linkGR_e$. By clause 37 (products of sections), clause 3 and
$(\mathrm P)$, these vectors satisfy $\linkGEF(\lfloor X\rfloor+o_S,t_T,D_1,c_*,e,n)$ with
$n=\rank S_g\cdot n_2$ (\lref{Formalism.MainNotions.Pf14.edgeFamOf}).
\end{enumerate}
It remains to check the bounds of $\linkGES$. Since $C\ge|C_\Lambda|$ and $N_s\le N_1$, $X\le t_0$ for the constant
$C'_t=\max(C_t,|C_{\deg}|)$, so $\lfloor X\rfloor+o_S\le t_0+\omega$; also $t_T\le t_0+|TT|c_0+\omega$ (the bound
for $t_0$ is monotone in the constant $C_t$) and $D_1\le C_D|TT|k^2\log_pa$. Finally, by the two counts,
\[
n=\rank S_g\cdot n_2\ \ge\ \Bigl(1-\frac{C_r|TT|r^2k}{Y}\Bigr)\Bigl(1-\frac{kc_0}{a}\Bigr)^{|TT|}
\frac{\dim\linkGV}{p^{n'r(k-r)}},
\]
because the products $\Pi=\prod_t\binom{d_t+k-1}{k-1}$ of binomial coefficients cancel and $C_{rk}\le C_r$: with
$\beta=|TT|r^2k/Y$, $u=(1-kc_0/a)^{|TT|}\ge0$ and $Q=p^{n'r(k-r)}$, if $1-C_r\beta\ge0$ then
$(1-C_r\beta)u\dim\linkGV\le(1-C_r\beta)\Pi n_2\le(1-C_{rk}\beta)\Pi n_2\le Q\rank S_g\,n_2$, and if $1-C_r\beta<0$ the
left side is negative. In the development the constants $C_2,C_3,C_t,C_D$ are provided by two hypotheses (realisability
of the sources, and the cloud statement), which hold under the clauses and the two auxiliary statements of the
criterion; the statement is otherwise the same.
\end{proof}

\begin{theorem}[The criterion]\label{pf:linkG:loc}
$\sfname{Loc}$ holds: for every datum $D$ and all constants, $\sfname{Adm}(D,C_\Lambda,C_k,C_{rk},C_{\deg})$ implies
$\mathsf{Concl}(D)$.
\formallabel{Formalism.L2.G.loc_proved}
\formallabel{Formalism.MainNotions.TraceSpectra.stmtE}
\formallabel{Formalism.MainNotions.G.criterion_restricted_of_closed}
\formallabel{Formalism.MainNotions.PairMembershipBlocks.stmtE_of_chain}
\formallabel{Formalism.MainNotions.TestSideCounting.stmtE_of_OpNS}
\formallabel{Formalism.MainNotions.TestSideCounting.theorem_II_closed_of_test_bound}
\formallabel{Formalism.MainNotions.Pf14.theorem_II_of_cloud_bound}
\formallabel{Formalism.MainNotions.Pf14.prvE_degree_five_of_adm_sources_realisable}
\formallabel{Formalism.MainNotions.Pf14.theorem_II_of_edge_bound}
\end{theorem}

\begin{proof}
We prove the stronger statement that $\sfname{Adm}$ alone implies $\mathsf{Concl}(D)$ (the auxiliary
statements in $\mathsf{Crit}$ are not needed). The formal proof is a chain of
reductions, each of which replaces the conclusion by explicit intermediate statements; we describe its structure.
\begin{enumerate}[label=(\arabic*)]
\item By \cref{pf:linkG:edgeassembly} and \cref{pf:linkG:edgestep}, $\mathsf{Concl}(D)$ follows from the clauses
(which contain $(\mathrm P)$ and $(\mathrm H)$), the source cleaning statement, the plane Gram statement, the
realisability of the sources (the bound $\mathsf{Real}$ for large constants) and the cloud statement (the bound
$\mathsf{Cloud}$ for large constants). The realisability of the sources is \cref{pf:linkG:srcreal}.
\item Under the clauses, the source cleaning statement holds with some constant $C\ge1$
(\lref{Formalism.MainNotions.TestSideCounting.cleaning_of_admitted}) and the plane Gram statement holds
(\lref{Formalism.MainNotions.TestSideCounting.plane_gram_of_admitted}); the source primitive integrality statement
holds by \cref{pf:linkG:prim}. The cloud statement is reduced, through two intermediate forms
(\lref{Formalism.MainNotions.Pf14.restCloud_of_X}, \lref{Formalism.MainNotions.Pf14.restCloudX_of_Y}, and \lref{Formalism.MainNotions.Pf14.restCloudY_of}:
reformulations which assemble the shift vectors, the cloud bound and the flatness bound into the statement $\mathsf{Cloud}$),
to the source primitive integrality and a statement of ``remaining inputs'' with the constants $(C,C+1,C,26,2,24,60)$.
\item The remaining inputs are obtained from two halves
(\lref{Formalism.MainNotions.TestSideCounting.test_bound_of_two_halves}: the remaining inputs are the conjunction of a
cloud bound and a flatness bound): the cloud bound $\mathsf{Cl}$ for large constants (\cref{pf:linkG:cloudb}), which
follows from $\mathsf{TO}$ (\cref{pf:linkG:tsop}) and the tensor operator statement (\cref{pf:linkG:tenop}); and the
flatness bound $\mathsf{Flat}$ for large constants (\cref{pf:linkG:flat}), which follows from $\Xi$ (\cref{pf:linkG:xi}),
which follows from $\mathsf{NS}$ for large constants. The latter follows from a membership statement and $\mathsf{P5}$ with the
constants $(26,2,24,60)$ (\cref{pf:linkG:prop5}) by \lref{Formalism.MainNotions.TestSideWords.restNS_of} (the membership
statement says that every $G$ as in $\mathsf{NS}$ satisfies $\mathsf{Mem}$).
\item The membership statement is proved from the word statement (\cref{pf:linkG:word}) and a monomial statement
(\lref{Formalism.MainNotions.TestSideWords.memCondition_of_theoremII}: the symbol polynomials satisfy
$\mathsf{Mem}(0,\frac18)$ by the monomial statement, words map this into $\mathsf{Mem}(\frac18,\frac38)$, and this gives $\mathsf{Mem}$
for the elements of the operator closure), using the one-operator
steps (\lref{Formalism.MainNotions.LetterSubstitution.noOpStep}), the plane coefficient chain
(\lref{Formalism.MainNotions.PairMembershipBlocks.restPl_of_restChain_of_restLink} and \lref{Formalism.MainNotions.TraceSpectra.restChain}:
valuation bounds for the plane coefficients of the level images, obtained from a chain statement and the level link
\cref{pf:linkG:link}) and an arithmetic condition valid for all constants $C\ge1$
(\lref{Formalism.MainNotions.PairMembershipBlocks.arithCondition_of_one_le}).
\end{enumerate}
Combining (1)--(4) gives $\mathsf{Concl}(D)$ for every $D$ under $\sfname{Adm}$, which is $\sfname{Loc}$.
\end{proof}

The steps (2)--(4) concern valuations of coefficients of polynomials in the Fock model of the tensor product over the
Frobenius orbit; the statements used in them are proved in \cref{sec:pf:linkG:inner}.

% ---------------------------------------------------------------------------
\subsection{The inner reductions of the criterion}\label{sec:pf:linkG:inner}

We now prove the statements used in steps (1)--(4) of the proof of \cref{pf:linkG:loc}. They concern the \emph{test
side}: polynomials in the variables $X_{t,i}$ ($t\in\Z/n'$, $1\le i\le k$), their images in the Fock model of the tensor
product over the Frobenius orbit of the dual Weyl modules of $\GL_k$, and valuations of pairings. Throughout, $D$ is a datum
and $C_\Lambda,C_k,C_{rk},C_{\deg}$ are constants.

\begin{definition}[Test-side data]\label{pf:linkG:testdata}
Let $C\ge0$. \emph{Test-side data with constant $C$} consist of: an unramified quadratic setting and parameters
$m,r,TT,L_1,L_2,\delta,h,x,\lambda_0$, $a_T,b_T,a,Y,c_{pp},c_0,D_p,c_*$ satisfying the hypotheses of $\linkGES$ with
constant $C$ (so $k=m+1$, $r\ge2$ and $\lambda_0$ is dominant); a Witt identification $\varphi$ compatible with $\sigma$; a
decomposition $k=r+K$ with $K\ge1$; a valid frame $F$ for $TT$ with the orbit parameters $N_1,G'$ of $(r,\delta,h,x)$;
degree bounds $E'_t\le\min(E_0,d_t)$ for $t\in TT$, where $d_t$ is the degree of the flat part of $\lambda_0$ at $t$ and
$E_0$ is the exponent bound of $\linkGES$; and a linearly independent family $(B_b)_{b<n_b}$ of generators of the symbol
spaces of clause 27 (for the multidegree $d-E'$) whose span has a complement. \emph{Edge data} consist in addition of: a
Hecke edge $e$ of type $r$ with target $W^k$; matrices $h\in\GL_k(Q_q)$ and $g$ (a Witt matrix) such that
$h(\varphi(g)^{-1})^T$ is an edge basis of $e$; \emph{realisations} $F^{\mathrm s}_b$, polynomials of multidegree $d_{TT}$
($d_{TT}=d$ on $TT$ and $0$ elsewhere) with $\mathrm{subst}_g(F^{\mathrm s}_b)=\mathrm{sym}_g(B_b)$; and a true theta
isomorphism $\Theta$ (clause 29) from the polynomials of multidegree $d_{TT}$ onto $\linkGV(\lambda_{\mathrm{fl}})$, where
$\lambda_{\mathrm{fl}}$ is the flat weight with degrees $d_{TT}$.
\formal{Formalism.MainNotions.TestSideWords.DataII}
\end{definition}

We write $\lambda_{\mathrm{rest}}=\lambda_0-\lambda_{\mathrm{fl}}$ (the complementary weight); products $y\beta$ of
$y\in\linkGV(\lambda_{\mathrm{fl}})$ and $\beta\in\linkGV(\lambda_{\mathrm{rest}})$ lie in $\linkGV(\lambda_0)$. For a polynomial
or a module of polynomials $F$, the \emph{operator closure} $\mathrm{Op}_{c_0}(F)$ is the span of the $w(F)$, $w$ a word in
the polarisation operators $E^{(t)}_{ij}=X_{t,i}\,\partial/\partial X_{t,j}$ with $t\in TT$ and at most $c_0$ letters at
every position $t$. A \emph{primitive combination} of a family $(S_x)$ is a sum $\sum_xc_xS_x$ with $c_x\in\OO$, some $c_x$
a unit. We write $b(u,T)=T(p^{A_r(\lambda)}u)$ for the normalised pairing, $\langle\,,\rangle$ for the Fock pairing,
$\kappa_{\mathrm{fl}}=\sum_sk\lfloor\log_pd_s\rfloor$, and $\delta_D=C_D|TT|k^2\log_pa$. For a dominant weight $\lambda$ let
$\ell_s(\lambda)$ be the difference between the largest and the smallest entry of $\lambda_s$, and
\[
\kappa_1(\lambda)=\sum_{s:\ \ell_s(\lambda)\ne0}\binom k2\bigl\lfloor\log_p(\ell_s(\lambda)+k-1)\bigr\rfloor .
\]
A \emph{duality with loss $\kappa_1(\lambda)$} is a linear isomorphism $\Psi:\linkGV(\lambda)^\vee\to\linkGV(\lambda)$ with
$\Psi(\linkGH_L^\vee)\subseteq\linkGH_{L^\vee}$ and $p^{\kappa_1(\lambda)}\linkGH_{L^\vee}\subseteq\Psi(\linkGH_L^\vee)$ for every
lattice $L$, and $\Psi(\linkGR'_e)\subseteq\linkGR_{e^\vee}$ and $p^{\kappa_1(\lambda)}\linkGR_{e^\vee}\subseteq\Psi(\linkGR'_e)$
for every Hecke edge $e$ of type $r$ ($L^\vee$ and $e^\vee$ are the dual lattice and the dual edge).

\begin{definition}[The test-side statements]\label{pf:linkG:teststeps}
Let $C,C_t,C_D,C_K\ge0$. Each statement below quantifies over all test-side data with constant $C$, and (a)--(e) also over
all edge data. Let $t_0$ be the bound of $\linkGES$ with the constants $C,C_t$, and $c'$ the integer of $\linkGES$.
\begin{enumerate}[label=(\alph*)]
\item $\mathsf{Prim}(C)$: there are $n_2$ shift vectors $\mathrm{sh}_l\in\linkGR_e(\lambda_{\mathrm{rest}})$, $n_2$ the number of
weight vectors of the complementary weight, such that no primitive combination of the products
$S_{b,l}=\Theta(F^{\mathrm s}_b)\,\mathrm{sh}_l$ lies in $\varpi\linkGR_e(\lambda_0)$.
\item $\mathsf{TO}(C)$: there are dualities $\Psi_0,\Psi_{\mathrm{fl}}$ for $\lambda_0^{TT}$ and $\lambda_{\mathrm{fl}}$ with losses
$\kappa_1(\lambda_0^{TT})$, $\kappa_1(\lambda_{\mathrm{fl}})$ ($\lambda^{TT}$ is $\lambda$ on $TT$ and $0$ elsewhere) such that for all
$\beta,s\in\linkGV(\lambda^{TT}_{\mathrm{rest}})$ and every $F$ of multidegree $d_{TT}$ there is $F'\in\mathrm{Op}_{c_0}(F)$ with
$p^{A_r(\lambda^{TT}_{\mathrm{rest}})+\kappa_1(\lambda_0^{TT})}\,\Psi_0^{-1}\bigl(\Psi_{\mathrm{fl}}(T)\beta\bigr)\bigl(\Theta(F)s\bigr)
=T(\Theta(F'))$ for every $T\in\linkGV(\lambda_{\mathrm{fl}})^\vee$.
\item $\mathsf{Cl}(C,C_K)$: there is $\kappa\le C_K|TT|k^2\log_pa$ such that for all shift vectors
$\mathrm{sh}_l\in\linkGR_e(\lambda_{\mathrm{rest}})$ and every combination $u=\sum c_xS_x\notin\varpi\linkGR_e(\lambda_0)$ of the
products $S_{b,l}$ there are $F_1\in\mathrm{Op}_{c_0}(\{F^{\mathrm s}_b\})$ with $\Theta(F_1)\in\linkGR_e(\lambda_{\mathrm{fl}})\smallsetminus
\varpi\linkGR_e(\lambda_{\mathrm{fl}})$, a linear map $\Lambda:\linkGV(\lambda_{\mathrm{fl}})^\vee\to\linkGV(\lambda_0)^\vee$ and
$c\in\OO$ such that $c\,b(u,\Lambda T)=p^\kappa\,b(\Theta(F_1),T)$ for all $T$; $\Lambda$ maps $\linkGR'_e(\lambda_{\mathrm{fl}})$
into $\linkGR'_e(\lambda_0)$, with $b(u,\Lambda T)\in\OO$ there; and $\Lambda$ maps $\varpi^i\linkGR'_{e'}(\lambda_{\mathrm{fl}})$ into
$\varpi^{i-\ell}\linkGR'_{e'}(\lambda_0)$ for every Hecke edge $e'$ of type $r$ with the target of $e$, and
$\varpi^i\linkGH_{e_t}(\lambda_{\mathrm{fl}})^\vee$ into $\varpi^{i-\ell}\linkGH_{e_t}(\lambda_0)^\vee$ ($i\in\Z$), where
$\ell=r|TT|c_0+2rn'c_{pp}$.
\item $\mathsf{Flat}(C,C_t,C_D)$: there are $t_F\le t_0+|TT|c_0$ and $D_F\le\delta_D$ such that for every
$F_1\in\mathrm{Op}_{c_0}(\{F^{\mathrm s}_b\})$ with $\Theta(F_1)\in\linkGR_e(\lambda_{\mathrm{fl}})\smallsetminus\varpi\linkGR_e(\lambda_{\mathrm{fl}})$
there is $T\in\varpi^{-t_F}\linkGH_{e_t}(\lambda_{\mathrm{fl}})^\vee\cap\linkGR'_e(\lambda_{\mathrm{fl}})$ with
$T\in\varpi^{c'}\linkGR'_{e'}(\lambda_{\mathrm{fl}})$ for every other Hecke edge $e'\ne e$ of type $r$ with the target of $e$, and
$v_p\bigl(b(\Theta(F_1),T)\bigr)\le D_F$.
\item $\Xi(C,C_t,C_D)$: there are $t_F\le t_0+|TT|c_0$ and $D_X\le\delta_D$ such that for every $F_1$ as in (d) there
are polynomials $\sigma_1,\dots,\sigma_q$ of multidegree $d_{TT}$, each a \emph{one-sided element} of some order $o_i\le
t_F$ for the frame (a test polynomial of the shape used in the one-sided lemma of the development), with
$v_p\bigl(p^{\kappa_{\mathrm{fl}}}\langle\mathrm{subst}_g(F_1),\sum_i\sigma_i\rangle/\prod_sd_s!\bigr)\le D_X$.
\item $\mathsf{NS}(C,C_D)$ with constants $(C_H,C_E,C_6,C_{18})$: there is $D_X\le\delta_D$ such that every $G$ in the operator
closure of the span of the symbol polynomials $\mathrm{sym}_P(B_b)$, of multidegree $d_{TT}$, with integral coefficients
not all in $\varpi\OO$, satisfies $\mathsf{NS}(G)$: there are signed chain data (with constant $C_H$) and a dual element
$F_2$ of vanishing $B$-valuation such that for every admissible perturbation $G_p+R$ of the plane coefficients
($G_p$ integral in the $Y$-space with constants $C_6,C_{18}$, $R$ with coefficients of valuation $\ge c-E(C_E)$, the
$B$-pairing with $F_2$ of controlled valuation) there are one-sided elements $\sigma_i$ of multidegree $d_{TT}$ with
$v_p\bigl(p^{\kappa_{\mathrm{fl}}}\langle G,\sum_i\sigma_i\rangle/\prod_sd_s!\bigr)\le D_X$.
\item $\mathsf{Word}(C)$: if a polynomial $F$ satisfies the membership condition $\mathsf{Mem}(0,\frac18)$, then $w(F)$
satisfies $\mathsf{Mem}(\frac18,\frac38)$ for every word $w$ in the operators $E^{(t)}_{ij}$, $t\in TT$, with at most $c_0$
letters at each position. Here $\mathsf{Mem}(s_c,s_L)$ means: for all signed chain data of the frame which satisfy the
twisted substitution growth bound with the parameters $(s_c,s_L)$, the plane coefficients of the projected level images
of $F$ lie in the $X$-space of the chain data; $\mathsf{Mem}$ is the analogous condition with the tie conditions of the
frame in place of the growth bound.
\item $\mathsf{P5}(C,C_D)$ with constants $(C_H,C_E,C_6,C_{18})$: there is $D_X\le\delta_D$ such that every $G$ as in (f)
which satisfies $\mathsf{Mem}$ satisfies $\mathsf{NS}(G)$.
\end{enumerate}
For each statement $\mathsf X$, ``$\mathsf X$ for large constants'' means: $\sfname{Adm}(D,C_\Lambda,C_k,C_{rk},C_{\deg})$ implies that
there are positive constants (those of $\mathsf X$ other than $C$, and $C_1$) such that $\mathsf X$ holds for all $C\ge C_1$.
\formal{Formalism.MainNotions.Pf14.PrimStep}
\formal{Formalism.MainNotions.TestSideCounting.TSOpStep}
\formal{Formalism.MainNotions.TestSideCounting.CloudBStep}
\formal{Formalism.MainNotions.TestSideCounting.FlatStep}
\formal{Formalism.MainNotions.TestSideCounting.XiStep}
\formal{Formalism.MainNotions.TestSideCounting.NSStep}
\formal{Formalism.MainNotions.TestSideWords.WordConditionAt}
\formal{Formalism.MainNotions.TestSideWords.Proposition5At}
\end{definition}

\paragraph{The sources.}

\begin{proposition}[Realisability of the sources]\label{pf:linkG:srcreal}
Assume $\sfname{Adm}(D,C_\Lambda,C_k,C_{rk},C_{\deg})$ and the source cleaning statement. Then $\mathsf{Real}(C,C_\Lambda)$ holds for
every $C\ge30$.
\formallabel{Formalism.MainNotions.Pf14.sources_realisable_holds}
\formallabel{Formalism.MainNotions.Pf14.symbol_realisable_general_rank}
\end{proposition}

\begin{proof}
Clauses 16 and 17 give the \emph{generator span inclusion} and clauses 18, 19 and 21 the statement that the image of a
generator under the symbol map is a symbol
(\lref{Formalism.MainNotions.G.symbolGeneratorData.generator_pair_conditions_of_admitted} and \lref{Formalism.MainNotions.G.symbolGeneratorData.generator_path_homogeneous_holds});
the monomial statement holds unconditionally (\lref{Formalism.MainNotions.SelectionLemmas.Pf.restMon}). Let $C\ge30$
and fix test-side data without the family $(B_b)$. By \lref{Formalism.MainNotions.Pf14.symbol_realisable_iff_general_rank}
(a reformulation), it suffices to find degree bounds $E'$ with $E'_t\le\min(E_0,d_t)$ on $TT$ such that for every Witt
matrix $g$, every $P$ in the module generated by the symbol spaces with multidegree $d-E'$ and every $T$ bounding the
degrees of the monomials of $P$, there are polynomials $F^{\mathrm s}_o$ ($o\le T$) of multidegree $d_{TT}$ such that each
$\mathrm{subst}_g(F^{\mathrm s}_o)$ has integral coefficients and is realisable of order $o$, and
$\sum_{o\le T}\mathrm{subst}_g(F^{\mathrm s}_o)=P\cdot\mathrm{subst}_g(\mathrm{SP})$, where $\mathrm{SP}$ is the true symbol
polynomial. Take for $E'_t$ the degree $e_{\mathrm{SP}}(t)$ of $\mathrm{SP}$ at $t$; it vanishes off $TT$, and for $C\ge30$ the
hypotheses of $\linkGES$ give $e_{\mathrm{SP}}(t)\le E_0$ and $e_{\mathrm{SP}}(t)\le d_t$
(\lref{Formalism.MainNotions.Pf14.spdeg}: a numerical consequence of the size conditions, the orbit window counts and
$p\ge13$, $r\ge2$, $c_*\ge8$). Then $d-E'+e_{\mathrm{SP}}=d_{TT}$, so $P\cdot\mathrm{SP}$ has multidegree $d_{TT}$, and the
required family is obtained by splitting $P\cdot\mathrm{subst}_g(\mathrm{SP})$ according to the order of its monomials and
realising every piece with the two generator properties and the monomial statement
(\lref{Formalism.MainNotions.Pf14.exists_pol_family_of_mem_gen_of_vsDeg_le}).
\end{proof}

\begin{proposition}[Primitivity of the products]\label{pf:linkG:prim}
For every datum and all constants, $\mathsf{Prim}(C)$ holds for every $C\ge1$. In particular the source primitive
integrality statement holds.
\formallabel{Formalism.MainNotions.TestSideCounting.rest6a}
\formallabel{Formalism.MainNotions.TestSideCounting.source_integrality_of_substitution}
\formallabel{Formalism.MainNotions.TestSideCounting.primStep_of}
\end{proposition}

\begin{proof}
Fix test-side and edge data. Let $\nu_s$ be the partition of the complementary weight at $s$. The integral dual Weyl
modules $\nabla(\nu_s)$ have standard monomial bases $(b_{s,y})$, triangular for the column-lexicographic order, and the
products $s_{\mathbf i}=\bigotimes_sb_{s,i_s}$, $\mathbf i\in I=\prod_sI_s$, lie in the tensor lattice; their number
$|I|$ is the number of weight vectors of the complementary weight
(\lref{Formalism.MainNotions.TestSideCounting.card_all}). The \emph{translation step}
(\lref{Formalism.MainNotions.TestSideCounting.transStep}) transports the $s_{\mathbf i}$ along the edge basis to shift
vectors $\mathrm{sh}_{\mathbf i}\in\linkGR_e(\lambda_{\mathrm{rest}})$ with the following property: if polynomials $P_{\mathbf i}$ of
multidegree $d_{TT}$ satisfy $\mathrm{subst}_g(P_{\mathbf i})=Q_{\mathbf i}$ with $Q_{\mathbf i}$ integral and
$\sum_{\mathbf i}\Theta(P_{\mathbf i})\mathrm{sh}_{\mathbf i}\in\varpi\linkGR_e(\lambda_0)$, then
$\sum_{\mathbf i}\theta(\bar Q_{\mathbf i})\prod_s\bar b_{s,i_s}=0$ over the residue field, where $\theta$ is the theta map of
the flat variables and bars denote reduction modulo $\varpi$. The true symbol polynomial has an integral form
$\mathrm{SP}_\OO$ with $\mathrm{subst}_g(\mathrm{SP})=\mathrm{SP}_\OO$ and $\overline{\mathrm{SP}_\OO}\ne0$
(\lref{Formalism.MainNotions.TestSideCounting.spStep}).

Suppose a primitive combination $\sum_{b,\mathbf i}c_{b,\mathbf i}S_{b,\mathbf i}$ lies in $\varpi\linkGR_e(\lambda_0)$, with
$c_{b_0,\mathbf i_0}$ a unit. Put $P_{\mathbf i}=\sum_bc_{b,\mathbf i}F^{\mathrm s}_b$ and $T_{\mathbf i}=\sum_bc_{b,\mathbf i}B_b$; then
$\sum c\,S=\sum_{\mathbf i}\Theta(P_{\mathbf i})\mathrm{sh}_{\mathbf i}$ and $\mathrm{subst}_g(P_{\mathbf i})=Q_{\mathbf i}:=T_{\mathbf i}
\mathrm{SP}_\OO$, because $\mathrm{subst}_g(F^{\mathrm s}_b)=\mathrm{sym}_g(B_b)=B_b\,\mathrm{SP}_\OO$. The $\bar Q_{\mathbf i}$ are
multihomogeneous of multidegree $d_{TT}$ (\lref{Formalism.MainNotions.TestSideCounting.substPol}), so the independence
statement \lref{Formalism.MainNotions.TestSideCounting.bigInd} (an induction on the leading monomials of the
$\bar b_{s,y}$, using that $\nu_s$ has odd parts only where allowed) gives $\bar Q_{\mathbf i}=0$ for all $\mathbf i$. As
$\overline{\mathrm{SP}_\OO}\ne0$ and the polynomial ring over the residue field is a domain, $\bar T_{\mathbf i_0}=0$, i.e.\
$\varpi$ divides $\sum_bc_{b,\mathbf i_0}B_b$. Since $(B_b)$ is linearly independent with a complement, $\varpi$ divides every
$c_{b,\mathbf i_0}$, contradicting that $c_{b_0,\mathbf i_0}$ is a unit. The constant $1$ comes from the fact that the
statements used hold for all $C$.
\end{proof}

\paragraph{Tensor operators and the cloud bound.}
The \emph{tensor operator statement} $\mathsf{Ten}$ is: for a field $K$ of characteristic $0$, $n'$, $m\ge1$, an orientation
vector defining the theta map $\theta$, a set $TT$, degrees
$d$, $c\in\N$ and partitions $\nu_t$ ($t\in\Z/n'$) with $\nu_t(1)\le c$ for $t\in TT$ and $\nu_t=0$ for $t\notin TT$, and for all
$\beta,s$ in the tensor product $\bigotimes_t\nabla(\nu_t)$, there are words $w_l$ in the operators $E^{(t)}_{ij}$ ($t\in TT$,
at most $c$ letters at each position) and $a_l\in K$ such that for all $G,H$ of multidegree $d$, the polynomial
$\sum_la_lw_l(G)$ has multidegree $d$ and
\[
\bigl\langle\theta(G)s,\theta(H)\beta\bigr\rangle=\Bigl\langle\theta\Bigl(\sum_la_lw_l(G)\Bigr),\theta(H)\Bigr\rangle .
\]
Thus multiplication by $s$ and $\beta$ on the two sides of the Fock pairing is realised by a bounded word in the
polarisation operators.

\begin{proposition}[Tensor operators]\label{pf:linkG:tenop}
$\mathsf{Ten}$ holds.
\formallabel{Formalism.MainNotions.TestSideWords.tensor_operators_holds}
\formallabel{Formalism.MainNotions.TestSideWords.tenOp_of_OB}
\end{proposition}

\begin{proof}
\emph{Operator words.} Let $W$ be a space of endomorphisms of the homogeneous polynomials of degree $d$ in $k$ variables
which is stable under the commutators with the operators $E_{ab}$ and spanned by weight vectors (eigenvectors of
$[E_{aa},\cdot\,]$) whose weights are at most $c$ in every coordinate. Then every element of $W$ is a linear combination of
words of length at most $c$ in the $E_{ab}$ (\lref{operator_words_bound}: the operator words bound for the action
of $\mathfrak{gl}_k$ by commutators). \emph{One position.} For one position ($n'=1$) and $\beta,s\in\nabla(\nu)$, the operator
$A_{s,\beta}$ defined by $\langle\theta(P)s,\theta(Q)\beta\rangle=\langle\theta(A_{s,\beta}P),\theta(Q)\rangle$ spans, as
$s,\beta$ vary, a space stable under commutators with the $E_{ab}$ (the theta map intertwines the polarisation operators,
\lref{Formalism.MainNotions.TestSideWords.letter_operators_holds}), spanned by weight vectors whose weights are bounded by
$\nu(1)\le c$ (the weights of $\nabla(\nu)$ are bounded by its first part,
\lref{Formalism.MainNotions.TestSideWords.dual_weyl_weights_bounded}); so the operator words bound gives the statement
at one position (\lref{Formalism.MainNotions.TestSideWords.onePos_of_OB}). \emph{Tensor products.} The Fock pairing and
the theta map are tensor products over the positions, and the statement for $\bigotimes_t\nabla(\nu_t)$ follows from the
statement at each position by expanding $s$ and $\beta$ in pure tensors and composing the words at the different positions
(\lref{Formalism.MainNotions.TestSideWords.tensor_operators_of_one_position}); positions outside $TT$ carry the trivial
module and contribute no letters.
\end{proof}

\begin{proposition}[Test-side operators]\label{pf:linkG:tsop}
If $\mathsf{Ten}$ holds, then for every datum $D$ and all constants, $\mathsf{TO}(C)$ holds for large constants, with $C_1=1$.
\formallabel{Formalism.MainNotions.TestSideCounting.restTSOp_of_ten}
\formallabel{Formalism.MainNotions.TestSideCounting.restOpF_of_ten}
\formallabel{Formalism.MainNotions.TestSideCounting.restTSOp_of_opF}
\formallabel{Formalism.MainNotions.TestSideCounting.restTSOp_of}
\formallabel{Formalism.MainNotions.TestSideCounting.tsOpAt_of}
\end{proposition}

\begin{proof}
\emph{Step 1.} For $C\ge1$, test-side data and a true theta isomorphism $\Theta$, every $\kappa\in K_f$, all
$\beta,s\in\linkGV(\lambda^{TT}_{\mathrm{rest}})$ and every $F$ of multidegree $d_{TT}$, there is $F'\in\mathrm{Op}_{c_0}(F)$ with
$\langle\Theta(F'),Y\rangle=\kappa\,\langle\Theta(F)s,Y\beta\rangle$ for all $Y\in\linkGV(\lambda_{\mathrm{fl}})$
(\lref{Formalism.MainNotions.TestSideCounting.opFAt_of_ten}). Indeed $\linkGV(\lambda_{\mathrm{rest}}^{TT})=\bigotimes_t\nabla(\nu_t)$
with $\nu_t$ the partition of the complementary weight at $t$, which vanishes off $TT$ and satisfies $\nu_t(1)\le c_0$ on $TT$
(the first entry of the complementary weight is at most the conductor,
\lref{Formalism.MainNotions.TestSideCounting.first_weight_le_conductor}); $\mathsf{Ten}$ with $c=c_0$ gives words and
coefficients, and $F'=\kappa\sum_la_lw_l(F)$. \emph{Step 2.} For every dominant weight $\lambda$ there are a duality $\Psi$ with
loss $\kappa_1(\lambda)$ and $c_\lambda\ne0$ with $\Psi^{-1}(y)(u)=c_\lambda\langle u,y\rangle$ for all $u,y$
(\lref{Formalism.L2.G.Model.test_side_explicit_identification}: the Fock pairing identifies the dual lattices up to the
loss $\kappa_1$). \emph{Step 3.} Take $\Psi_0,\Psi_{\mathrm{fl}}$ from Step 2 for $\lambda_0^{TT}$ and $\lambda_{\mathrm{fl}}$. For
$T\in\linkGV(\lambda_{\mathrm{fl}})^\vee$ put $Y=\Psi_{\mathrm{fl}}(T)$, so that $T(\Theta F')=c_{\mathrm{fl}}\langle\Theta F',Y\rangle$; with
$\kappa=p^{A_r+\kappa_1(\lambda^{TT}_0)}c_0'/c_{\mathrm{fl}}$ ($c'_0$ the constant of $\Psi_0$) Step 1 gives $F'$ with
$p^{A_r+\kappa_1}\Psi_0^{-1}(Y\beta)(\Theta(F)s)=p^{A_r+\kappa_1}c'_0\langle\Theta(F)s,Y\beta\rangle=c_{\mathrm{fl}}\langle\Theta
F',Y\rangle=T(\Theta F')$.
\end{proof}

\begin{lemma}[Size of the losses]\label{pf:linkG:logsize}
Let $p>1$, $a\ge1$, and let $\lambda$ be a weight with $\ell_s(\lambda)=0$ for $s\notin TT$ and $\ell_s(\lambda)+k\le a^2$ for
$s\in TT$. Then $\kappa_1(\lambda)\le2|TT|k^2\log_pa$. Consequently, for $C\ge2$ and test-side data with constant $C$,
$\kappa_1(\lambda_0^{TT})+\kappa_1(\lambda_{\mathrm{fl}})\le4|TT|k^2\log_pa$.
\formallabel{Formalism.MainNotions.TestSideCounting.log_size_bound_le}
\formallabel{Formalism.MainNotions.TestSideCounting.sizeStep}
\formallabel{Formalism.MainNotions.TestSideCounting.restSize}
\end{lemma}

\begin{proof}
Only $s\in TT$ with $\ell_s\ne0$ contribute. For them $\binom k2\le k^2$ and, as $0<\ell_s+k-1\le a^2$,
$\lfloor\log_p(\ell_s+k-1)\rfloor\le\log_p(a^2)=2\log_pa$; the terms are non-negative since $a\ge1$. Summing over $TT$ gives
the first claim. For the second, the size hypothesis of $\linkGES$ (with $C\ge2$, $k\ge1$, $r\ge2$, $|TT|\ge1$, $c_*\ge8$)
gives $a\ge8k$, hence $a\ge8$, $kc_0\le a$ and $4a\le a^2$. The top weight $\lambda_0^{TT}$ has width $\ell_s=a_T(s)+b_T(s)$ at
$s\in TT$, and the bounds for $a_T(s)$ and $b_T(s)\le c_0\le a$ in the hypotheses give $a_T(s)+b_T(s)+k\le a^2$; the flat weight
has width $\ell_s=d_s$, and the bounds for the flat degrees give $d_s+k\le a^2$. Apply the first claim to both weights.
\end{proof}

\begin{proposition}[The cloud bound]\label{pf:linkG:cloudb}
If $\mathsf{TO}(C)$ holds for large constants, then $\mathsf{Cl}(C,C_K)$ holds for large constants.
\formallabel{Formalism.MainNotions.TestSideCounting.restCloudB_of_op}
\formallabel{Formalism.MainNotions.TestSideCounting.restCloudB_of_two}
\formallabel{Formalism.MainNotions.TestSideCounting.restCloudB_of_three}
\formallabel{Formalism.MainNotions.TestSideCounting.restPairTop_of}
\end{proposition}

\begin{proof}
Assume $\sfname{Adm}$, and let $C_2$ be such that $\mathsf{TO}(C)$ holds for $C\ge C_2$, and let $C\ge\max(C_2,2)$. Three facts
are combined.
\begin{enumerate}[label=(\roman*)]
\item The dualities $\Psi_0,\Psi_{\mathrm{fl}}$ and the operator identity of $\mathsf{TO}(C)$.
\item The losses satisfy $\kappa_0+\kappa_{\mathrm{fl}}\le4|TT|k^2\log_pa$, where $\kappa_0=\kappa_1(\lambda_0^{TT})$ and
$\kappa_{\mathrm{fl}}=\kappa_1(\lambda_{\mathrm{fl}})$ (\cref{pf:linkG:logsize}).
\item \emph{Products span}: for $\lambda=\lambda_{\mathrm{fl}}+\lambda_{\mathrm{rest}}$ with $\lambda_{\mathrm{rest}}$ dominant and every Hecke
edge $e$ of type $r$, $\linkGR_e(\lambda)$ is spanned over $\OO$ by the products $y\beta$ with $y\in\linkGR_e(\lambda_{\mathrm{fl}})$
and $\beta\in\linkGR_e(\lambda_{\mathrm{rest}})$ (\lref{Formalism.MainNotions.TestSideCounting.ontoAll}: the
multiplication of tensor lattices of dual Weyl modules is surjective).
\end{enumerate}
By \lref{Formalism.MainNotions.TestSideCounting.pairAt_of_tsop}, (i)--(iii) give the \emph{pair statement} for the
top part $\lambda_0^{TT}$: dualities with losses $\kappa_0,\kappa_{\mathrm{fl}}$, $\kappa_0+\kappa_{\mathrm{fl}}\le4|TT|k^2\log_pa$, mapping
the dual lattices of the target and the lattices of the edges with the target of $e$ as in the definition of a duality,
the spanning of $\linkGR_{e^\vee}$ by products, and the operator identity. From the pair statement the cloud bound is
derived (\lref{Formalism.MainNotions.TestSideCounting.restCloudB_of_pair}: the weight $\lambda_0$ is the sum of its top
part and its part outside $TT$, and the map $\Lambda$ and the element $F_1$ of the cloud bound are constructed from the
dualities, the spanning by products and the operator identity; the exponent is bounded by the bound $C_K|TT|k^2\log_pa$
for $\kappa_0+\kappa_{\mathrm{fl}}$, and the loss $\ell$ accounts for the operators, at most $r|TT|c_0$, and for the part of
$\lambda_0$ outside $TT$, at most $2rn'c_{pp}$).
\end{proof}

\paragraph{Flatness.}

\begin{proposition}[From $\Xi$ to flatness]\label{pf:linkG:flat}
Suppose that $\sfname{Adm}$ and the plane Gram statement with constants $(C_H,C_E,C_6,C_{18})$ imply that there are
$C_t,C_D,C_2>0$ with $\Xi(C,C_t,C_D)$ for $C\ge C_2$. Then $\sfname{Adm}$, the source cleaning statement and the plane Gram
statement imply that there are $C_t,C_D,C_1>0$ with $\mathsf{Flat}(C,C_t,C_D)$ for $C\ge C_1$.
\formallabel{Formalism.MainNotions.TestSideCounting.restFlat_of_Xi}
\formallabel{Formalism.MainNotions.TestSideCounting.restFlat_of}
\end{proposition}

\begin{proof}
Keep the constants $C_t,C_D$ of $\Xi$. Three further inputs are used: the \emph{realisation of test polynomials}
(\lref{Formalism.MainNotions.TestSideCounting.restReal}: for $C$ beyond a constant there is a Witt matrix $g_v$ such that
every integral polynomial $\sigma$ of multidegree $d_{TT}$ which is realisable of order $o$ for $g_v$ is represented by a
functional $T_\sigma\in\varpi^{-o}\linkGH_{e_t}(\lambda_{\mathrm{fl}})^\vee\cap\linkGR'_e(\lambda_{\mathrm{fl}})$, lying in
$\varpi^{c'}\linkGR'_{e'}(\lambda_{\mathrm{fl}})$ for the other edges $e'$ with the target of $e$, in the sense that
$b(\Theta(F),T_\sigma)$ differs from $p^{\kappa_{\mathrm{fl}}}\langle\mathrm{subst}_g(F),\sigma\rangle/\prod_sd_s!$ by an element of
$\varpi^{c'}\OO$ for every $F$ with $\Theta(F)\in\linkGR_e(\lambda_{\mathrm{fl}})$);
the one-sided lemma for the admitted data, by which one-sided elements are realisable
(\lref{Formalism.MainNotions.TestSideCounting.one_sided_lemma4_of_admitted_theorem_II}); and the hypotheses of the edge
step at the given edge, which make $c'$ larger than the valuation bounds in question
(\lref{Formalism.MainNotions.Pf14.hypII_at}). Let $C$ be at least the constants of $\Xi$, of the realisation, of the
edge hypotheses and $C_D$. Given $F_1$ with
$\Theta(F_1)\in\linkGR_e(\lambda_{\mathrm{fl}})\smallsetminus\varpi\linkGR_e(\lambda_{\mathrm{fl}})$, $\Xi$ gives one-sided elements
$\sigma_i$ of orders $o_i\le t_F$ with $v_p(p^{\kappa_{\mathrm{fl}}}\langle\mathrm{subst}_g(F_1),\sum_i\sigma_i\rangle/\prod_sd_s!)\le
D_X$. Put $T=\sum_iT_{\sigma_i}$. Then $T\in\varpi^{-t_F}\linkGH_{e_t}(\lambda_{\mathrm{fl}})^\vee\cap\linkGR'_e(\lambda_{\mathrm{fl}})$ (as
$o_i\le t_F$), $T$ lies in $\varpi^{c'}\linkGR'_{e'}$ for the other edges, and, the Fock pairing being additive in its right
argument, $b(\Theta(F_1),T)=u+\varpi^{c'}z$ with $z\in\OO$ and $u=p^{\kappa_{\mathrm{fl}}}\langle\mathrm{subst}_g(F_1),\sum_i\sigma_i\rangle/
\prod_sd_s!$; by the valuation bound, $u\in\OO$ and $uu'=p^{D_X}$ for some $u'\in\OO$. Since $\varpi^2=p$ and $2D_X<c'$ (a
consequence of the hypotheses of the edge step), $\varpi^{c'}=p^{D_X}\varpi^{c'-2D_X}=uu'\varpi^{c'-2D_X}$, so
$b(\Theta(F_1),T)=u\,(1+u'\varpi^{c'-2D_X}z)$, and the second factor is a unit of the local ring $\OO$. Hence
$v_p(b(\Theta(F_1),T))=v_p(u)\le D_X\le\delta_D$.
\end{proof}

\begin{proposition}[From the numerical statement to $\Xi$]\label{pf:linkG:xi}
Suppose that $\sfname{Adm}$ implies that there are $C_D,C_2>0$ with $\mathsf{NS}(C,C_D)$ (constants $(C_H,C_E,C_6,C_{18})$) for
$C\ge C_2$. Then $\sfname{Adm}$ and the plane Gram statement with constants $(C_H,C_E,C_6,C_{18})$ imply that there are
$C_t,C_D,C_1>0$ with $\Xi(C,C_t,C_D)$ for $C\ge C_1$.
\formallabel{Formalism.MainNotions.TestSideCounting.restXi_of_NS}
\formallabel{Formalism.MainNotions.TestSideCounting.restXi_of_P}
\formallabel{Formalism.MainNotions.TestSideCounting.restXi_of_N}
\end{proposition}

\begin{proof}
\emph{Step 1} (\lref{Formalism.MainNotions.TestSideCounting.restXiP_of_NS}). The plane Gram statement applies to the
discrete valuation ring of the setting, whose fraction field has characteristic $0$; combined with $\mathsf{NS}(C,C_D)$ it gives,
for every polynomial $G$ in the operator closure of the span of the $\mathrm{sym}_P(B_b)$ with integral coefficients not all in
$\varpi\OO$, one-sided elements $\sigma_i$ with the valuation bound of $\Xi$ (the perturbation hypotheses of $\mathsf{NS}(G)$ are
supplied by the plane Gram statement). \emph{Step 2} (\lref{Formalism.MainNotions.TestSideCounting.restXiN_of_P}). This is transferred to the polynomials $F_1$
of the operator closure of the realisations with $\Theta(F_1)\in\linkGR_e(\lambda_{\mathrm{fl}})\smallsetminus\varpi\linkGR_e(\lambda_{\mathrm{fl}})$:
$G=\mathrm{subst}_g(F_1)$ lies in the operator closure of the span of the symbol polynomials (the substitution is compatible
with the polarisation operators and maps the realisations to the symbols), and by the true theta property $G$ has integral
coefficients, not all in $\varpi\OO$. \emph{Step 3} (\lref{Formalism.MainNotions.TestSideCounting.restND_of_adm}). The symbol space
statement (clause 27, with constant $C_g$) bounds the orders: for $C\ge\max(C_g,|C_\Lambda|)$ and with
$C_t=\max(1,|C_{\deg}|)$, the realisations have degrees at most $t_0$, so the orders $o_i$ are at most $t_0+|TT|c_0$, the
words contributing at most $c_0$ per position of $TT$. \emph{Step 4}
(\lref{Formalism.MainNotions.TestSideCounting.restXi_of}). Steps 1--3 give $\Xi(C,C_t,C_D)$ for $C\ge C_1$, the maximum
of the constants involved.
\end{proof}

\paragraph{Proposition 5 and the words.}

\begin{proposition}[Proposition 5]\label{pf:linkG:prop5}
For every datum and all constants, $\mathsf{P5}(C,C_D)$ with the constants $(C_H,C_E,C_6,C_{18})=(26,2,24,60)$ holds for large
constants.
\formallabel{Formalism.MainNotions.Proposition5Numerics.proposition5_holds}
\formallabel{Formalism.MainNotions.Proposition5Numerics.proposition5_of_three_parts}
\end{proposition}

\begin{proof}
The statement is assembled from three parts, each of which holds under $\sfname{Adm}$ for $C$ beyond a constant; fix data of
the type of $\mathsf{P5}$ and write $L=|TT|k^2\log_pa$.
\begin{enumerate}[label=(\roman*)]
\item \emph{Choice} (\cref{pf:linkG:pick}, with the constant $C_{Dn}=9$): for every $G$ as
in $\mathsf{P5}$ there are signed chain data satisfying the degree condition of the frame, block data, a dual element $F_2$
in the $X$-space of the chain data with vanishing $B$-valuation, a norm test $N_0$ and $k_0\in K_f$ with
$v_p(k_0)\le\nu_2(d)-n_0+C_{Dn}L$ ($\nu_2(d)$ a valuation of a product of factorials of $d$, $n_0$ a parameter of the chain
data), such that for every integral $A$ of the plane multidegree,
$v_p\bigl(\langle G,\mathrm{asm}(A)\rangle-k_0\,B(F_2,A)\bigr)\ge c'+\nu_2(d)-\kappa_{\mathrm{fl}}$, where $\mathrm{asm}(A)$ is the
test polynomial assembled from $A$, the block data and $N_0$.
\item \emph{Membership} (\cref{pf:linkG:memb}): for such chain and block data and
every integral $G_p$ in the $Y$-space with the constants $(24,60)$, the assembled polynomial $\mathrm{asm}(G_p)$ is a sum of
one-sided elements of multidegree $d_{TT}$.
\item \emph{Numerics} (\lref{Formalism.MainNotions.Proposition5Numerics.restNum}, with a constant $C_{Dm}$): $L\ge1$,
$\kappa_{\mathrm{fl}}\le C_{Dm}L$, and the chain data of (i) satisfy the hypothesis with constant $26$, $4(D_m+D_l+3k_p)<c$,
$k_p\ge0$ and $D_m+D_l+3k_p\le C_{Dm}L$.
\end{enumerate}
Given $G$, take the chain data and $F_2$ of (i); by (iii) they satisfy the requirements of $\mathsf{NS}(G)$ with
$C_H=26$. For an admissible perturbation $G_p+R$ (with $(C_6,C_{18})=(24,60)$ and $C_E=2$), (ii) writes the assembled test
polynomial as a sum of one-sided elements $\sigma_i$, and by (i) the valuation of
$p^{\kappa_{\mathrm{fl}}}\langle G,\sum_i\sigma_i\rangle/\prod_sd_s!$ is controlled by $v_p(k_0)$ and the valuation of the
$B$-pairing of $F_2$ with the perturbation, which is bounded in the hypotheses of $\mathsf{NS}(G)$; the sum of the losses is at
most $(C_{Dn}+2C_{Dm}+1)L$ by (iii) (\lref{Formalism.MainNotions.Proposition5Numerics.value_core}). Hence $\mathsf{P5}(C,C_D)$
holds with $C_D=C_{Dn}+2C_{Dm}+1$ for $C\ge\max(C_1,C_2,C_3,C_{Dn}+3C_{Dm})$, where $C_1,C_2,C_3$ are the constants of
(i)--(iii).
\end{proof}

\begin{proposition}[The word statement]\label{pf:linkG:word}
Assume the ten one-operator steps of the letter substitution (the effect of a single operator $E^{(t)}_{ij}$ on the plane
and non-plane variables, the derivative and multiplication rules, the Euler operator and two geometric conditions on the
tops). Then $\mathsf{Word}(C)$ holds for large constants.
\formallabel{Formalism.MainNotions.PairMembershipBlocks.restWord_of_steps}
\formallabel{Formalism.MainNotions.TestSideWords.restWord_of}
\formallabel{Formalism.MainNotions.PairMembershipBlocks.restOneOp_of}
\end{proposition}

\begin{proof}
Write $\mathsf{Mem}_\le(s_c,s_L;\epsilon)$ for the refinement of $\mathsf{Mem}(s_c,s_L)$ in which the growth bound is weakened by a
loss $\epsilon(j,u)\ge0$ depending on the level $j$ and the position $u$. Four statements are combined
(\lref{Formalism.MainNotions.TestSideWords.restWordStep_of}).
\begin{enumerate}[label=(\roman*)]
\item \emph{Refinement} (\lref{Formalism.MainNotions.PairMembershipBlocks.restLeOfG}): $\mathsf{Mem}(s_c,s_L)$ implies
$\mathsf{Mem}_\le(s_c,s_L;0)$.
\item \emph{One operator} (\lref{Formalism.MainNotions.PairMembershipBlocks.oneOpStep_of}, for all $C\ge1$): if $F$ satisfies
$\mathsf{Mem}_\le(0,\frac18;\epsilon)$, then for $s\in TT$ and all $i,j$ the polynomial $E^{(s)}_{ij}F$ satisfies
$\mathsf{Mem}_\le(0,\frac18;\epsilon+\mathrm{cost}_s)$, where $\mathrm{cost}_s(j,u)$ is the cost of one operator at the position $s$;
the proof follows the level images of $E^{(s)}_{ij}F$ through the ten one-operator steps.
\item \emph{Word bound} (\lref{Formalism.MainNotions.PairMembershipBlocks.wordBoundStep}, for $C\ge16$): for a word $w$ with
letters in $TT$ and at most $c_0$ letters at each position, the total cost $\mathrm{cost}_w(j,u)=\sum_{\text{letters}}
\mathrm{cost}_s(j,u)$ is at most $\frac18$ of the slope of the frame at the position $u$ of the level $j$.
\item \emph{Conversion}: a loss bounded by $\frac18$ of the slope turns $\mathsf{Mem}_\le(0,\frac18;\mathrm{cost}_w)$ into
$\mathsf{Mem}(\frac18,\frac38)$.
\end{enumerate}
Given $F$ with $\mathsf{Mem}(0,\frac18)$, (i) and induction on the length of $w$ with (ii) give
$\mathsf{Mem}_\le(0,\frac18;\mathrm{cost}_w)$ for $w(F)$; (iii) and (iv) give $\mathsf{Mem}(\frac18,\frac38)$. This holds for
$C\ge16$.
\end{proof}

\begin{proposition}[The level link]\label{pf:linkG:link}
The \emph{level link statement} holds: let $F$ be a valid frame with $n'$ even, $J\ge2$ and signed chain data tied to $F$,
let $A_i$ be polynomials in the letter variables of $F$ ($A_i$ involving only the letters of the segment $i$, $i<J$), $l$ a
label and $\gamma$ a plane-even exponent, and let $t(j,u)$ be lower bounds for the valuations of $A_{\mathrm{seg}(j)}$ on the
variables below the positions $u<n_L(j)+r$. Then for every $j$ and every $\mu<r-1$, the valuation on the true letters of the
$T$-set $T(j,\mu)$ of the compressed level image of the polynomial $z(A,l,\gamma)$ is at least
$t(j,n_L(j)+r-1-\mu)-\mathrm{dw}(j)$, where $\mathrm{dw}(j)$ is the degree weight of $l+\gamma$ at $j$.
\formallabel{Formalism.MainNotions.PairMembershipBlocks.restLink}
\formallabel{Formalism.MainNotions.PairMembershipBlocks.level_link_of_degree_letters}
\end{proposition}

\begin{proof}
The statement is reduced to three lemmas on the level map $\Phi_j$, which sends polynomials in the letter variables of the
segment of $j$ to the variables of the levels
(\lref{Formalism.MainNotions.PairMembershipBlocks.restLink_of}, which also uses the homogeneous decomposition of
polynomials and the additivity of the valuations).
\emph{Letters} (\lref{Formalism.MainNotions.PairMembershipBlocks.linkL}): if $P$ involves only the letters of the
segment of $j$, its monomials have weighted degree at most $\mathrm{dw}$, and its valuation on the variables below the
position $n_L(j)+r-1-\mu$ is at least $\theta$, then the valuation of $\Phi_j(P)$ on the true letters of $T(j,\mu)$ is at least
$\theta-\mathrm{dw}$. \emph{Degrees} (\lref{Formalism.MainNotions.PairMembershipBlocks.lkDeg}): if $P$ is homogeneous of
multidegree $\delta$, every monomial of $\Phi_j(P)$ has degree $P_{j-1}(m')\,\delta(f(m'))$ in the plane variable $(j-1,m')$, where
$P_{j-1}$ and $f$ are data of the chain and of the frame. \emph{Variables}
(\lref{Formalism.MainNotions.PairMembershipBlocks.lkVars}): $\Phi_j(P)$ involves only the letter variables of the level $j$
and plane variables. The compressed level image of $z(A,l,\gamma)$ is obtained by applying $\Phi_j$ to $A_{\mathrm{seg}(j)}$ and
multiplying by the images of the label monomials; the lemmas on degrees and variables show that this multiplication shifts
the valuation on the true letters by at most the degree weight $\mathrm{dw}(j)$, and the lemma on letters gives the bound
$t(j,n_L(j)+r-1-\mu)-\mathrm{dw}(j)$.
\end{proof}

% ---------------------------------------------------------------------------
\subsection{Verification of the clauses; proof of the local lattice theorem}\label{sec:pf:linkG:verif}

Most clauses are restatements of components of a package of \emph{tool statements} about $\GL_n$-representations
over discrete valuation rings, proved in the development under $(\mathrm E)$
(\lref{Formalism.MainNotions.G.tool_statements_hold}): it consists of statements on invariant projections and
their denominators, integral invariants, isotypic projectors, reducibility, three degree bounds, a support lemma,
multiplicity and second decompositions, and a lattice index statement, and of the hyperplane rank statement
(clause 15). One clause is a count of subspaces over a finite field, for which the following classical bound is
used.

\begin{lemma}[Gaussian binomials]\label{pf:linkG:gauss}
For integers $q\ge0$ and $0\le r\le k$, the Gaussian binomial coefficient satisfies $\binom{k}{r}_q\ge q^{r(k-r)}$.
Consequently, for a finite field $\F$, the number of $r$-dimensional subspaces of $\F^k$ is at least
$|\F|^{r(k-r)}$.
\formallabel{Formalism.MainNotions.G.gauss_binomial_ge}
\formallabel{Formalism.MainNotions.G.rational_planes_count}
\end{lemma}

\begin{proof}
Induction on $k$, using $\binom{k}{0}_q=1$ and the $q$-Pascal rule
$\binom{k+1}{r+1}_q=\binom{k}{r}_q+q^{r+1}\binom{k}{r+1}_q$. For $r+1=k+1$ the first term is $\binom kk_q\ge1=q^0$. For
$r+1\le k$, write $k=r+1+d$; by induction $\binom{k}{r+1}_q\ge q^{(r+1)d}$, so
$\binom{k+1}{r+1}_q\ge q^{r+1}q^{(r+1)d}=q^{(r+1)(d+1)}$. The second statement follows since the number of
$r$-dimensional subspaces of $\F^k$ is $\binom kr_{|\F|}$ (\lref{Formalism.G0.Top.subspaceCount}).
\end{proof}

\begin{lemma}[Products of sections]\label{pf:linkG:products}
Let $\OO\subseteq K_f$ be a ring and its fraction field, $\Lambda\subseteq K_f^{m+1}$ an $\OO$-submodule and $\lambda,\mu\in\Z^{m+1}$
non-increasing. Then $FG\in H^0(\lambda+\mu,\Lambda)$ for $F\in H^0(\lambda,\Lambda)$ and $G\in H^0(\mu,\Lambda)$.
Consequently, for dominant weights $\lambda_0,\mu_0$ and modules $\Lambda_s$, the product of an element of
$\bigotimes_sH^0(\lambda_s,\Lambda_s)$ and an element of $\bigotimes_sH^0(\mu_s,\Lambda_s)$ lies in
$\bigotimes_sH^0(\lambda_s+\mu_s,\Lambda_s)$; this is clause 37.
\formallabel{Formalism.L2.G.Model.mul_mem_riemannRochSpace_add}
\formallabel{Formalism.L2.G.Model.tensorLat_mul}
\formallabel{Formalism.L2.G.Model.products_sections_integral}
\end{lemma}

\begin{proof}
Let $g$ be a matrix whose columns span $\Lambda$, and write $F=\det(g)^{-\lambda_m}F_0$, $G=\det(g)^{-\mu_m}G_0$ with
$\rho(g)F_0\in\nabla_\OO(\lambda-\lambda_m)$ and $\rho(g)G_0\in\nabla_\OO(\mu-\mu_m)$. The substitution $\rho(g)$ and the
column operations of \cref{pf:linkG:weightlat}(a) are ring homomorphisms, and column weights add under multiplication;
hence the product of an element of $\nabla_\OO(\nu)$ and an element of $\nabla_\OO(\nu')$ lies in $\nabla_\OO(\nu+\nu')$,
and $\rho(g)(F_0G_0)\in\nabla_\OO(\lambda+\mu-(\lambda+\mu)_m)$. Thus
$FG=\det(g)^{-(\lambda+\mu)_m}F_0G_0\in\nabla^{\mathrm{tw}}_\OO(\lambda+\mu,g)$, and intersecting over $g$ gives the first
claim. For the second, the tensor lattices are spanned by products over $s$ of elements of the factors, the product of two
such generators is the product over $s$ of the products of the factors (the variable sets for different $s$ are disjoint),
and the weight attached to $\lambda_0+\mu_0$ at $s$ is $\lambda_s+\mu_s$ (also for odd $s$, where the weight is reversed and
negated); bilinearity extends the claim to the spans.
\end{proof}

\begin{proposition}[Plane count theorem with constants]\label{pf:linkG:planecount}
Assume $(\mathrm E)$ and $(\mathrm L6)$. For every datum $D$ there are $C_\Lambda,C_k,C_{rk},C_{\deg}$ such that clause
27 holds with these constants.
\formallabel{Formalism.MainNotions.G.plane_count_theorem_with_constants}
\end{proposition}

\begin{proof}
The ``second part'' of $(\mathrm L6)$ says: if a list of thirteen statements holds, then there is
$C_0$ such that for all $C,C'\ge C_0$ there are $C_{rk},C_{\deg}$ with clause 27 for $(C,C',C_{rk},C_{\deg})$. The list
consists of a symbol space orbit statement
(\lref{Formalism.MainNotions.G.TopChoiceBounds.symbol_space_orbit_holds}), a valuation transport
(\lref{Formalism.MainNotions.G.Joints.valuation_transport_of_lemma}), the three Gelfand--Tsetlin counting forms of
$(\mathrm L6)$, the subspace count of \cref{pf:linkG:gauss}, five statements obtained from the tool statements (degree
bounds, reducibility, isotypic projectors), the space statement
(\cref{pf:linkG:spacestatement}) and a transport statement
(\lref{Formalism.MainNotions.G.Joints.transport_statement_holds}). Take $C_\Lambda=C_k=C_0$.
\end{proof}

\begin{proposition}[The local statement for every datum]\label{pf:linkG:closedeverywhere}
Assume $(\mathrm E)$ and $(\mathrm L6)$. Then $\mathsf{Concl}(D)$ holds for every datum $D$.
\formallabel{Formalism.MainNotions.G.second_theorem_closed_everywhere}
\end{proposition}

\begin{proof}
Let $D$ be a datum and let $(C_\Lambda,C_k,C_{rk},C_{\deg})$ be as in \cref{pf:linkG:planecount}. We verify
$\sfname{Adm}(D,C_\Lambda,C_k,C_{rk},C_{\deg})$, clause by clause (numbering of \cref{tab:pf:linkG:clauses}):
clauses 1 and 2 are proved in
\lref{Formalism.MainNotions.G.TopChoiceBounds.orbit_positions_admitted_holds} and \lref{Formalism.MainNotions.G.TopChoiceBounds.orbit_counts_admitted_holds};
clauses 3, 4, 5, 12, 23, 28--32, 36 and 37 are components of $(\mathrm L6)$; clause 6 is $(\mathrm E)$; clauses 7--11, 13--22,
24, 25 and 26 are components of the tool statements (each transported by a lemma of three lines, except clause 26,
which is the valuation transport cited above); clause 27 is \cref{pf:linkG:planecount}; clause 33 is the space statement;
clauses 34 and 35 are the symbol space orbit and transport statements. The source cleaning statement holds with
positive constants by \lref{Formalism.MainNotions.G.derived_statements_with_constants} (from the tool statements,
$(\mathrm E)$ and clause 27, through \lref{Formalism.MainNotions.G.symbolGeneratorData.stmtE_proved}), and the
plane Gram statement holds for all $C_H$ beyond a constant by \lref{Formalism.MainNotions.PlaneGram.stmtE}, whose
hypotheses are again supplied by the tool statements
(\lref{Formalism.MainNotions.G.split_theorem_with_constants}). The criterion, a component of $(\mathrm L6)$, gives
$\mathsf{Concl}(D)$.
\end{proof}

\begin{proof}[Proof of \cref{pf:linkG:localthm}]
The formal statement is: there is $C_0\ge0$ such that $\linkGLL(C_0)$ holds \emph{for every datum $D$} (recall that
$\linkGLL$ is formally indexed by a datum, \cref{pf:linkG:LocLat}). If there is a datum $D_0$, then
\cref{pf:linkG:closedeverywhere} gives $C_0>0$ with $\linkGLL(C)$ for all $C\ge C_0$, in particular $\linkGLL(C_0)$; as
the statement does not depend on the datum, it holds for every datum. If there is no datum, $C_0=0$ satisfies the
statement vacuously.
\end{proof}

\begin{proposition}[The tool packages hold]\label{pf:linkG:locmodel}
$(\mathrm E)$ holds, and $\sfname{Loc}$ implies $(\mathrm L6)$. Consequently, by \cref{pf:linkG:loc}, both $(\mathrm E)$
and $(\mathrm L6)$ hold, and the conclusion of \cref{pf:linkG:localthm} holds unconditionally.
\formallabel{Formalism.L2.G.extremal_projector_tools_of_loc}
\formallabel{Formalism.L2.G.local_model_results}
\end{proposition}

\begin{proof}
$(\mathrm E)$ is \lref{Formalism.MainNotions.G.ExtremalProjector.extremal_projector_statement}. The components of
$(\mathrm L6)$ are: the criterion, which is $\sfname{Loc}$; the second part,
\cref{pf:linkG:secondpart} (restriction estimates, good path weights, a
generator modulo the uniformiser and the non-vanishing of a discriminant); the Weyl dimension formula, branching rule
and Pieri rule in counting form
(\lref{Formalism.MainNotions.G.Joints.Lit.weyl_dimension_formula_index_sets}, \lref{Formalism.MainNotions.G.Joints.Lit.branching_rule_count_form}, and \lref{Formalism.MainNotions.G.Joints.Lit.pieri_rule_count_form};
classical facts on Gelfand--Tsetlin patterns, proved in the development from its own counting
\cite{GelfandTsetlin1950,FultonHarris1991}); the seven lattice statements
(\lref{Formalism.L2.G.Model.dualInclusionLattices}, \lref{Formalism.L2.G.Model.bePerfect}, \lref{Formalism.L2.G.Model.lPOneEvery}, \lref{Formalism.L2.G.Model.flatTop}, \lref{Formalism.L2.G.Model.dualHeightBasis}, \lref{Formalism.L2.G.Model.hlSpans}, and \lref{Formalism.L2.G.Model.testSide});
the Witt identification (\lref{Formalism.MainNotions.G.WittPI.wittIdentity_holds}); and clauses 12, 23, 36 and 37
(\lref{Formalism.MainNotions.G.Joints.joint_types_corollary_sentences}, \lref{Formalism.MainNotions.G.Joints.coarse_units_lemma}, and \lref{Formalism.L2.G.Model.fewBadOrbits},
and \cref{pf:linkG:products}).
\end{proof}

% ---------------------------------------------------------------------------
\subsection{The second part of the tool package}\label{sec:pf:linkG:second}

We prove the two components of $(\mathrm L6)$ which concern symbol spaces: the second part (the symbol space statement of
clause 27 from a list of thirteen statements) and the space statement (clause 33).

\begin{theorem}[The second part]\label{pf:linkG:secondpart}
Let $D$ be a datum and assume the thirteen statements of the second-part list (those listed in the proof of
\cref{pf:linkG:planecount}; among them the subspace count $|\mathrm{Gr}_r(\F_q^{r+K})|\ge q^{rK}$ of \cref{pf:linkG:gauss} and the
space statement of \cref{pf:linkG:spacestatement}). Then there is $C_0>0$ such that for all $C_\Lambda,C_k\ge C_0$ there are
$C_{rk},C_{\deg}>0$ for which clause 27 holds with the constants $(C_\Lambda,C_k,C_{rk},C_{\deg})$; one can take $C_0=1088$,
$C_{rk}=1$ and $C_{\deg}=10$.
\formallabel{Formalism.MainNotions.SelectionLemmas.Pf.secondPartClosed}
\formallabel{Formalism.MainNotions.SelectionLemmas.Pf.second_estimate_of_restrictions}
\formallabel{Formalism.MainNotions.SelectionLemmas.Pf.second_estimate_of_rank_bound}
\formallabel{Formalism.MainNotions.SelectionLemmas.Pf.restRank_of}
\formallabel{Formalism.MainNotions.SelectionLemmas.Pf.restCount_of}
\formallabel{Formalism.MainNotions.SelectionLemmas.Pf.restSel}
\end{theorem}

Recall the content of clause 27. The data are: an unramified quadratic setting with $p\ge13$, integers $r\ge2$, $K\ge1$,
$k=r+K$, a valid frame $F$ for a set $TT$, degrees $\deg_t$ with $a\le\deg_t\le2a$ on $TT$, integers $N_s,E_0,m_*$, degree
bounds $E'_t\le E_0$ and reals $a,Y\ge1,Q$ satisfying the size conditions of the clause (every segment has at most $N_s$
letters, $N_s\le N_1$, $N_s\le2|TT|$, $N_s+r\le K$, $p^{\mathrm{dist}}\le Q$, $C_kN_s^2\log a\le k$, $128k^2N_sY^2\le
p^{L_1-1}$, $32pYk^2N_sQ^2(E_0+m_*+1)\le a$). With $\Lambda_a=C_\Lambda N_s^2\log a$ they determine the segment modules of the
frame (with explicit multiplicity, valuation and degree parameters), the set $\mathrm{Gen}$ of \emph{generators} (products over
the segments of elements of the segment modules, multiplied by monomials of the labels, of multidegree
$d=\deg-E'$ on $TT$), and the module $NN$ they span. The conclusion is: there is a free submodule $S\subseteq NN$, with a
complement and a basis consisting of generators, of rank at least $p^{-n'rK}\beta$, where
\[
\beta=\Bigl(1-\frac{C_{rk}|TT|r^2k}{Y}\Bigr)\prod_{t\in TT}\binom{\deg_t+k-1}{k-1}
\]
(and at least $\beta/|\mathrm{Gr}_r(\F_q^k)|$ for every finite field with $q=p^{n'}$ elements); and every monomial of every
element of $NN$ has degree at most $C_{\deg}|TT|\bigl(\Lambda_aa/(pk)+ap^{-L_1}+kQ\bigr)$.

\begin{proof}
Fix the data and write $\bar P$ for the reduction of a polynomial $P$ modulo $\varpi$, $\F$ for the residue field, and
$\overline{NN}=\operatorname{span}_\F\{\bar g:g\in\mathrm{Gen}\}$.

\emph{Step 1: degrees} (\lref{Formalism.MainNotions.SelectionLemmas.Pf.restDeg}). Every monomial of every element of $NN$
has degree at most $10|TT|(\Lambda_aa/(pk)+ap^{-L_1}+kQ)$, for $C_\Lambda\ge0$ (a bound for the degrees of the segment factors
and of the label monomials of the generators).

\emph{Step 2: selection.} Let $s$ be a finite set of exponents and $G$ a set of polynomials over $\OO$ supported in $s$. Then
there is a submodule $S\subseteq\operatorname{span}G$ with a complement, with a basis consisting of elements of $G$, of rank
$\dim_\F\operatorname{span}\{\bar g:g\in G\}$. Indeed, choose
$g_1,\dots,g_n\in G$ whose reductions form a basis of $\operatorname{span}\bar G$. The matrix of the $g_i$ in the monomial basis
of the free module of polynomials supported in $s$ has a minor of size $n$ which is a unit, as its reduction has rank $n$; so
the $g_i$ are linearly independent, and together with the monomials outside the columns of that minor and the
polynomials supported outside $s$ they span the polynomial ring, which gives the complement.

\emph{Step 3: the rank modulo $\varpi$.} We show $\dim_\F\overline{NN}\ge p^{-n'rK}\beta$ when $C_\Lambda,C_k\ge1088$ and $C_{rk}=1$.
Let $\Delta$ be the \emph{discriminant}, the product over the segments $j$ of the images of the products of the $r\times r$
determinants $\det(X_{\mathrm{row}(j,x,a),b})$ of plane variables over the late letters $x$ of $j$, raised to explicit powers.
Let the \emph{$\lambda$-generators} be the products over the segments of elements $y_j$ of the segment modules $M_j(b_j)$ with
the reduced degree parameters, multiplied by monomials of the coarse labels, of multidegree $a'=d-\deg\Delta$, and let
$\bar\Lambda$ be the span of their reductions. Three facts are used.
\begin{enumerate}[label=(\alph*)]
\item $\bar\Delta\ne0$ (\lref{Formalism.MainNotions.SelectionLemmas.Pf.discriminant_nonzero_mod_uniformiser}: the
determinants have pairwise distinct rows, as late vectors have distinct tops).
\item For every $\lambda$-generator $u$ there is $g\in\mathrm{Gen}$ with $\overline{\Delta u}=\bar g$
(\lref{Formalism.MainNotions.SelectionLemmas.Pf.generator_modulo_uniformiser_of_adm}, for $C_\Lambda,C_k\ge1088$: column
conditions from the numerical bounds, a factorisation of coarse labels and a comparison of top degrees).
\item $\dim_\F\bar\Lambda\ge p^{-n'rK}\beta$ (Step 4).
\end{enumerate}
Multiplication by $\bar\Delta$ is injective on the polynomial ring over $\F$, a domain, by (a); by (b) it maps $\bar\Lambda$ into
$\overline{NN}$. As $\overline{NN}$ is finite-dimensional (its elements are supported on the finite set of exponents of
multidegree $d$), $\dim\overline{NN}\ge\dim\bar\Delta\bar\Lambda=\dim\bar\Lambda\ge p^{-n'rK}\beta$.

\emph{Step 4: the count.} For a label datum $\lambda$ (in an explicit finite set $\mathrm{LA}$ of residue data) and a residue
choice $e\in\mathrm{EE}(\lambda)$ let $b(\lambda,e)_j$ be the induced degree parameters on the segment $j$. Then
\[
p^{-n'rK}\beta\ \overset{(4)}{\le}\ \sum_{\lambda\in\mathrm{LA}}\sum_{e\in\mathrm{EE}(\lambda)}\prod_jG_j(b(\lambda,e)_j)
\ \overset{(3)}{\le}\ \sum_{\lambda,e}\prod_j\rank M_j(b(\lambda,e)_j)\ \overset{(1)}{\le}\ \dim_\F\bar\Lambda,
\]
where the functions $G_j\ge0$ and the three inequalities are as follows.
\begin{enumerate}[label=(\arabic*)]
\item The $\lambda$-generators attached to different $(\lambda,e)$ and to bases of the $M_j(b(\lambda,e)_j)$ have linearly
independent reductions (\lref{Formalism.MainNotions.SelectionLemmas.Pf.restLam_of}).
\item For every segment $j$ there is a function $G_j$ of the degree parameters, depending only on the parameters of the late
letters, such that $G_j(b)\le\rank M_j(b)$ for admissible $b$, and such that the \emph{segment bound} holds: for admissible $b$
and every residue vector $\rho$,
\[
\Bigl(1-a^{-N_s^2}\Bigr)\prod_{x}\binom{b_x+k-1}{k-1}\prod_{i<r}\binom{B_w(j,i)+k-1}{k-1}\Bigl(Q_j^{r-1}\max\bigl(0,1-\tfrac r{2Y}\bigr)\Bigr)^r
\le\sum_{e\equiv\rho}G_j(\tau_e b),
\]
the product over the letters $x$ of the segment, $B_w(j,i)$ explicit bounds, $Q_j$ the residue modulus and $\tau_eb$ the
parameters $b$ with the values $e$ at the top letters
(\lref{Formalism.MainNotions.SelectionLemmas.Pf.restSeg_of}; this is the core of the second part: $G_j$ is defined
through the paths of the segment, and its properties are proved from a membership statement for the segment modules
(\lref{Formalism.MainNotions.SelectionLemmas.Pf.restMem_of}), an independence statement
(\lref{Formalism.MainNotions.SelectionLemmas.Pf.restInd}), a monomial statement
(\lref{Formalism.MainNotions.SelectionLemmas.Pf.restMon}) and the good path weight bound
(\cref{pf:linkG:gpw}), the largest single estimate of the development).
\item The parameters $b(\lambda,e)_j$ are admissible (\lref{Formalism.MainNotions.SelectionLemmas.Pf.restPad}); with (2)
this gives the middle inequality termwise.
\item Summing the segment bounds over the label data gives the first inequality
(\lref{Formalism.MainNotions.SelectionLemmas.Pf.restSum}: a summation over the label data, valid for every family of
functions $G_j$ with the properties of (2), with $C_{rk}=1$).
\end{enumerate}

\emph{Step 5: assembly.} Put $C_0=1088$, $C_{rk}=1$, $C_{\deg}=10$, and let $C_\Lambda,C_k\ge C_0$. Apply Step 2 to $G=\mathrm{Gen}$,
whose elements are supported on the finite set of exponents of multidegree $d$; it gives $S\subseteq NN$ with a complement
and a basis of generators, with $\rank S=\dim\overline{NN}\ge p^{-n'rK}\beta$ by Step 3. For a finite field with $q=p^{n'}$
elements, the plane count of the list gives $|\mathrm{Gr}_r(\F_q^k)|\ge q^{rK}$; if $\beta\le0$ the bound
$\beta/|\mathrm{Gr}_r(\F_q^k)|\le\rank S$ is trivial, and otherwise $\beta/|\mathrm{Gr}_r(\F_q^k)|\le q^{-rK}\beta\le\rank S$. Step 1
gives the degree bound with $C_{\deg}=10$.
\end{proof}

\begin{theorem}[The space statement]\label{pf:linkG:spacestatement}
Assume the hyperplane rank statement, the definite rank-one statement, the rotated sum bound and the third conclusion of
the covariant degree theory (four statements of the development on Gelfand--Tsetlin paths and joint types). Then clause
33 holds: for every uniformised discrete valuation ring $\OO$ over $p$, every valid frame $F$, every segment $j$ with
$\mathrm{seg}_N(F,j)+r\le K$ and all parameters $m\ge0$, $\theta$, $o$, the segment module $XX(F,j,m,\theta,o)$
\begin{enumerate}[label=(\roman*)]
\item is the image under the renaming of the segment of the module spanned by the joint types of the admissible pairs;
\item contains the linearly independent family of the products $\Psi_{\mathrm{Pl}}\Psi_{\mathrm{Tan}}$ of path polynomials
attached to the standard pairs;
\item is saturated, and, if $\OO$ has characteristic $0$, consists of the polynomials whose image over $K_f$ lies in the span of
the joint types;
\item contains every product of a plane polynomial $A$ and a non-plane polynomial $G$ in the letters of the segment which
satisfies the valuation conditions ($\lambda_{1,S}(G)\le M_x$ and $\mathrm{val}_S(A)\ge M_x+m_x$ below every non-$W$ letter
$x$, with the parity convention, and $\mathrm{val}_S(A)\ge\theta_x$) and the degree conditions ($\deg_x\le o_x$);
\item and the same over the residue field: every such product over $\F$ is the reduction of an element of the module.
\end{enumerate}
\formallabel{Formalism.MainNotions.G.Joints.space_statement_of_lemmas}
\formallabel{Formalism.MainNotions.G.Joints.StandardMonomialSpace.space_statement_of_lemmas}
\end{theorem}

\begin{proof}
The statement is equivalent to its version for each discrete valuation ring separately
(\lref{Formalism.MainNotions.G.Joints.StandardMonomialSpace.space_statement_iff_forall_dvr}). Fix $\OO$. The definite
rank-one statement gives the \emph{standard basis data} over $\OO$ and over the residue field: the products of path
polynomials attached to standard pairs form bases of the spaces of polynomials with prescribed weights. From these data and
the hyperplane rank statement, (i)--(iii) follow
(\lref{Formalism.MainNotions.G.Joints.StandardMonomialSpace.space_statement_at_ring_of_basis}): the standard products of
the admissible pairs span the module of (i) and are linearly independent
(\lref{Formalism.MainNotions.G.Joints.StandardMonomialSpace.space_statement_linear_independent}); the module is spanned by
part of a basis of the polynomial ring, hence saturated
(\lref{Formalism.MainNotions.G.Joints.StandardMonomialSpace.space_statement_saturated}); and in characteristic $0$ a
polynomial whose image lies in the rational span of the joint types lies in the integral span by saturation. The membership
criteria (iv) and (v) are the fourth clause
(\lref{Formalism.MainNotions.G.Joints.StandardMonomialSpace.fourth_clause_of_lemmas}), proved from the rotated sum bound
and the last inequality clause (a consequence of the third conclusion): a product $AG$ with the stated valuation and degree
conditions has its joint type among the admissible pairs (a comparison of joint types with bi-polynomial degrees, and a
description of the admissible pairs by the valuations), so it lies in the span of (i), over $\OO$ and over the residue
field.
\end{proof}

% ---------------------------------------------------------------------------
\subsection{From local to global: lengths of lattices of modular forms}\label{sec:pf:linkG:lengths}

In this subsection $R$ is a discrete valuation ring with uniformiser $\pi$, and $\len$ denotes the length of an
$R$-module. We first prove a module-theoretic lemma converting ``all but few'' statements into length bounds, and then
apply it to the lattices of algebraic modular forms of \cref{pf:linkG:gatp}.

\begin{definition}[All but few]\label{pf:linkG:allbut}
Let $M,P$ be submodules of an $R$-module $X$, $a\in R$ and $g,E\in\N$. $\linkGAB_a(M,P,g,E)$ means: there is a finite set
$C\subseteq M$ with $|C|\le g$ such that $a^Ex\in P+\operatorname{span}_R(C)$ for every $x\in M$.
\formal{Formalism.G0.AllButLE}
\end{definition}

\begin{lemma}[Passing to submodules]\label{pf:linkG:invfac}
Let $P\subseteq M'\subseteq M$ be submodules of an $R$-module $X$. Then $\linkGAB_\pi(M,P,g,E)$ implies
$\linkGAB_\pi(M',P,g,E)$.
\formallabel{Formalism.InvariantFactorPaths.G0.invariant_factors_submodule}
\end{lemma}

\begin{proof}
Let $C$ be as in the hypothesis, and let $f:R^C\to X/P$ be $f(a)=\sum_ca_c[c]$. Let $W\subseteq X/P$ be the image of
$M'$ under $x\mapsto[\pi^Ex]$ and $U=f^{-1}(W)\subseteq R^C$. As a submodule of a free module of rank $|C|$ over a
principal ideal domain, $U$ is free with a basis $b_1,\dots,b_n$, $n\le|C|$. Choose $m_i\in M'$ with
$[\pi^Em_i]=f(b_i)$ and put $C'=\{m_1,\dots,m_n\}$. For $x\in M'$ write $\pi^Ex=y+\sum_ca_cc$ with $y\in P$; then
$f(a)=[\pi^Ex]\in W$, so $a\in U$, $a=\sum_ir_ib_i$, and $[\pi^Ex]=\sum_ir_if(b_i)=[\pi^E\sum_ir_im_i]$. Hence
$\pi^Ex\in P+\sum_ir_i\pi^Em_i\subseteq P+\operatorname{span}(C')$.
\end{proof}

\begin{lemma}[Shift]\label{pf:linkG:invshift}
Let $A,B,N_A,N_B$ be submodules of $X$ with $B\subseteq A$, $N_A\subseteq A$ and $\pi^eN_A\subseteq N_B\subseteq B$. Then
$\linkGAB_\pi(A,N_A,g,E)$ implies $\linkGAB_\pi(B,N_B,g,E+e)$.
\formallabel{Formalism.InvariantFactorPaths.G0.inverse_shift}
\end{lemma}

\begin{proof}
By \cref{pf:linkG:invfac} applied to $N_A\subseteq B+N_A\subseteq A$ there is $C\subseteq B+N_A$, $|C|\le g$, with
$\pi^Ey\in N_A+\operatorname{span}(C)$ for $y\in B+N_A$. Put $C''=\pi^eC$. For $c=b+n\in C$ ($b\in B$, $n\in N_A$),
$\pi^ec=\pi^eb+\pi^en\in B+N_B=B$, so $C''\subseteq B$ and $|C''|\le g$. For $x\in B$ write $\pi^Ex=n+z$ with $n\in
N_A$ and $z\in\operatorname{span}(C)$; then $\pi^{E+e}x=\pi^en+\pi^ez\in N_B+\operatorname{span}(C'')$.
\end{proof}

\begin{lemma}[Three length bounds]\label{pf:linkG:lenlemmas}
\begin{enumerate}[label=(\alph*)]
\item If an $R$-module $Q$ is generated by $n$ elements and $\pi^kQ=0$, then $\len(Q)\le kn$.
\item If $E,E'$ are submodules of an $R$-module with $\pi^SE\subseteq E'$, and $E$ is contained in the span of $n$
elements, then $\len\bigl(E/(E\cap E')\bigr)\le Sn$.
\item If $E_0\supseteq E_1\supseteq\dots\supseteq E_V=0$ are submodules, then
$\len(E_0)\le\sum_{j<V}\len(E_j/E_{j+1})$.
\end{enumerate}
\formallabel{Formalism.G0.CyclicSpanLengthLemmas.length_le_mul_of_span_fin_pow_smul_eq_zero}
\formallabel{Formalism.G0.CyclicSpanLengthLemmas.length_le_of_cyclic_pow_smul_eq_zero}
\formallabel{Formalism.G0.CyclicSpanLengthLemmas.length_quot_comap_le_mul_of_pow_smul_mem}
\formallabel{Formalism.G0.CyclicSpanLengthLemmas.length_le_sum_of_antitone_chain_quot_le}
\end{lemma}

\begin{proof}
(a) For a cyclic module $Rx$ with $\pi^kx=0$, induction on $k$: $\len(Rx)=\len(R\pi x)+\len(Rx/R\pi x)\le(k-1)+1$. For
$n$ generators, induction on $n$: $\len(Q)=\len(Rg_1)+\len(Q/Rg_1)\le k+k(n-1)$. (b) Let $\psi:R^n\to E''$ be the map to
the ambient module sending the unit vectors to the $n$ elements. The submodule $\psi^{-1}(E)$ of $R^n$ is free of rank
at most $n$ and maps onto $E$, so $E/(E\cap E')$ is generated by at most $n$ elements and is killed by $\pi^S$; apply
(a). (c) Induction on $V$, using $\len(E_0)=\len(E_1)+\len(E_0/E_1)$.
\end{proof}

\begin{lemma}[Lengths from all-but statements]\label{pf:linkG:lengthpart2}
Let $M\subseteq X$ be a free $R$-module of finite rank $d$, $N\subseteq X$ a submodule, $\tau,S,V\in\N$ and $0\le s\le
\frac12$. Assume
\begin{enumerate}[label=(\roman*)]
\item for every $\beta<V$ there is a finite set $C_\beta\subseteq M$ with $|C_\beta|\le ds^{\beta+1}$ and
$\pi^{\tau+\beta S}x\in N+\operatorname{span}(C_\beta)$ for all $x\in M$;
\item $\pi^{VS}M\subseteq N$.
\end{enumerate}
Then $\len\bigl(M/(N\cap\pi^\tau M)\bigr)\le d(\tau+2Ss)$.
\formallabel{Formalism.G0.CyclicSpanLengthLemmas.lengthPart2}
\end{lemma}

\begin{proof}
Put $K=N\cap\pi^\tau M\subseteq M$, $Q=M/K$ and $A_0=\pi^\tau Q$. Since $K\subseteq\pi^\tau M$, $Q/A_0=M/\pi^\tau M$, which
is generated by $d$ elements and killed by $\pi^\tau$; so $\len(Q/A_0)\le\tau d$ by \cref{pf:linkG:lenlemmas}(a).
Let $\theta:Q\to X/N$ be induced by the inclusion $M\subseteq X$ (well defined as $K\subseteq N$), and let
$E_j\subseteq X/N$ be the image of $\pi^{\tau+jS}M$. The restriction of $\theta$ to $A_0$ is injective with image $E_0$:
if $\theta(\pi^\tau x)=0$ then $\pi^\tau x\in N\cap\pi^\tau M=K$. Hence $\len(A_0)=\len(E_0)$. We have
$E_{j+1}\subseteq E_j$, $\pi^SE_j\subseteq E_{j+1}$, $E_V=0$ by (ii), and by (i) $E_j$ is contained in the span of the
images of the elements of $C_j$ for $j<V$. By \cref{pf:linkG:lenlemmas}(b) and (c),
\[
\len(E_0)\le\sum_{j<V}S|C_j|\le Sd\sum_{j<V}s^{j+1}\le2Sds,
\]
since $\sum_{j<V}s^{j+1}\le s/(1-s)\le2s$ for $0\le s\le\frac12$. Therefore $\len(Q)=\len(A_0)+\len(Q/A_0)\le
d\tau+2Sds$.
\end{proof}

\begin{corollary}\label{pf:linkG:lengthallbut}
Let $a\in R$ be non-zero and assume the hypotheses of \cref{pf:linkG:lengthpart2} with $a$ in place of $\pi$. Then
$\len\bigl(M/(N\cap a^\tau M)\bigr)\le\len(R/aR)\cdot d(\tau+2Ss)$.
\formallabel{Formalism.MainNotions.G.LatticeSystems.lengthFromAllBut}
\formallabel{Formalism.MainNotions.G.LatticeSystems.ln_length_pow}
\end{corollary}

\begin{proof}
Write $a=u\pi^\epsilon$ with $u$ a unit; then $\len(R/aR)=\epsilon$ and $a^jM=\pi^{\epsilon j}M$. Apply
\cref{pf:linkG:lengthpart2} with $(\epsilon\tau,\epsilon S)$ in place of $(\tau,S)$.
\end{proof}

\begin{lemma}[Elementary symmetric functions]\label{pf:linkG:esymm}
Let $\kappa_v\ge0$ for $v$ in a finite set $VV$, and let $e_k(\kappa)=\sum_{B\subseteq VV,\ |B|=k}\prod_{v\in B}\kappa_v$.
Then $d\,m^ke_k(\kappa)\le d\,(m\sum_v\kappa_v)^k/k!$ for all $d,m,k\in\N$.
\formallabel{Formalism.MainNotions.G.LatticeSystems.esymm_clause}
\end{lemma}

\begin{proof}
Expanding $(\sum_v\kappa_v)^k$ as a sum over $k$-tuples, every $k$-subset $B$ occurs $k!$ times (as the tuples of
distinct elements enumerating $B$), and all other terms are non-negative; so $k!\,e_k(\kappa)\le(\sum_v\kappa_v)^k$.
\end{proof}

\paragraph{Hecke products.}
Let $\mathcal G$ be a local group datum (\cref{pf:linkG:gatp}) over $p$, with modules $M\subseteq M^{\mathrm{na}}$,
$d=\rank M$, $d^{\mathrm{na}}=\rank M^{\mathrm{na}}$ and Hecke operators $y_v$. For a set $VV$ of places, $m\in\N$ and
$S_v\in\N$ ($v\in VV$), and a submodule $M_0$, let $N(M_0)=N_{VV,m,S}(M_0)$ be the span of the elements
\[
p^{\sum_{v\in VV\smallsetminus B}S_v}\Bigl(\prod_{v\in B}y_v^m\Bigr)F,\qquad B\subseteq VV,\ F\in M_0,
\]
and for $t:VV\to\N$ and $D\in\N$ put $\tau(VV,m,D,t)=\sum_{v\in VV}m(t_v+D+1)$. The Hecke operators commute with each
other (\lref{Formalism.MainNotions.G.LatticeSystems.exponent_condition_holds}).

\begin{definition}[Slot systems]\label{pf:linkG:slots}
A \emph{slot assignment} is a map $\sigma$ from the places of $\mathcal G$ to $\{\mathrm{base},\mathrm{ball},\mathrm{dual}\}$;
$\sigma[v\mapsto s]$ is $\sigma$ with the value at $v$ replaced by $s$, and for $B\subseteq VV$ the \emph{vertex} $\sigma_B$
takes the value $\mathrm{dual}$ on $B$ and $\mathrm{base}$ elsewhere. Let $t$, $D$ and $\kappa_v$ ($v\in VV$) be given.
A \emph{slot system} is a family of $\OO$-submodules $\mathbf L(\sigma)$ of the space of functions of $\mathcal G$ with
$\mathbf L(\sigma_\emptyset)=M^{\mathrm{na}}$ such that for all $v\in VV$ and all $\sigma$:
\begin{gather*}
\mathbf L(\sigma[v\mapsto\mathrm{base}])\subseteq\mathbf L(\sigma[v\mapsto\mathrm{ball}]),\qquad
\mathbf L(\sigma[v\mapsto\mathrm{dual}])\subseteq\mathbf L(\sigma[v\mapsto\mathrm{base}]),\\
\varpi^{t_v}\mathbf L(\sigma[v\mapsto\mathrm{ball}])\subseteq\mathbf L(\sigma[v\mapsto\mathrm{base}]),\qquad
\varpi^{t_v}\mathbf L(\sigma[v\mapsto\mathrm{base}])\subseteq\mathbf L(\sigma[v\mapsto\mathrm{dual}]),\\
y_v\mathbf L(\sigma[v\mapsto\mathrm{ball}])\subseteq\mathbf L(\sigma[v\mapsto\mathrm{dual}]),
\end{gather*}
and (the \emph{one-step statement}) for every non-empty $B\subseteq VV$ there is a finite set $C\subseteq\mathbf L(\sigma_B)$
with $|C|\le d^{\mathrm{na}}\prod_{v\in B}\kappa_v$ and
\[
p^{D+1}\mathbf L(\sigma_B)\subseteq\sum_{v\in B}\operatorname{span}\,y_v\bigl(\mathbf L(\sigma_B[v\mapsto\mathrm{ball}])\bigr)+\operatorname{span}(C).
\]
For a slot assignment $\sigma$ and a tuple of lattices $\mathbf L=(L_v)$ let $T_\sigma(\mathbf L)\subseteq\mathcal X$ be the
tensor product over the places $v$ of the \emph{slot lattices} $\linkGH_{L_v}$, $\linkGNs(t_v,L_v)$ or
$\{x\in\linkGV:\ G(x)\in\OO\text{ for all }G\in\linkGNt(t_v,L_v)\}$, according as $\sigma(v)$ is $\mathrm{base}$, $\mathrm{ball}$
or $\mathrm{dual}$. The \emph{local slot modules} are the modules $\mathcal L(\sigma)$ of functions $F$ as in
\cref{pf:linkG:gatp} (equivariant, and zero on tuples which are not tuples of lattices) with $F(i,\mathbf L)\in
T_\sigma(\mathbf L)$ for every component $i$ and every tuple of lattices $\mathbf L$; thus $\mathcal L(\sigma_\emptyset)=M^{\mathrm{na}}$.
\formal{Formalism.MainNotions.G.LatticeSystems.RestWalk}
\formal{Formalism.MainNotions.G.LatticeSystems.latSys}
\end{definition}

\paragraph{The walk lemma.}
The first ingredient is an abstract statement on modules indexed by slot assignments. For $V\in\N$ write
$[V]=\{0,\dots,V-1\}$; slot assignments $\sigma:[V]\to\{\mathrm{base},\mathrm{ball},\mathrm{dual}\}$, $\sigma[v\mapsto s]$ and
the vertices $\sigma_B$ ($B\subseteq[V]$) are defined as in \cref{pf:linkG:slots}.

\begin{theorem}[Walk lemma]\label{pf:linkG:walkabs}
Let $R$ be a commutative ring, $X$ an $R$-module, $\pi\in R$ and $V\in\N$. Let $L(\sigma)\subseteq X$ be submodules indexed by
the slot assignments $\sigma$ of $[V]$, let $y_0,\dots,y_{V-1}$ be pairwise commuting endomorphisms of $X$, and let $m\ge1$,
$\beta,c_0,S\in\N$ and $a_v,b_v,c_v,S_v\in\N$ with $S_v\le S$ ($v<V$). Assume that for all $\sigma$ and $v<V$
\begin{gather*}
L(\sigma[v\mapsto\mathrm{base}])\subseteq L(\sigma[v\mapsto\mathrm{ball}]),\qquad
L(\sigma[v\mapsto\mathrm{dual}])\subseteq L(\sigma[v\mapsto\mathrm{base}]),\\
\pi^{a_v}L(\sigma[v\mapsto\mathrm{ball}])\subseteq L(\sigma[v\mapsto\mathrm{base}]),\qquad
\pi^{b_v}L(\sigma[v\mapsto\mathrm{base}])\subseteq L(\sigma[v\mapsto\mathrm{dual}]),\\
\pi^{c_v}L(\sigma[v\mapsto\mathrm{ball}])\subseteq L(\sigma[v\mapsto\mathrm{dual}]),\qquad
y_vL(\sigma[v\mapsto\mathrm{ball}])\subseteq L(\sigma[v\mapsto\mathrm{dual}]),
\end{gather*}
and that for every $B\subseteq[V]$ with $|B|=\beta+1$ there is a finite set $G(B)\subseteq L(\sigma_B)$ with $|G(B)|\le
g(B)$ and
\[
\pi^{c_0}L(\sigma_B)\subseteq\sum_{v\in B}y_v\bigl(L(\sigma_B[v\mapsto\mathrm{ball}])\bigr)+\operatorname{span}G(B).
\]
Let $N$ be the submodule spanned by the elements $\pi^{\sum_{v\notin B'}S_v}\bigl(\prod_{v\in B'}y_v^m\bigr)x$ with
$B'\subseteq[V]$ and $x\in L(\sigma_\emptyset)$. Then, with $E=\sum_{v<V}\bigl(a_v+b_v+(m-1)c_v+mc_0\bigr)+\beta S$,
\[
\linkGAB_\pi\Bigl(L(\sigma_\emptyset),N,\sum_{|B|=\beta+1}m^{\beta+1}g(B),E\Bigr)\qquad\text{and}\qquad
\linkGAB_\pi\Bigl(L(\sigma_\emptyset),N,0,\sum_{v<V}S_v\Bigr).
\]
\formallabel{Formalism.G0.WalkSharpP.main}
\formallabel{Formalism.G0.walkSharp_true}
\formallabel{Formalism.G0.Top.walkSharp}
\end{theorem}

\begin{proof}
Write $M_0=L(\sigma_\emptyset)$ and $f_v=a_v+b_v+(m-1)c_v+mc_0$, so $E=\sum_vf_v+\beta S$. Every $y_u$ preserves $M_0$, since
$M_0\subseteq L(\sigma_\emptyset[u\mapsto\mathrm{ball}])$ and $y_u$ maps this module into
$L(\sigma_\emptyset[u\mapsto\mathrm{dual}])\subseteq M_0$; likewise $L(\sigma_B)\subseteq M_0$ for every $B$ (replace the slot
$\mathrm{dual}$ by $\mathrm{base}$ at the places of $B$ one after the other). Taking $B'=\emptyset$ gives
$\pi^{\sum_vS_v}M_0\subseteq N$, which is the second statement (with the empty set). If $V<\beta+1$ there is no $B$ with
$|B|=\beta+1$, and $\sum_vS_v\le VS\le\beta S\le E$, so $\pi^EM_0\subseteq N$ and the first statement holds with the empty
set. Assume from now on $V\ge\beta+1$.

\emph{States.} A \emph{state} is a pair $(B,i)$ with $B\subseteq[V]$, $|B|=\beta+1$, and $i:B\to\{0,\dots,m-1\}$. Let
$\top(B)=\max B$. A place $u$ is \emph{done} in the state if $u<\top(B)$ and $u\notin B$, and \emph{active} if $u\in B$; the
places $u>\top(B)$ are \emph{future}. Put
\[
Y_{B,i}=\prod_{u\ \mathrm{done}}y_u^m\prod_{u\in B}y_u^{i_u},\qquad
\mathrm{acc}(B,i)=\sum_{u\ \mathrm{done}}f_u+\sum_{u\in B}\bigl(b_u+i_u(c_u+c_0)\bigr),
\]
and let $\mathrm{ex}(B,i)\le mV$ be the total exponent $m\cdot\#\{\text{done places}\}+\sum_{u\in B}i_u$. Since
$f_u-b_u-i_u(c_u+c_0)\ge a_u+(m-i_u)c_0$ for $u\in B$ and future places contribute $f_u\ge0$ to $\sum_vf_v$, we have
\begin{equation}\label{eq:pf:linkG:walkcost}
\mathrm{acc}(B,i)+c_0\le\sum_vf_v,\qquad\text{and}\qquad\mathrm{acc}(B,i)+a_v+c_0\le\sum_uf_u\ \text{ if }v\in B,\ i_v=m-1.
\end{equation}
Let $C=\bigcup_{(B,i)}Y_{B,i}(G(B))$. Then $C\subseteq M_0$, and since there are $m^{\beta+1}$ maps $i$ for each $B$,
$|C|\le\sum_{|B|=\beta+1}m^{\beta+1}g(B)$.

\emph{Claim: for every state $(B,i)$ and every $z\in L(\sigma_B)$, $\pi^{E-\mathrm{acc}(B,i)}Y_{B,i}z\in N+\operatorname{span}(C)$.}
We argue by descending induction on $\mathrm{ex}(B,i)$. By hypothesis $\pi^{c_0}z=\sum_{v\in B}y_vz_v+g$ with
$z_v\in L(\sigma_B[v\mapsto\mathrm{ball}])$ and $g\in\operatorname{span}G(B)$. By \eqref{eq:pf:linkG:walkcost},
$E-\mathrm{acc}(B,i)-c_0\ge0$ and
\[
\pi^{E-\mathrm{acc}(B,i)}Y_{B,i}z=\pi^{E-\mathrm{acc}(B,i)-c_0}Y_{B,i}\,g+\sum_{v\in B}\pi^{E-\mathrm{acc}(B,i)-c_0}Y_{B,i}\,y_vz_v .
\]
The first term lies in $\operatorname{span}(C)$. Fix $v\in B$; there are three cases.
\begin{enumerate}[label=(\roman*)]
\item $i_v+1<m$. Let $i'$ agree with $i$ except $i'_v=i_v+1$. Then $Y_{B,i}y_v=Y_{B,i'}$,
$\mathrm{ex}(B,i')=\mathrm{ex}(B,i)+1$, $\mathrm{acc}(B,i')=\mathrm{acc}(B,i)+c_v+c_0\le E$ by \eqref{eq:pf:linkG:walkcost}, and
$\pi^{c_v}z_v\in L(\sigma_B)$ because $\sigma_B(v)=\mathrm{dual}$. The term equals
$\pi^{E-\mathrm{acc}(B,i')}Y_{B,i'}(\pi^{c_v}z_v)$, which lies in $N+\operatorname{span}(C)$ by induction.
\item $i_v=m-1$ and $w=\top(B)+1<V$. Let $B'=(B\smallsetminus\{v\})\cup\{w\}$, $i'=i$ on $B\smallsetminus\{v\}$ and $i'_w=0$. In
the new state $v$ is done and no other place changes its role except $w$, so $Y_{B,i}y_v=Y_{B',i'}$,
$\mathrm{ex}(B',i')=\mathrm{ex}(B,i)+1$ and $\mathrm{acc}(B',i')=\mathrm{acc}(B,i)+a_v+c_0+b_w\le E$. The element
$z'=\pi^{b_w}\pi^{a_v}z_v$ lies in $L(\sigma_{B'})$: $\pi^{a_v}z_v\in L(\sigma_B[v\mapsto\mathrm{base}])$, which has the slot
$\mathrm{base}$ at $w$, and $\pi^{b_w}$ moves the slot at $w$ to $\mathrm{dual}$. The term equals
$\pi^{E-\mathrm{acc}(B',i')}Y_{B',i'}z'$, which lies in $N+\operatorname{span}(C)$ by induction.
\item $i_v=m-1$ and $\top(B)=V-1$. Then every place outside $B$ is done, so
$Y_{B,i}y_v=\prod_{u\notin B\smallsetminus\{v\}}y_u^m\prod_{u\in B\smallsetminus\{v\}}y_u^{i_u}$. The element
$x=\prod_{u\in B\smallsetminus\{v\}}y_u^{i_u}(\pi^{a_v}z_v)$ lies in $M_0$, since $\pi^{a_v}z_v\in L(\sigma_B[v\mapsto\mathrm{base}])
\subseteq M_0$. By \eqref{eq:pf:linkG:walkcost}, $E-\mathrm{acc}(B,i)-c_0-a_v\ge\beta S\ge\sum_{u\in B\smallsetminus\{v\}}S_u$, as
$|B\smallsetminus\{v\}|=\beta$. Hence the term is a multiple of
$\pi^{\sum_{u\in B\smallsetminus\{v\}}S_u}\prod_{u\notin B\smallsetminus\{v\}}y_u^m\,x$, a generator of $N$ (with
$B'=[V]\smallsetminus(B\smallsetminus\{v\})$).
\end{enumerate}
This proves the claim. Finally let $B_0=\{0,\dots,\beta\}$ and $i=0$. No place is done, so $Y_{B_0,0}=1$ and
$\mathrm{acc}(B_0,0)=\sum_{u\in B_0}b_u$. For $x\in M_0$ the element $z=\pi^{\sum_{u\in B_0}b_u}x$ lies in $L(\sigma_{B_0})$
(apply the inclusions $\pi^{b_u}$, $u\in B_0$, one after the other), and the claim gives $\pi^Ex\in N+\operatorname{span}(C)$.
\end{proof}

\begin{corollary}[Walk lemma for slot systems]\label{pf:linkG:walk}
Let $\mathcal G$ be a local group datum whose Hecke operators commute and with $\varpi^2=p$, let $VV$, $t$, $D$ and
$\kappa_v$ be as in \cref{pf:linkG:slots}, $m\ge1$, and $S_v\le S$ for $v\in VV$. If a slot system exists, then for every
$\beta\ge0$
\begin{equation}\label{eq:pf:linkG:manyorbits}
\linkGAB_p\bigl(M^{\mathrm{na}},N(M^{\mathrm{na}}),\lfloor d^{\mathrm{na}}m^{\beta+1}e_{\beta+1}(\kappa)\rfloor,\tau(VV,m,D,t)+\beta S\bigr).
\end{equation}
\formallabel{Formalism.MainNotions.G.LatticeSystems.walk_clause}
\formallabel{Formalism.MainNotions.G.LatticeSystems.walk_allBut_bound}
\end{corollary}

\begin{proof}
Enumerate $VV=\{v_0<\dots<v_{V-1}\}$, and for a slot assignment $\sigma$ of $[V]$ let $L(\sigma)=\mathbf L(\tilde\sigma)$,
where $\tilde\sigma(v_j)=\sigma(j)$ and $\tilde\sigma=\mathrm{base}$ outside $VV$; thus $L(\sigma_\emptyset)=M^{\mathrm{na}}$ and
$L(\sigma_B)=\mathbf L(\sigma_{\tilde B})$ with $\tilde B=\{v_j:j\in B\}$. Apply \cref{pf:linkG:walkabs} with $\pi=\varpi$,
$y_j=y_{v_j}$, $a_j=b_j=t_{v_j}$, $c_j=2t_{v_j}$, $c_0=2(D+1)$, $2S$ and $2S_{v_j}$ in place of $S$ and $S_j$, and $G(B)$ the
set $C$ of the one-step statement for $\tilde B$, $g(B)=|G(B)|$. The inclusions with $a_j,b_j$ and the Hecke inclusions are
those of the slot system; the inclusion with $c_j$ is the composite of the inclusions with $a_j$ and $b_j$; and the one-step
statement is the remaining hypothesis, because $\varpi^{2(D+1)}=p^{D+1}$. Since $\varpi^{2S_v}=p^{S_v}$, the module $N$ of the
theorem is $N(M^{\mathrm{na}})$. The exponent is
$\sum_{v\in VV}\bigl(2t_v+2(m-1)t_v+2m(D+1)\bigr)+2\beta S=2(\tau(VV,m,D,t)+\beta S)$, and
$\varpi^{2(\tau+\beta S)}=p^{\tau+\beta S}$. Finally
$\sum_{|B|=\beta+1}m^{\beta+1}|G(B)|\le d^{\mathrm{na}}m^{\beta+1}\sum_{|B|=\beta+1}\prod_{v\in\tilde B}\kappa_v
=d^{\mathrm{na}}m^{\beta+1}e_{\beta+1}(\kappa)$, and the left side is an integer.
\end{proof}

\paragraph{Lemma J.}
The inclusions of a slot system hold for the local slot modules: they follow from the local inclusions
$\linkGH_{e_s}\subseteq\linkGR_e$, $\linkGH_{e_t}^\vee\subseteq\linkGR'_e$ and the reflexivity of $\linkGH_L$, and from the fact
that $y_v$ maps the slot $\mathrm{ball}$ at $v$ to the slot $\mathrm{dual}$
(\lref{Formalism.MainNotions.G.LatticeSystems.restLocD} and \lref{Formalism.MainNotions.G.LatticeSystems.restWalk_of_local_J}:
these local facts hold at every place, and together with the one-step statement they give a slot system). The one-step
statement is lemma J.

\begin{theorem}[Lemma J]\label{pf:linkG:lemmaJ}
Let $\mathcal G$ be a local group datum with trivial stabilisers and finite classes, $VV$ a non-empty set of places,
$t:VV\to\N$, $D,c_*\in\N$ and $\kappa_v\in\R$ such that $\linkGLP(t_v,D,\kappa_v,c_*)$ holds at every $v\in VV$ (for the Hida
setting, $r$ and the weight $\lambda_v$ of $v$). Then for every non-empty $B\subseteq VV$ there is a finite set
$C\subseteq\mathcal L(\sigma_B)$ with $|C|\le d^{\mathrm{na}}\prod_{v\in B}\kappa_v$ and
\[
p^{D+1}\mathcal L(\sigma_B)\subseteq J_B+\operatorname{span}(C),\qquad
J_B=\sum_{v\in B}\operatorname{span}\,y_v\bigl(\mathcal L(\sigma_B[v\mapsto\mathrm{ball}])\bigr).
\]
\formallabel{Formalism.MainNotions.G.LatticeSystems.restJAll}
\formallabel{Formalism.MainNotions.G.LatticeSystems.restJAll_of}
\formallabel{Formalism.MainNotions.G.LatticeSystems.ja_restJ}
\formallabel{Formalism.MainNotions.G.LatticeSystems.jDelta}
\formallabel{Formalism.MainNotions.G.LatticeSystems.jRank}
\formallabel{Formalism.MainNotions.G.LatticeSystems.jSmith}
\formallabel{Formalism.MainNotions.G.LatticeSystems.jm_core}
\formallabel{Formalism.MainNotions.G.LatticeSystems.jm_count}
\formallabel{Formalism.MainNotions.G.LatticeSystems.jStep1}
\formallabel{Formalism.MainNotions.G.LatticeSystems.jl_good}
\formallabel{Formalism.MainNotions.G.LatticeSystems.jLocal}
\end{theorem}

\begin{proof}
Fix $B$, write $M_B=\mathcal L(\sigma_B)$, $A_v=A_r(\lambda_v)$, and $\dim_v=\dim\linkGV(\lambda_v)$.

\emph{Step 1: delta functions.} The slot lattices are equivariant: for a component $i$, $\gamma$ in its level, a tuple of
lattices $\mathbf L$ and $\xi\in T_\sigma(\mathbf L)$ we have $\chi(\gamma)\gamma\xi\in T_\sigma(\gamma\mathbf L)$
(\lref{Formalism.MainNotions.G.LatticeSystems.jEquiv}: the local lattices $\linkGH_L$, $\linkGNs(t,L)$, $\linkGNt(t,L)$ are
built from $L$ and the Hecke edges at $L$ compatibly with the action of $\GL_{m+1}(Q_v)$, which permutes the Hecke edges of
type $r$). For a component $i$, a tuple of lattices $\mathbf L_0$ and $\xi\in\mathcal X$ let $\delta_{i,\mathbf L_0}(\xi)$ be the
function with value $\chi(\gamma)\gamma\xi$ at $(i,\gamma\mathbf L_0)$ and $0$ at all pairs (component, tuple) not of this
form. By trivial stabilisers $\gamma$ is determined by $\gamma\mathbf L_0$, so $\delta_{i,\mathbf L_0}(\xi)$ is well defined; it
is $\OO$-linear in $\xi$, equivariant, $\delta_{i,\mathbf L_0}(\xi)(i,\mathbf L_0)=\xi$, and by equivariance of the slot lattices
$\delta_{i,\mathbf L_0}(\xi)\in\mathcal L(\sigma)$ whenever $\xi\in T_\sigma(\mathbf L_0)$. By finite classes there is a finite set
$\mathcal R$ of representatives $(i,\mathbf L)$ of the orbits of the levels on pairs (component, tuple of lattices), and
$F=\sum_{(i,\mathbf L)\in\mathcal R}\delta_{i,\mathbf L}(F(i,\mathbf L))$ for every $F\in\mathcal L(\sigma)$, as both sides are
equivariant and agree at the representatives.

\emph{Step 2: counting representatives.} We show $|\mathcal R|\prod_v\dim_v\le d^{\mathrm{na}}$. For $(i,\mathbf L)\in\mathcal R$
the lattice $T_{\sigma_\emptyset}(\mathbf L)=\linkGH(\mathbf L)$, a tensor product of free lattices of ranks $\dim_v$, contains
$\prod_v\dim_v$ linearly independent elements $Q_{i,\mathbf L,j}$. The functions $\delta_{i,\mathbf L}(Q_{i,\mathbf L,j})$ lie in
$M^{\mathrm{na}}$ and are linearly independent: evaluating a vanishing combination at a representative $(i_0,\mathbf L_0)$ kills
the terms of the other representatives (which lie in other orbits) and leaves a vanishing combination of the
$Q_{i_0,\mathbf L_0,j}$. Since $M^{\mathrm{na}}$ is free of rank $d^{\mathrm{na}}$, the bound follows.

\emph{Step 3: Smith normal form at one place.} Let $v\in VV$, $L$ a lattice at $v$, and let $P$ be the dual slot lattice at $L$.
We show: there are Hecke edges $(L''_k\subseteq L)$ of type $r$ and $F_k\in\linkGNs(t_v,L''_k)$ ($k\le K$) such that
$p^{A_v}F_k\in\varpi^{c_*}P(L')$ for every target $L'\ne L$ of a Hecke edge of type $r$ with source $L''_k$ (where $P(L')$ is
the dual slot lattice at $L'$), a basis $(b_a)$ of $P$ and a set $\mathrm{bad}$ of at most $\kappa_v\dim_v$ indices with
$p^Db_a\in\operatorname{span}\{p^{A_v}F_k\}$ for $a\notin\mathrm{bad}$. The property $\linkGLPat(t_v,D,\kappa_v,c_*,L)$ provides
the edges, the $F_k$ (with $F_k\in\varpi^{c_*}\linkGR_{e'}$ for the other edges $e'$ with source $L''_k$), an integer
$n\ge(1-\kappa_v)\dim_v$, a linear map $u:\OO^n\to\linkGNt(t_v,L)$ and a matrix $(M_{jk})$ with
$\sum_kM_{jk}\,u(x)(f_k)=p^Dx_j$, where $f_k=p^{A_v}F_k$. The local lattice chain gives
$\varpi^{t_v}\linkGH_L\subseteq P\subseteq\linkGH_L$, so $P$ is free of rank $\dim_v$, and $p^{A_v}\linkGR_e\subseteq P(e_t)$ for
every Hecke edge $e$ (\lref{Formalism.MainNotions.G.LatticeSystems.ch_locChain} and \lref{Formalism.MainNotions.G.LatticeSystems.jm_Re}:
consequences of the local inclusions above); hence $f_k\in P$ and $f_k\in\varpi^{c_*}P(L')$ for the other targets $L'$. Take
a Smith normal form of $\operatorname{span}\{f_k\}\subseteq P$: a basis $(b_a)$ of $P$, an injection $\iota$ of an index set
$\{1,\dots,n_S\}$ into the indices of the basis and $\alpha_l\in\OO$ such that $(\alpha_lb_{\iota(l)})_l$ is a basis of
$\operatorname{span}\{f_k\}$. Writing $\sum_kM_{jk}f_k=\sum_lc_{jl}\alpha_lb_{\iota(l)}$ and $X_{lk'}\in\OO$ for the value of
$u(e_{k'})$ at $b_{\iota(l)}$ (integral, as $\linkGNt$ pairs integrally with $P$), the identity gives the matrix equation
$\sum_lc_{jl}\alpha_lX_{lk'}=p^D\delta_{jk'}$. Let $\mathrm{Good}=\{l:\alpha_l\mid p^D\}$. For $l\notin\mathrm{Good}$ we have
$\alpha_l=p^D\mu_l$ with $\mu_l$ in the maximal ideal $\mathfrak m$ of $\OO$ ($\OO$ is a discrete valuation ring), so the matrix
equation reads $\sum_{l\in\mathrm{Good}}c_{jl}\alpha_lX_{lk'}=p^D(1-Z)_{jk'}$, where $Z=\bigl(\sum_{l\notin\mathrm{Good}}c_{jl}\mu_l
X_{lk'}\bigr)$ has entries in $\mathfrak m$. As $1-Z$ is invertible over the local ring $\OO$, the matrix $p^D\cdot1_n$ factors
through $\OO^{\mathrm{Good}}$, and comparing ranks over $K_f$ gives $n\le|\mathrm{Good}|$. Put $\mathrm{bad}$ equal to the
complement of $\iota(\mathrm{Good})$; then
$|\mathrm{bad}|\le\dim_v-n\le\kappa_v\dim_v$, and for $a=\iota(l)$ with $l\in\mathrm{Good}$,
$p^Db_a=(p^D/\alpha_l)\,\alpha_lb_a\in\operatorname{span}\{f_k\}$.

\emph{Step 4: the Hecke step.} Let $v\in B$, $i$ a component, $\mathbf L$ a tuple of lattices, $(L''\subseteq L_v)$ a Hecke edge
of type $r$, $E\in\N$, and $\xi\in T_{\sigma_B[v\mapsto\mathrm{ball}]}(\mathbf L[v\mapsto L''])$ such that
$p^{A_v}\xi\in p^ET_{\sigma_B}(\mathbf L[v\mapsto L'])$ for every target $L'\ne L_v$ of a Hecke edge of type $r$ with source
$L''$. We show $\delta_{i,\mathbf L}(p^{A_v}\xi)\in J_B+p^EM_B$. Let $F=\delta_{i,\mathbf L[v\mapsto L'']}(\xi)$; by Step 1,
$F\in\mathcal L(\sigma_B[v\mapsto\mathrm{ball}])$, so $y_vF\in J_B$, and $y_vF\in M_B$ (the Hecke operator maps the slot
$\mathrm{ball}$ at $v$ to $\mathrm{dual}$). Also $\delta_{i,\mathbf L}(p^{A_v}\xi)\in M_B$, because $p^{A_v}\xi\in
T_{\sigma_B}(\mathbf L)$: the factor of $\xi$ at $v$ lies in $\linkGNs(t_v,L'')\subseteq\linkGR_e$ for the edge $e=(L''\subseteq
L_v)$, and $p^{A_v}\linkGR_e$ lies in the dual slot lattice at $L_v$ (Step 3). We compare the values of the two functions at a
tuple $\mathbf L_1$ and the component $i$. By definition,
$(y_vF)(i,\mathbf L_1)=p^{A_v}\sum_AF(i,\mathbf L_1[v\mapsto A])$ over the sources $A$ of the Hecke edges $(A\subseteq L_{1,v})$. A
term is non-zero only if $\mathbf L_1[v\mapsto A]=\gamma(\mathbf L[v\mapsto L''])$ for some $\gamma$; then $A=\gamma L''$, and
$\gamma^{-1}$ transports the edge $(A\subseteq L_{1,v})$ to a Hecke edge $(L''\subseteq L')$ with $L'=\gamma^{-1}L_{1,v}$. If
$L'=L_v$, then $\mathbf L_1=\gamma\mathbf L$ and the term is $\chi(\gamma)\gamma(p^{A_v}\xi)=\delta_{i,\mathbf L}(p^{A_v}\xi)(i,\mathbf L_1)$;
by trivial stabilisers there is at most one such term, and conversely, if $\mathbf L_1=\gamma\mathbf L$, the edge
$(\gamma L''\subseteq L_{1,v})$ contributes it. If $L'\ne L_v$, then $p^{A_v}\xi=p^E\eta$ with $\eta\in T_{\sigma_B}(\mathbf
L[v\mapsto L'])$, and the term $\chi(\gamma)\gamma(p^E\eta)$ lies in $p^ET_{\sigma_B}(\mathbf L_1)$ by equivariance. At the other
components both functions vanish. Hence all
values of $\delta_{i,\mathbf L}(p^{A_v}\xi)-y_vF$ lie in $p^ET_{\sigma_B}(\cdot)$; as $\OO$ is a domain, a function of $M_B$
with this property is $p^E$ times a function of $M_B$. So $\delta_{i,\mathbf L}(p^{A_v}\xi)=y_vF+\bigl(\delta_{i,\mathbf
L}(p^{A_v}\xi)-y_vF\bigr)\in J_B+p^EM_B$.

\emph{Step 5: the local statement.} Let $i$ be a component and $\mathbf L$ a tuple of lattices. We show: there is
$C_{i,\mathbf L}\subseteq T_{\sigma_B}(\mathbf L)$ with $|C_{i,\mathbf L}|\le\prod_{v\in B}\kappa_v\prod_v\dim_v$ such that every
$\xi\in T_{\sigma_B}(\mathbf L)$ satisfies $p^D\xi=g+c$ with $\delta_{i,\mathbf L}(g)\in J_B+p^{D+1}M_B$ and
$c\in\operatorname{span}(C_{i,\mathbf L})$. For $v\in B$ take the basis $(b_{v,a})$ of the dual slot lattice at $L_v$ and the set
$\mathrm{bad}_v$ of Step 3; for $v\notin B$ take a spanning family $(b_{v,a})$ of $\linkGH_{L_v}$ with $\dim_v$ members
(\lref{Formalism.MainNotions.G.LatticeSystems.jBase}). Then $T_{\sigma_B}(\mathbf L)$ is spanned by the products
$b_{\mathbf a}=\bigotimes_vb_{v,a_v}$. Let $C_{i,\mathbf L}$ be the set of $b_{\mathbf a}$ with $a_v\in\mathrm{bad}_v$ for all
$v\in B$; its size is at most $\prod_{v\in B}\kappa_v\dim_v\prod_{v\notin B}\dim_v$. If $\mathbf a$ has $a_v\notin\mathrm{bad}_v$
for some $v\in B$, then $p^Db_{v,a_v}=\sum_kr_k\,p^{A_v}F_k$ by Step 3, and $\delta_{i,\mathbf L}(p^Db_{\mathbf a})$ is the sum of
the terms $r_k\,\delta_{i,\mathbf L}\bigl(p^{A_v}\xi_k\bigr)$, $\xi_k=F_k\otimes\bigotimes_{u\ne v}b_{u,a_u}\in
T_{\sigma_B[v\mapsto\mathrm{ball}]}(\mathbf L[v\mapsto L''_k])$. For a target $L'\ne L_v$ of an edge with source $L''_k$ we have
$p^{A_v}F_k=\varpi^{c_*}y$ with $y$ in the dual slot lattice at $L'$, hence, using $\varpi^2=p$ and $c_*\ge4(D+1)\ge2(D+1)$,
$p^{A_v}\xi_k=p^{D+1}\varpi^{c_*-2(D+1)}\bigl(y\otimes\bigotimes_{u\ne v}b_{u,a_u}\bigr)\in p^{D+1}T_{\sigma_B}(\mathbf L[v\mapsto L'])$.
Step 4 with $E=D+1$ gives $\delta_{i,\mathbf L}(p^{A_v}\xi_k)\in J_B+p^{D+1}M_B$. Writing $\xi=\sum_{\mathbf a}r_{\mathbf a}b_{\mathbf a}$
and splitting the sum according to whether $b_{\mathbf a}\in C_{i,\mathbf L}$ gives $g$ and $c$.

\emph{Step 6: assembly.} Let $C=\bigcup_{(i,\mathbf L)\in\mathcal R}\delta_{i,\mathbf L}(C_{i,\mathbf L})\subseteq M_B$. By Steps 5
and 2, and $\kappa_v\ge0$ (\lref{Formalism.MainNotions.G.LatticeSystems.kappa_nonneg'}),
$|C|\le|\mathcal R|\prod_{v\in B}\kappa_v\prod_v\dim_v\le d^{\mathrm{na}}\prod_{v\in B}\kappa_v$. Let $T=J_B+\operatorname{span}(C)$.
For $x\in M_B$, Step 1 gives $p^Dx=\sum_{(i,\mathbf L)\in\mathcal R}\delta_{i,\mathbf L}(p^Dx(i,\mathbf L))$, and Step 5 applied to
$\xi=x(i,\mathbf L)$ shows $p^Dx\in T+p^{D+1}M_B$. Thus the submodule $N=p^DM_B$ satisfies $N\subseteq T+pN$. Since $M_B\subseteq
M^{\mathrm{na}}$ is finitely generated over the noetherian ring $\OO$, so is $N$, and $p$ lies in the maximal ideal of the local
ring $\OO$; by Nakayama's lemma $N\subseteq T$. Hence $p^{D+1}x=p\cdot p^Dx\in T$ for $x\in M_B$.
\end{proof}

\begin{proposition}[The many-orbits statement]\label{pf:linkG:manyorbits}
Let $\mathcal G$ have trivial stabilisers and finite classes. Let $VV$, $t$, $S_v\le S$, $D$, $c_*$, $m\ge1$ and $\kappa_v$ be
such that $\linkGLP(t_v,D,\kappa_v,c_*)$ holds for $v\in VV$, and put $N^{\mathrm{na}}=N(M^{\mathrm{na}})$,
$\tau=\tau(VV,m,D,t)$, $\Sigma=m\sum_{v\in VV}\kappa_v$. Then \eqref{eq:pf:linkG:manyorbits} holds for every $\beta\ge0$,
$p^{\sum_vS_v}M^{\mathrm{na}}\subseteq N^{\mathrm{na}}$, and if $\Sigma\le\frac12$,
\[
\len\bigl(M^{\mathrm{na}}/(N^{\mathrm{na}}\cap p^{\tau}M^{\mathrm{na}})\bigr)\le\len(\OO/p\OO)\cdot d^{\mathrm{na}}(\tau+2S\Sigma).
\]
\formallabel{Formalism.MainNotions.G.LatticeSystems.counting_statement_of_block_form}
\formallabel{Formalism.MainNotions.G.LatticeSystems.theta_kernel_expansion}
\formallabel{Formalism.MainNotions.G.stability_of_trivial_stabilisers}
\end{proposition}

\begin{proof}
Taking $B=\emptyset$ in the definition of $N^{\mathrm{na}}$ gives $p^{\sum_vS_v}M^{\mathrm{na}}\subseteq N^{\mathrm{na}}$. If
$VV=\emptyset$, then $\tau=0$ and $M^{\mathrm{na}}\subseteq N^{\mathrm{na}}$, so \eqref{eq:pf:linkG:manyorbits} holds with $C=\emptyset$
and the quotient in the length bound is $0$. If $VV\ne\emptyset$, the local slot modules form a slot system (lemma J,
\cref{pf:linkG:lemmaJ}, together with the inclusions recalled before it), the Hecke operators commute,
and $\varpi^2=p$ (Hida setting of any $v\in VV$), so \cref{pf:linkG:walk} gives \eqref{eq:pf:linkG:manyorbits}. For the length bound,
$\kappa_v\ge0$ (\lref{Formalism.MainNotions.G.LatticeSystems.kappa_nonneg'}), $M^{\mathrm{na}}$ is free of finite rank
$d^{\mathrm{na}}$ (\lref{Formalism.MainNotions.G.LatticeSystems.mna_finite} and \lref{Formalism.MainNotions.G.LatticeSystems.mna_free}),
and by \cref{pf:linkG:esymm} the exceptional sets have at most $d^{\mathrm{na}}\Sigma^{\beta+1}$ elements; since also
$p^{|VV|S}M^{\mathrm{na}}\subseteq N^{\mathrm{na}}$, \cref{pf:linkG:lengthallbut} with $a=p$, $V=|VV|$ and $s=\Sigma$ applies.
\end{proof}

The proposition says: up to the span of an exceptional set whose size is controlled by the elementary symmetric functions
of the defects $\kappa_v$, the module $p^{\tau+\beta S}M^{\mathrm{na}}$ is contained in the span of the Hecke products.

\begin{proposition}[Length bound for Hecke products]\label{pf:linkG:core}
Let $\mathcal G$ be a local group datum over $p$ such that: $p^eM^{\mathrm{na}}\subseteq M\subseteq M^{\mathrm{na}}$;
\eqref{eq:pf:linkG:manyorbits} holds for the data below and all $\beta$ (for instance by \cref{pf:linkG:manyorbits}); every
$y_v$ preserves $M^{\mathrm{na}}$ and $M$; and $\mathcal G$ has finite classes. Let
$VV$, $t$, $S_v\le S$, $D$, $c_*$, $m\ge1$ and $\kappa_v$ be such that $\linkGLP(t_v,D,\kappa_v,c_*)$ holds for $v\in VV$,
let $V'$ be a submodule such that $M\cap V'$ is saturated in $M$ (if $x\in M$, $0\ne c\in\OO$ and $cx\in M\cap V'$, then
$x\in M\cap V'$), and assume $\Sigma:=m\sum_{v\in VV}\kappa_v\le\frac12$. Put $N=N(M)$ and $\tau=\tau(VV,m,D,t)$. Then
\begin{enumerate}[label=(\alph*)]
\item $N\cap p^{e+\tau}M\cap V'\subseteq p^{e+\tau}(M\cap V')$;
\item $\len\bigl((M\cap V')/(N\cap p^{e+\tau}M\cap V')\bigr)\le\len(\OO/p\OO)\cdot d\,\bigl(e+\tau+2S\Sigma\bigr)$.
\end{enumerate}
\formallabel{Formalism.MainNotions.FormalGroupSystem.Pf.core}
\formallabel{Formalism.MainNotions.G.LatticeSystems.sat_clause}
\formallabel{Formalism.MainNotions.G.LatticeSystems.inj_clause}
\end{proposition}

\begin{proof}
The local lattice theorem in the form $\linkGLP(t_v,D,\kappa_v,c_*)$ implies $\kappa_v\ge0$
(\lref{Formalism.MainNotions.G.LatticeSystems.kappa_nonneg'}), so $\Sigma\ge0$. The module $M$ is finitely
generated (finite classes and $M\subseteq M^{\mathrm{na}}$) and free, of rank $d\ge d^{\mathrm{na}}$ (by the sandwich).
The modules $N^{\mathrm{na}}=N(M^{\mathrm{na}})$ and $N=N(M)$ satisfy $N^{\mathrm{na}}\subseteq M^{\mathrm{na}}$ and
$N\subseteq M$ (the $y_v$ preserve $M^{\mathrm{na}}$ and $M$), and $p^eN^{\mathrm{na}}\subseteq N$, because
$p^eM^{\mathrm{na}}\subseteq M$ and the generators of $N(\cdot)$ are linear in $F$. By \eqref{eq:pf:linkG:manyorbits} and
\cref{pf:linkG:invshift} (with $A=M^{\mathrm{na}}$, $B=M$),
$\linkGAB_p(M,N,\lfloor d^{\mathrm{na}}m^{\beta+1}e_{\beta+1}(\kappa)\rfloor,(\tau+\beta S)+e)$ for all $\beta$, and by
\cref{pf:linkG:esymm}
\[
\lfloor d^{\mathrm{na}}m^{\beta+1}e_{\beta+1}(\kappa)\rfloor\le\frac{d^{\mathrm{na}}\Sigma^{\beta+1}}{(\beta+1)!}\le d\,\Sigma^{\beta+1}.
\]
Moreover, taking $B=\emptyset$ in the definition of $N(M)$ shows $p^{\sum_vS_v}M\subseteq N$, hence
$p^{|VV|S}M\subseteq N$. \Cref{pf:linkG:lengthallbut} with $a=p$, $e+\tau$ in place of $\tau$, $V=|VV|$ and $s=\Sigma$
gives
\[
\len\bigl(M/(N\cap p^{e+\tau}M)\bigr)\le\len(\OO/p\OO)\cdot d\,(e+\tau+2S\Sigma).
\]
(a) If $x\in N\cap p^{e+\tau}M\cap V'$, write $x=p^{e+\tau}y$ with $y\in M$; saturation gives $y\in M\cap V'$. (b) The
natural map $(M\cap V')/(N\cap p^{e+\tau}M\cap V')\to M/(N\cap p^{e+\tau}M)$ is injective, so (b) follows from the
displayed bound.
\end{proof}

% ---------------------------------------------------------------------------
\subsection{The family argument}\label{sec:pf:linkG:family}

Let $\mathcal F$ be a family. For a member $(D,L)$ with scale $s$ we use the parameters
\[
h_G=2+8\eps,\quad r_G=\lceil8(h_G+2)/\eps\rceil+1,\quad N_1=N_1(r_G,\delta_1,h_G,x),\quad G'_s=G'(r_G,\delta_1,h_G,x),
\]
$c_G=4S_c$ (with $S_c$ a size quantity of $L$), $L_1(p)=\lceil\log_p((kx^A)^4)\rceil$ for an exponent $A$ of the
scale, and $L_2=2\sup\{L_1(p):\ p\text{ in the inert class}\}$. At a good orbit $t_0$ of a candidate prime $p$ the
parameter instance $X$ and the link constants $K$ determine the \emph{local parameters} $t_v(p,t_0)$, $D_p$ and
$\kappa_v(p,t_0)$ of \cref{pf:linkG:orbitdata}: for $p<p_0(k)$ (\emph{small} primes) these are $t_{II}$, $D_{II}$ and
$\kappa_{II}=\min\bigl(1,k_{II}Nr_G^2k\,x^{-A}\bigr)$ with $N=N_v(p,t_0)$, explicit expressions in $K$ and $L$, with
$r=r_G$; for
$p_0(k)\le p<X'$ (\emph{medium} primes) they are $t_I$, $D_I$ and a defect $\kappa_I$ chosen by $X$, with $r=1$.

\begin{definition}[The admissibility conditions]\label{pf:linkG:admcond}
The proof of $\sfname{LinkG}$ for a family uses nine eventual conditions on its members, each a statement about the
orbit data of \cref{pf:linkG:orbitdata}:
\begin{itemize}
\item $\sfname{Adm1F}$: at every prime $p$ in the inert class with $\ord(h_p)\ge x^{\eps/4}$, the good orbits have
$\ord(h_p)\delta_1/(C\log x)\le N_v\le40\ord(h_p)\delta_1$ for a constant $C>0$, and the tops on them satisfy the
spacing conditions of clause 1;
\item $\sfname{Adm2F}$: the condition \eqref{eq:pf:linkG:adm2};
\item $\sfname{Adm3F}$: $\sum\log p/p\le x^{-h_G+\eps}$ over the primes $p<X''(n)$ in the inert class with
$\ord(h_p)<x^{\eps/4}$;
\item $\sfname{Adm4F}$: bounds for the band widths $\rho\le\rho_B\le x^{-2}$ and $\sum\log p/p\le x^{-1}$ over the primes in the
bands near $X'$;
\item $\sfname{Adm5F}$: for every $\delta>0$, eventually $x^{\eps-\delta}\le\Lambda_KN_1^2\log a\le x^{\eps+\delta}$ and
bounds for $\log(p^{G'_s})$ at the small and medium primes;
\item $\sfname{Adm6F}$: at small primes $p\ge13$, for every unramified quadratic setting with $n'$, $TT$ satisfying the
orbit hypothesis with the parameters $(r_G,L_1(p),L_2,\delta_1,h_G,x)$, and every dominant base weight with entries in
the windows of the scale, $\linkGLP\bigl(\lceil t_{II}\rceil,D_{II},\kappa_{II},8(D_{II}+1)\bigr)$ holds with $r=r_G$,
provided a list of explicit size conditions holds;
\item $\sfname{Adm7F}$: the analogous statement at medium primes with $r=1$ and the defect $\kappa_I$;
\item $\sfname{Adm8F}$: $\kappa_I\ge0$, and the mean of $\kappa_I$ over the weight boxes is at most
$C\,|TT|N_1/x^{A/2}$;
\item $\sfname{Adm9F}$: $e_p\le(C+1)fk^2\log_pa$ for $p\le a$, and the sandwich $p^{e_p}M^{\mathrm{na}}\subseteq M\subseteq
M^{\mathrm{na}}$ at every $(u,w)$ with $p=p(w)$.
\end{itemize}
Together with an eventual unit bound at the twisted places, the length condition $\mathsf{Len}$ of
\cref{pf:linkG:admg1} and the size conditions of the type primes, these are the twelve conditions $\mathsf{Adm}_{\mathcal F}$ of the
\emph{Kudla family statement}.
\formal{Formalism.MainNotions.AuxBlock8.admitted}
\formal{Formalism.MainNotions.G.Adm6F}
\end{definition}

\begin{theorem}[Kudla family statement]\label{pf:linkG:stmtF}
Let $\mathcal F$ be a family, $K$ non-negative link constants, and $X,Y$ parameter and lattice instances for its
members. Suppose that there are $C_{31}>0$, $C_e,C_m\ge0$ such that for every $\delta>0$, eventually, every member
satisfies the twelve conditions $\mathsf{Adm}_{\mathcal F}$ with the constants $(C_{31},C_e,C_m,\delta)$. Then there are $C$ and a
function $\delta(x)\to0$ such that, eventually, every member satisfies $\mathsf{Gain}(D,L,X,K,C,b^{\delta(x)})$:
\begin{enumerate}[label=(\alph*)]
\item at every prime $p\le X''(n)$ in the inert or signal class: $0\le\mathrm{gain}_P(p,j)\le4|TT|a/p$ for $1\le j\le
m$, and
\begin{equation}\label{eq:pf:linkG:conc1}
\sum_{j=1}^m\sum_{s\in TT_P(p)}\bigl(\Lambda^{G}_{j,s}(p)-\mathrm{lo}(p)\bigr)-m\,\mathrm{fl}(p)\le
\sum_{j=1}^m\mathrm{gain}_P(p,j)-\tau(p),
\end{equation}
where $TT_P(p)$ is the set of tops on the served orbits (all of $TT$ if $p$ is in the signal class or $p\ge X'$),
$\Lambda^G_{j,s}(p)=\sum_{e\text{ odd}}\max\bigl(0,a_T(s)/p^e-\phi_{j}(s,e)\bigr)$ for explicit filling values
$\phi_j(s,e)$, and $\mathrm{lo}(p)=\mathrm{lo}(K,L,C,b^{\delta(x)},p)$, $\mathrm{fl}(p)$ are explicit error terms;
\item the mean-$\kappa$ bound of \cref{pf:linkG:linkGat}(iv);
\item a points lemma, the existence of weight vectors, a comparison $\Lambda^P\le\Lambda^G\le\Lambda^P+N(m)\log_p(2a)$
of the class sums for the fillings with and without the shift $(k+(2j-1)(t+1))/2$, and a weighted prime sum bound.
\end{enumerate}
\formallabel{Formalism.MainNotions.AuxBlock8.Pf.stmtF}
\formallabel{Formalism.MainNotions.AuxBlock8.Pf.stmtF_of_stmtEF}
\end{theorem}

\begin{proof}
By \cref{pf:linkG:stmtEF} (the closed form below), for non-negative link constants $K$ and the constants
$C_{31}>0$, $C_e,C_m\ge0$ of the hypothesis there is $C$ such that, for all instances $X,Y$ whose members eventually satisfy
$\mathsf{Adm}_{\mathcal F}$ with the constants $(C_{31},C_e,C_m,\delta)$ for every $\delta>0$, the conclusion holds with this $C$
and some function $\delta(x)\to0$. This is the theorem.
\end{proof}

This theorem is the analytic heart of the gains: the inequality \eqref{eq:pf:linkG:conc1} says that the divisibility
gains produced by the served orbits at $p$ dominate the von Mangoldt class sums at $p$. Its closed form is proved from
six component statements, which estimate sums over the primes of the inert classes by the orbit counts, the gain bounds at
every candidate prime and lattice chains of the served orbits.

\begin{theorem}[Closed form of the Kudla family statement]\label{pf:linkG:stmtEF}
Let $\mathcal F$ be a family with parameter $\eps$. For all non-negative link constants $K$ and all $C_{31}>0$,
$C_e,C_m\ge0$ there is $C\ge0$ such that for all parameter and lattice instances $X,Y$: if for every $\delta>0$, eventually,
every member satisfies $\mathsf{Adm}_{\mathcal F}$ with the constants $(C_{31},C_e,C_m,\delta)$, then, with
$\delta(x)=x^{-\eps/2}\to0$, eventually every member satisfies $\mathsf{Gain}(D,L,X,K,C,b^{\delta(x)})$.
\formallabel{Formalism.MainNotions.AuxBlock8.Pf.stmtEF}
\formallabel{Formalism.MainNotions.AuxBlock8.Pf.stmtEF_of}
\formallabel{Formalism.MainNotions.AuxBlock8.Pf.remainingCountingStatement_of_restStar}
\end{theorem}

\begin{proof}
Fix $K$, $C_{31},C_e,C_m$ and write $\beta(x)=b^{x^{-\eps/2}}$. For a member, a prime $p$, a lift $j$ and tops $s$, let
$\mathrm{ask}_P(p,j,s)$ and $\mathrm{ask}_G(p,j,s)$ be the class sums of the fillings without and with the shift (so
$\sum_s\mathrm{ask}_G$ is the left side of \eqref{eq:pf:linkG:conc1} before subtracting the error terms), and $(a,b)$ range over
the pairs of top data in the window of the scale ($a\le a_t\le a+a/(kx^{A/2})$ and $b'\le b_t$ on $TT$). Six statements are
combined, each valid eventually in $x$ under the hypothesis.
\begin{enumerate}[label=(\arabic*)]
\item \emph{Star bounds} (\cref{pf:linkG:starbounds}): there is $C_0\ge0$ such that for $C\ge C_0$,
at every candidate prime $p$, every good representative $t_0$ and every lift $j$,
\[
\sum_{s\text{ on the orbit of }t_0}\bigl(\mathrm{ask}_P(p,j,s)-\mathrm{lo}(p)\bigr)\le g_U(t_0)+\mathrm{dl}(t_0)-\Pr(t_0)-\lceil
t_v(t_0)\rceil-D_p-1,
\]
where $g_U(t_0)$ is the gain of the orbit and the subtracted terms are its costs; this is proved orbit by orbit from the
local lattice data of the family (the conditions $\sfname{Adm1F}$--$\sfname{Adm9F}$).
\item \emph{Comparison of fillings} (\cref{pf:linkG:asks}): for $a_t\le2a$ and $b_t\ge0$,
$\mathrm{ask}_P\le\mathrm{ask}_G\le\mathrm{ask}_P+N(m)\log_p(2a)$.
\item \emph{Gain bounds}: let $C_1=\max(C_0,3,2\mathrm{dl}_1+C_e+1,4\mathrm{dl}_2)$ and $C=C_1+1$. First, the version of the gain
bound with $\mathrm{ask}_P$ in place of $\mathrm{ask}_G$ holds with the constant $C_1$
(\cref{pf:linkG:gainstar}: at a candidate prime the left side is bounded by summing
the star bounds (1) over the served orbits; the remaining primes are handled with the unit bound of $\sfname{Adm9F}$). Then
(2) gives, at every prime $p\le X''(n)$ in the inert or signal class, $0\le\mathrm{gain}_P(p,j)\le4|TT|a/p$ and
\eqref{eq:pf:linkG:conc1} with the constant $C$
(\cref{pf:linkG:addone}: the extra term from (2) is
absorbed by increasing the constant by $1$).
\item \emph{Mean $\kappa$} (\cref{pf:linkG:meanw}): the mean-$\kappa$ bound of (b), from
$\sfname{Adm8F}$.
\item \emph{Points and weight vectors}
(\cref{pf:linkG:points,pf:linkG:meanw}): given the gain bounds,
the points lemma holds (from a lattice chain statement for the served orbits and numerical facts on prime sums along the
family), and the weight vectors of (c) exist.
\item \emph{Weighted prime sums} (\cref{pf:linkG:points}): the weighted prime sum bound of (c),
for every $C\ge0$, from the third estimate along the family.
\end{enumerate}
Now let $X,Y$ satisfy the hypothesis and let $x_0$ be the maximum of the thresholds of (1)--(6). For $x\ge x_0$ and every
member, (3) gives part (a) of $\mathsf{Gain}$ for the top data of the member (which lie in the window by the definition of
size data), (4) gives (b), and (5), (2), (6) give the four statements of (c). Finally $\delta(x)=x^{-\eps/2}\to0$ as $\eps>0$.
\end{proof}

\begin{lemma}[Specialisation of the local lattice theorem]\label{pf:linkG:adm6}
Let $C\ge0$ be at most each of the link constants $k_{II},t_{II},\Lambda_K,D_{II},s_{II},k'_{II}$ of $\mathcal F$, and
suppose that, eventually, $\linkGLL(C)$ holds for the members. Then $\sfname{Adm6F}(\mathcal F)$ holds.
\formallabel{Formalism.MainNotions.G.family_second_bound_of_theorem}
\formallabel{Formalism.MainNotions.G.Joints.type_two_hypothesis_of_counts}
\formallabel{Formalism.MainNotions.G.Joints.LPprime_min_one}
\formallabel{Formalism.MainNotions.G.Joints.LPprime_mono_kappa}
\end{lemma}

\begin{proof}
Fix a member and a small prime $p$ with the data and hypotheses of $\sfname{Adm6F}$. Apply $\linkGLL(C)$ with
$m+1=k$, $r=r_G$, $(L_1,L_2,\delta,h,x)=(L_1(p),L_2,\delta_1,h_G,x)$, $a-c_G$ in place of $a$, $Y=x^A$,
$c_{pp}=\lfloor c_G\rfloor$, $c_0=\lfloor2b'+2c_G\rfloor$, $D_p=D_{II}$ and $c_*=8(D_{II}+1)$. The hypotheses
\ref{pf:linkG:H1}--\ref{pf:linkG:H5} follow from those of $\sfname{Adm6F}$: each is an inequality between explicit
quantities, and the inequalities $C\le k_{II},\dots,k'_{II}$ allow one to replace $C$ by the corresponding link constant;
the facts used are $c_G\ge0$, $a-c_G\ge1$ (from the size data), $\log(a-c_G)\le\log a$, $\lfloor c_G\rfloor\le c_G$,
and $k^2\le N_1^3+r_G^5+k^2$. In the same way $t_{\mathrm B}\le\lceil t_{II}\rceil$, so $\linkGLL(C)$ gives
$\linkGLP(\lceil t_{II}\rceil,D_{II},C|TT|r_G^2k/x^A,c_*)$. Since $\min(1,C|TT|r_G^2k\,x^{-A})\le\kappa_{II}$ (as $C\le
k_{II}$, with $\kappa_{II}$ computed with $N=|TT|$), the monotonicity of $\linkGLP$ in $\kappa$ and its invariance under
$\kappa\mapsto\min(1,\kappa)$ give the claim.
\end{proof}

\begin{lemma}[Closedness of the local input]\label{pf:linkG:secondclosed}
If there is $C_0\ge0$ such that $\linkGLL(C_0)$ holds for every datum, then on every family class the property
$\mathsf T_{II}$ is closed: there is $C$ with $0\le C\le k_{II},t_{II},\Lambda_K,D_{II},s_{II},k'_{II}$ such
that, eventually, $\linkGLL(C)$ holds for the members.
\formallabel{Formalism.MainNotions.G.second_theorem_closed_form}
\end{lemma}

\begin{proof}
Take $\eps_1=1$ and, for every $\eps$, the staged threshold with $K_{\mathrm{mid}}\equiv K_{\mathrm{last}}\equiv C_0$. If
$\mathcal F$ is large, then the middle constant $\Lambda_K$ and the last constants $k_{II},t_{II},D_{II},s_{II},k'_{II}$
are at least $C_0$; take $C=C_0$.
\end{proof}

\begin{lemma}[Further closed properties]\label{pf:linkG:otherclosed}
On every family class the following properties are closed: $1\le\mathrm{Pr}_1$, $1\le C_N$ and $1\le C_e$; and,
eventually, $N_1\le C_N\log x$.
\formallabel{Formalism.MainNotions.G.ones_closed}
\formallabel{Formalism.MainNotions.G.cb_closed}
\end{lemma}

\begin{proof}
By the definition of a family, $\mathrm{Pr}_1=1$ and $C_e=1$; $C_N$ is a last constant, so the staged threshold with
$K_{\mathrm{last}}\equiv1$ gives $C_N\ge1$. For the second property take $\eps_1=1/104$. If a datum $D_0$ exists, the orbit
counts (\lref{Formalism.MainNotions.G.TopChoiceBounds.master_orbit_sizes_hold}, applied with $h=2+8\eps$ and
$r=\lceil8(h+2)/\eps\rceil+1$) give $C_e'$ with $N_1(r,\delta,h,x)\le C_e'\log x$ for all $x\ge3$ and $0<\delta\le1$; the
staged threshold with $K_{\mathrm{last}}\equiv C'_e$ forces $C_N\ge C'_e$, and $0<\delta_1\le1$, so $N_1\le C_N\log x$ for
$x\ge\max(x_1,3)$ (here $h_G=2+8\eps$ and $r_G$ are the parameters of the member). A family has a member at $x=0$, so if
there is no datum there is no family and the property holds vacuously.
\end{proof}

In the same way the following properties are closed on every class: $\Lambda_K>0$ and $C_e\ge0$ (the constant
$\Lambda_K$ is a middle constant and $C_e=1$ by the definition of a family;
\lref{Formalism.MainNotions.G.signs_closed}); a window logarithm bound
(\lref{Formalism.MainNotions.G.TopChoiceBounds.family_log_bound_closed}); and the \emph{first theorem}
$\mathsf T_{I}$ (\lref{Formalism.MainNotions.G.first_theorem_closed_form}), a bound for the mean of the
medium-prime defects $\kappa_I$, obtained from threshold functions
(\lref{Formalism.MainNotions.G.first_theorem_threshold_functions}) whose hypotheses are supplied, under
$\sfname{Loc}$, by $(\mathrm E)$, $(\mathrm L6)$ and a package of local flatness results
(\lref{Formalism.L2.G.local_flatness_results} and \lref{Formalism.MainNotions.G.loc8All_of}).

\begin{definition}[Lengths at the served orbits]\label{pf:linkG:admg1}
For a member, a place $w$ with $p=p(w)$ and a word $u$, let $\mathcal G=\mathcal G(u,w)$, let $V_{\mathrm{ram}}$ be the
ramified subspace of the space of functions of $\mathcal G$ chosen by the lattice instance, $M_{\mathrm{ram}}=M\cap
V_{\mathrm{ram}}$, and let $N_P\subseteq M_{\mathrm{ram}}$ be the intersection of the span of the Hecke products
$\bigl(\prod_{t_0\in V'}y_{v(t_0)}^m\bigr)\bigl(p^{\sum_{t_0\in\linkGVV(p)\smallsetminus V'}\lceil S_v(p,t_0)\rceil}F\bigr)$
($V'\subseteq\linkGVV(p)$, $F\in M$) with $p^{\lceil\tau(p)\rceil}M$ and $V_{\mathrm{ram}}$; here $v(t_0)$ is the place of
$\mathcal G$ attached to the orbit $t_0$. The condition $\mathsf{Len}$ asserts, for all $(u,w)$: $N_P\subseteq
p^{\lceil\tau(p)\rceil}M_{\mathrm{ram}}$; $\len(M_{\mathrm{ram}}/N_P)\le v_\OO(p)\,d\,(\tau(p)+\iota(p))$; and a valuation
bound $\ge\sum_j\mathrm{gain}_P(p,j)+\mathrm{low}(u,w)$ for the elements of $N_P$ under the chain map of $Y$.
\formal{Formalism.MainNotions.AuxBlock8.admG1}
\end{definition}

\begin{proposition}[Admissibility without kernels]\label{pf:linkG:afnk}
Let $\mathcal F$ be a family satisfying $\sfname{Adm1F}$--$\sfname{Adm9F}$, and suppose that eventually: the local group
data of the members have trivial stabilisers and finite classes; the Hecke operators preserve $M^{\mathrm{na}}$ and $M$;
the \emph{orbit-lattice bounds} hold ($\linkGLP(\lceil t_v\rceil,D_p,\kappa_v,8(D_p+1))$ at the place $v(t_0)$ of every
served orbit $t_0$); the lattice sandwich holds with the exponents $e_p$; and $\mathsf{Desc}(\mathcal F)$ holds. Let
$Y_0$ be the lattice instance $Y$ with its chain maps replaced by $0$. Then the conditions $\mathsf{Adm}_{\mathcal F}$ hold for
$(\mathcal F,G_K,G_x,Y_0)$, and eventually the unit bound and $\mathsf{Len}$ hold for $Y_0$.
\formallabel{Formalism.L2.G.admissible_family_no_kernel}
\formallabel{Formalism.MainNotions.FormalGroupSystem.Pf.stmt6G}
\formallabel{Formalism.MainNotions.FormalGroupSystem.Pf.coreT}
\formallabel{Formalism.L2.G.vKge_zero}
\end{proposition}

\begin{proof}
For $Y_0$ the chain maps are $0$, so every valuation statement ``$\ge g$'' for an image under the chain map holds
trivially ($0$ lies in every lattice); this gives the unit bound at the twisted places and the third part of
$\mathsf{Len}$. Fix a member (eventually) and $(u,w)$. By trivial stabilisers and finite classes, the many-orbits
statement holds for $\mathcal G(u,w)$ (\cref{pf:linkG:manyorbits}). If
$\linkGVV(p)=\emptyset$, then $\tau(p)=\iota(p)=0$ and $N_P=M_{\mathrm{ram}}$ (take $V'=\emptyset$), so the first two parts
of $\mathsf{Len}$ hold. Otherwise apply \cref{pf:linkG:core} with the served orbits as $VV$ (through the place map
$t_0\mapsto v(t_0)$, injective on representatives), $t_v=\lceil t_v(p,t_0)\rceil$, $D=D_p$, $c_*=8(D_p+1)$, $\kappa=\kappa_v$,
$S_v=\lceil S_v(p,t_0)\rceil$, $e=e_p$ and $V'=V_{\mathrm{ram}}$: the hypotheses are the sandwich, the orbit-lattice
bounds, the preservation of the lattices, finite classes, the saturation of $M_{\mathrm{ram}}$ in $M$
(\lref{Formalism.MainNotions.G.saturated_ramified_lattice}, from the comparison data of $\mathrm{ex}(x,i)$) and
$\Sigma(p)\le\frac12$
(eventually, since $\Sigma(p)\le x^{-\eps}$ when $\linkGVV(p)\ne\emptyset$).
Since $\tau(p)=e_p+\tau(VV,m,D_p,t)$ and $\iota(p)=2\Sigma(p)\max S_v$, \cref{pf:linkG:core} gives the first two parts of
$\mathsf{Len}$ (\lref{Formalism.MainNotions.G.TopChoiceBounds.gluing_statement_holds}). The remaining conditions of
$\mathsf{Adm}_{\mathcal F}$ are the hypotheses $\sfname{Adm1F}$--$\sfname{Adm9F}$ and the size conditions of the type primes, which
hold eventually for every family (\lref{Formalism.MainNotions.G.eventually_sizes_admitted}); the constants
$C_{31},C_e,C_m$ are those of $\sfname{Adm1F},\sfname{Adm9F},\sfname{Adm8F}$.
\end{proof}

\paragraph{Fill identities.}
A family carries, as part of its definition (the record $\sfname{FamP}$), the following identities and inclusions for
every member with $x>x_1$ and every place $w$: with $p=p(w)$,
\begin{gather}
\tau(w)=\tau(p),\qquad\iota(w)=\iota(p),\qquad TT_w=TT_P(p),\qquad \Lambda_K=\Lambda_G(L),\notag\\
\mathrm{gain}(j,w)=\mathrm{gain}_P(p,j)\quad\text{if }w\in\supp L\text{ and }p\text{ is inert or signal},
\label{eq:pf:linkG:fill}
\end{gather}
the class sums $\Lambda_{j,t}(p)$ of $L$ equal the sums $\Lambda^P_{j,t}(p)$ for the filling without shift, and every
place with $p(w)\le X''(n)$ lies in the support of $L$. Moreover the lattice $N_u$ and the lengths
$\mathrm{len}_{MN}(u,w)$ are tied to the local data: $N_u\subseteq M^{\mathrm{ram}}_u$ locally at every $w$ with $p(w)\ne
T_u$ (part of the definition of $N_u$), and, under $\mathsf{Desc}$, eventually
$\len(M_{\mathrm{ram}}/N_P)=\mathrm{len}_{MN}(u,w)\,v_\OO(p)$ for $w\in\supp L$, and $\mathrm{len}_{MN}(u,w)=0$ if
$\linkGVV(p(w))=\emptyset$ and $p(w)\ne T_u$ (\lref{Formalism.MainNotions.G.sizeDataLat_ev}). Finally, eventually: the served orbits
give admissible indices for every word, $S_v(p,t_0)\le2mN_v(p,t_0)a/p$, and $N_u\subseteq\mathfrak p_w^{\lfloor
e_w\tau(w)\rfloor}M^{\mathrm{ram}}_u$ locally at $w$ when $\linkGVV(p(w))\ne\emptyset$ and $p(w)\ne T_u$
(\lref{Formalism.MainNotions.G.servedGIdx_ev}, \lref{Formalism.MainNotions.G.svBound_ev}, and \lref{Formalism.MainNotions.G.nuServed_ev}).

\begin{proposition}[The link at one member]\label{pf:linkG:linkgatcore}
Let $(D,L)$ be a member with $x>x_1$, satisfying the fill identities, the lattice-length identities and the eventual
properties of the previous paragraph, and assume parts (a), (b) and the comparison of class sums of
$\mathsf{Gain}(D,L,X,K,C,b_o)$, and $\mathsf{Len}$ for a lattice instance. Then
$\sfname{LinkG}_{\mathrm{at}}(X,K,4|C|+1,b_o)$ holds.
\formallabel{Formalism.MainNotions.G.linkG_at_core}
\formallabel{Formalism.MainNotions.G.linkG_at_of}
\end{proposition}

\begin{proof}
(i) Let $u,w$ with $p=p(w)\ne T_u$. If $\linkGVV(p)=\emptyset$, then $\tau(w)=\tau(p)=0$ and $\iota(w)=\iota(p)=0$ by
\eqref{eq:pf:linkG:fill}, $\mathrm{len}_{MN}(u,w)=0$, and the local containment at level $\lfloor e_w\cdot0\rfloor=0$ is
part of the definition of $N_u$; all three assertions follow. If $\linkGVV(p)\ne\emptyset$, the local containment is the last eventual
property. Since $p$ is then a candidate, $p<X'\le X''(n)$, so $w\in\supp L$; the served orbits give an admissible index
for $(u,w)$, and $\mathsf{Len}$ gives an integer $l=\len(M_{\mathrm{ram}}/N_P)\le
v_\OO(p)\,d\,(\tau(p)+\iota(p))$, while $l=\mathrm{len}_{MN}(u,w)\,v_\OO(p)$. As $d=d_u$ and $v_\OO(p)>0$, dividing by
$v_\OO(p)$ gives $\mathrm{len}_{MN}(u,w)\le d_u(\tau(w)+\iota(w))$.

(ii) is the bound $S_v\le2mN_va/p$.

(iii) Let $w$ with $p=p(w)$ in the inert or signal class and $p\le X''(n)$. Then $w\in\supp L$, so
$\mathrm{gain}(j,w)=\mathrm{gain}_P(p,j)$, and part (a) of $\mathsf{Gain}$ gives $0\le\mathrm{gain}(j,w)\le4|TT|a/p$
and \eqref{eq:pf:linkG:conc1}. Using $TT_w=TT_P(p)$, $\tau(w)=\tau(p)$ and the identity of class sums, the left side of
\eqref{eq:pf:linkG:vonMangoldt} (with the constant $4|C|+1$) is
$\sum_j\sum_{s\in TT_P(p)}(\Lambda^P_{j,s}(p)-\mathrm{lo}_G(p))-m\,\mathrm{fl}_G(p)$. Termwise,
$\Lambda^P_{j,s}(p)\le\Lambda^G_{j,s}(p)$ by the comparison of class sums (whose hypotheses $a_T(s)\le2a$ and $b_T(s)\ge0$
follow from the windows of the size data), $\mathrm{lo}(p)\le\mathrm{lo}_G(p)$
(\lref{Formalism.MainNotions.G.lo_le_loG}, using $\Lambda_K=\Lambda_G(L)$ and $K\ge0$), and $\mathrm{fl}(p)\le
\mathrm{fl}_G(p)$ (\lref{Formalism.MainNotions.G.fl_le_flG}). Summing and using $m\ge0$ gives
\eqref{eq:pf:linkG:vonMangoldt}.

(iv) The containment of served orbits holds by the definition of $\linkGVV$, and the mean-$\kappa$ bound is part (b) of
$\mathsf{Gain}$.
\end{proof}

\begin{theorem}[The link without kernels]\label{pf:linkG:withoutkernels}
Let $\mathcal F$ satisfy the hypotheses of \cref{pf:linkG:afnk}. Then $\sfname{LinkG}(\mathcal F)$ holds.
\formallabel{Formalism.L2.G.linkG_without_kernels}
\end{theorem}

\begin{proof}
By \cref{pf:linkG:afnk} the hypotheses of \cref{pf:linkG:stmtF} hold for $(\mathcal F,G_K,G_x,Y_0)$, which gives $C$
and $\delta(x)\to0$ with $\mathsf{Gain}(D,L,G_x,G_K,C,b^{\delta(x)})$ eventually. Also eventually $\mathsf{Len}$
holds for $Y_0$, and the eventual properties of the fill paragraph hold. Take $x_0$ beyond all these thresholds and
$x_1$. For $x\ge x_0$ and $i\in I_x$, \cref{pf:linkG:linkgatcore} (with the lattice instance $Y_0$, which has the
same place maps and ramified subspaces as $Y$) gives $\sfname{LinkG}_{\mathrm{at}}(G_x,G_K,4|C|+1,b^{\delta(x)})$. Hence
$\sfname{LinkG}(\mathcal F)$ with the constant $4|C|+1$, the threshold $x_0$ and the function $\delta$.
\end{proof}

\begin{theorem}[The link for a family]\label{pf:linkG:fam}
Let $\mathcal F$ be a family with $\mathsf T_{II}$, $\mathsf T_{I}$, $\Lambda_K>0$,
$C_e\ge1$, eventually $N_1\le C_N\log x$, eventually trivial stabilisers and finite classes, eventually the
standard-subgroup property $\mathsf{Std}$, $\sfname{Adm2F}$, $\mathsf{Desc}$ and the window logarithm bound.
Then $\sfname{LinkG}(\mathcal F)$ holds.
\formallabel{Formalism.L2.G.linkG_fam_noKer}
\end{theorem}

\begin{proof}
We verify the hypotheses of \cref{pf:linkG:withoutkernels}. $\sfname{Adm6F}$ is \cref{pf:linkG:adm6}. $\sfname{Adm7F}$
holds because $\kappa_I$ is defined by the parameter instance as the infimum of the set of defects for which the
medium-prime statement holds, and this infimum belongs to the set (\lref{Formalism.MainNotions.G.Joints.kappa_share_of_least_serving}).
$\sfname{Adm8F}$ follows from the first theorem (\lref{Formalism.MainNotions.G.Joints.type_one_hypothesis_of_bounds}).
$\sfname{Adm1F}$, $\sfname{Adm4F}$ and $\sfname{Adm5F}$ hold for every family with the stated constants
(\lref{Formalism.MainNotions.G.family_first_lemma}, \lref{Formalism.MainNotions.G.family_orders_bound_of_lemma}, and \lref{Formalism.MainNotions.G.family_sizes_bound}),
and $\sfname{Adm3F}$ follows from an order estimate for the Frobenius elements
(\lref{Formalism.MainNotions.G.family_orders_first_part}, with
\lref{Formalism.MainNotions.G.TopChoiceBounds.family_order_bound_holds} and the field lemmas
\lref{Formalism.MainNotions.UnusedSlot33.field_lemmas_except_seven}). $\sfname{Adm9F}$ (with the constant $C_e$) is the
Weyl sandwich: $p^{e_p}M^{\mathrm{na}}\subseteq M\subseteq M^{\mathrm{na}}$ with $e_p\le(C_e+1)fk^2\log_pa$, from the
comparison of Weyl and dual Weyl lattices in the model and the definition of $e_p$
(\lref{Formalism.MainNotions.G.Joints.weyl_sandwich_of_factors}, \lref{Formalism.MainNotions.G.Joints.weylSandwich_of_model}, and \lref{Formalism.MainNotions.G.Joints.expBound_of}),
and it gives the lattice sandwich (\lref{Formalism.MainNotions.G.TopChoiceBounds.sandwich_of_exponent_clause}). The Hecke
operators preserve the lattices (\lref{Formalism.L2.G.Model.operators_preserve_lattices}). The orbit-lattice bounds
follow from $\sfname{Adm6F}$ at small primes, $\sfname{Adm7F}$ at medium primes, $\mathsf{Std}$ (which presents the
rings of every place as a base change of an unramified quadratic setting), the base change property of $\linkGLP$
(\lref{Formalism.L2.G.Model.lpBaseChange}), two size hypotheses of the members
(\lref{Formalism.L2.G.hypIAt_ev} and \lref{Formalism.MainNotions.G.Joints.hypIIAt_fam}) and $\rho\le\rho_B$ (from
$N_1\le C_N\log x$); this is \lref{Formalism.MainNotions.G.Joints.linear_programming_input_of_hypotheses}.
\end{proof}

% ---------------------------------------------------------------------------
\subsection{Proof of \texorpdfstring{\cref{thm:linkG}}{the local lattice link}}\label{sec:pf:linkG:proof}

\begin{lemma}[Consequences of group compatibility]\label{pf:linkG:groupcompat}
Let $\mathcal F$ be a family which satisfies the group compatibility condition
$\mathsf{GC}$. Then eventually the local group data
of the members have trivial stabilisers and finite classes, eventually the members satisfy $\mathsf{Std}$, and
$\mathsf{Desc}(\mathcal F)$ holds.
\formallabel{Formalism.L2.G.gatEF_of_tied}
\formallabel{Formalism.L2.G.stdSubF_of_tied}
\formallabel{Formalism.L2.G.npDescF_of_tied}
\end{lemma}

\begin{proof}
Take $x$ beyond the threshold of the group compatibility and beyond $x_1$, and a group compatibility record for the
member. Its hermitian level datum has trivial stabilisers at every word, because the gram matrix of the hermitian
space is totally definite and the two type primes are distinct primes
(\lref{Formalism.L2.G.Triv.trivStabU}); together with the description of the local group data as data attached
to the definite unitary group, this gives trivial stabilisers and finite classes
(\lref{Formalism.L2.G.gatE_of_gatIs}). The ring isomorphisms of the record give $\mathsf{Std}$
(\lref{Formalism.L2.G.stdSubAt_of_ring}). The comparison record gives a valuation-rationality property of the
comparison maps (\lref{Formalism.L2.G.Cmp.CmpTied.valRat}), from which $\mathsf{Desc}$ follows
(\lref{Formalism.L2.G.Cmp.npDescF_of_valRatF}).
\end{proof}

\begin{theorem}[Ten closed properties]\label{pf:linkG:tenclosed}
Assume $\sfname{Loc}$. Let $\mathcal C$ be a family class on which the following are closed: eventually trivial
stabilisers and finite classes; eventually $\mathsf{Std}$; $\sfname{Adm2F}$; $\mathsf{Desc}$. Then
$\sfname{LinkG}$ is closed on $\mathcal C$.
\formallabel{Formalism.L2.G.link_of_local_closed}
\end{theorem}

\begin{proof}
By \cref{pf:linkG:locmodel}, $(\mathrm E)$ and $(\mathrm L6)$ hold under $\sfname{Loc}$, so \cref{pf:linkG:localthm}
gives $C_0\ge0$ with $\linkGLL(C_0)$ for every datum, and \cref{pf:linkG:secondclosed} shows that
$\mathsf T_{II}$ is closed on $\mathcal C$. The first theorem, $\Lambda_K>0\wedge C_e\ge0$,
$1\le\mathrm{Pr}_1\wedge1\le C_N\wedge1\le C_e$, eventually $N_1\le C_N\log x$, and the window logarithm bound are
closed on every class (\cref{pf:linkG:otherclosed} and the paragraph after it). The conjunction of these ten closed properties
is closed, and for every family with $\eps>0$ it implies the hypotheses of \cref{pf:linkG:fam}; hence $\sfname{LinkG}$ is
closed on $\mathcal C$.
\end{proof}

\begin{theorem}[\cref{thm:linkG}]\label{pf:linkG:main}
Let $\mathcal C$ be a family class on which group compatibility $\mathsf{GC}$ and $\sfname{Adm2F}$ are closed. Then
$\sfname{LinkG}$ is closed on $\mathcal C$.
\formallabel{Formalism.L2.G.linkG_concrete5}
\formallabel{Formalism.L2.G.linkG_concrete_noKer}
\end{theorem}

\begin{proof}
By \cref{pf:linkG:groupcompat}, the closedness of $\mathsf{GC}$ implies the closedness of eventual
trivial stabilisers and finite classes, of eventual $\mathsf{Std}$ and of $\mathsf{Desc}$ (closedness is preserved
under implications). $\sfname{Loc}$ holds by \cref{pf:linkG:loc}. \Cref{pf:linkG:tenclosed} gives the claim.
\end{proof}


% ---------------------------------------------------------------------------
\subsection{Counting good paths}\label{sec:pf:linkG:goodpaths}

We prove the good path weight bound used in Step 4 of the proof of \cref{pf:linkG:secondpart}
(\cref{pf:linkG:gpw} below). It is a lower bound for the number of pairs consisting of a \emph{good} path of weights and a
standard family for its final weight. The total number of pairs is a product of binomial coefficients, by the rules of
Pieri and of branching; the good paths include the \emph{typical} ones; and the atypical paths carry an exponentially small
part of the total, by a comparison of the logarithm of Weyl's dimension formula with sums of logarithms of binomial
coefficients. Throughout, a weight $\lambda\in\Z^n$ is \emph{dominant} if $\lambda(1)\ge\dots\ge\lambda(n)$, and for weights
$\lambda,\nu\in\Z^n$, $\nu/\lambda$ is a \emph{horizontal strip of size $d$} if $\lambda(m)\le\nu(m)$ for all $m$,
$\nu(m+1)\le\lambda(m)$ for $m<n$, and $\sum_m(\nu(m)-\lambda(m))=d$. For a dominant $\tau\in\Z^n$ let
\[
w_n(\tau)=\prod_{1\le i<j\le n}\frac{\tau(i)-\tau(j)+j-i}{j-i}
\]
be Weyl's dimension of the irreducible representation of $\GL_n$ with highest weight $\tau$. We write
$\lambda^+=\sum_m\max(\lambda(m),0)$ and $\lambda^-=\sum_m\max(-\lambda(m),0)$.

\paragraph{The three classical rules.} We use three facts on Weyl dimensions \cite{FultonHarris1991,GelfandTsetlin1950},
which are components of the second-part list (the Weyl dimension formula, branching rule and Pieri rule in counting form,
\lref{Formalism.MainNotions.G.Joints.Lit.weyl_dimension_formula_index_sets}, \lref{Formalism.MainNotions.G.Joints.Lit.branching_rule_count_form}, and \lref{Formalism.MainNotions.G.Joints.Lit.pieri_rule_count_form}):
\begin{enumerate}[label=(\roman*)]
\item \emph{Standard families.} For a dominant $\tau\in\Z^n$, the \emph{standard families} of shape $\tau$ (families
$(I_c)_{c\in\Z}$ of subsets of $\{1,\dots,n\}$ with $I_c=\emptyset$ outside the column range of $\tau$,
$|I_c|=\#\{m:\tau(m)\ge c\}$ for $c\ge1$ and $|I_c|=n-\#\{m:\tau(m)\ge c\}$ for $c\le0$, and the column condition: for all $c$
and $t$, the column $c+1$ has at most as many entries $\le t$ as the column $c$) are $w_n(\tau)$ in number; $w_n(0)=1$.
\item \emph{Pieri.} For dominant $\lambda\in\Z^n$ and $B\in\N$,
$\sum_{\nu}w_n(\nu)=w_n(\lambda)\binom{B+n-1}{n-1}$, the sum over the $\nu$ with $\nu/\lambda$ a horizontal strip of size $B$,
and also over the $\nu$ with $\lambda/\nu$ a horizontal strip of size $B$.
\item \emph{Branching.} For dominant $\nu\in\Z^{n+1}$, $\sum_\tau w_n(\tau)=w_{n+1}(\nu)$, the sum over the $\tau\in\Z^n$
interlacing $\nu$ ($\nu(l+1)\le\tau(l)\le\nu(l)$ for $l\le n$).
\end{enumerate}

\begin{definition}[Letter systems, paths and good paths]\label{pf:linkG:paths}
Let $N,n_0,r,K\in\N$ with $N=n_0+r$, $r,K\ge1$, and $k=r+K$. A \emph{letter system} assigns to every $i\in\{1,\dots,N\}$ a type,
\emph{vector} or \emph{covector}, and a size $B_i\in\N$; put $\beta=\sum_{i\le n_0}B_i$ and
\[
\rho=\prod_{i=1}^{n_0}\binom{B_i+k-1}{k-1}.
\]
A \emph{path} is a sequence $\mu=(\mu_0,\dots,\mu_N)$ of dominant weights in $\Z^k$ with $\mu_0=0$ such that for every
$i\le N$: $\mu_i/\mu_{i-1}$ is a horizontal strip if $i$ is a vector letter, and $\mu_{i-1}/\mu_i$ is a horizontal strip if
$i$ is a covector letter; the strip has size $B_i$ if $i\le n_0$; and for $n_0<i\le N$ the first $i-n_0$ entries of $\mu_i$
equal $\beta$. The \emph{plane part} $\pi(\lambda)$ and the \emph{tangent part} $\tau(\lambda)$ of $\lambda\in\Z^k$ are its first $r$
and its last $K$ entries. The \emph{weight} of a path is $w(\mu)=w_K(\tau(\mu_N))$, and $S_{\mathrm{all}}$ is the set of pairs
$(\mu,I)$ of a path and a standard family of shape $\tau(\mu_N)$.

Let moreover $p$ be a prime, $a,\Lambda>0$, numbers $b_i>0$ and indices $o(i)<i$ for $r<i\le N$, and put
$X_i=\Lambda b_i/(pk)$. A path is \emph{typical} if for all $i\le N$:
\begin{enumerate}[label=(T\arabic*)]
\item if $i$ is a covector letter, $\mu_{i-1}^{+,i-1}-\mu_i^{+,i}\le X_i$;
\item if $i>r$ is a vector letter, $\mu_{i-1}^{-,i-1}-\mu_i^{-,i}\le X_i$;
\item if $i>r$, $\frac23B_{o(i)}\le\mu_{i-1}(r)$;
\end{enumerate}
here $\lambda^{\pm,i}$ denotes $\lambda^{\pm}$ computed without the first $\max(i-n_0,0)$ entries (which equal $\beta$ by the head
condition). A path is \emph{good} if for all $r<i\le N$:
\begin{enumerate}[label=(G\arabic*)]
\item $\frac23B_{o(i)}\le\mu_{i-1}(r)$;
\item if $i$ is a vector letter, $\sum_c\bigl(\tau(\mu_i)(c)-\tau(\mu_{i-1})(c)\bigr)\le X_i+2B_{o(i)}$;
\item if $i$ is a covector letter, $\sum_c\bigl(\pi(\mu_{i-1})(c)-\pi(\mu_i)(c)\bigr)\le X_i$.
\end{enumerate}
$S_{\mathrm{good}}\subseteq S_{\mathrm{all}}$ is the set of pairs whose path is good. The data are \emph{admissible} if
\begin{enumerate}[label=(A\arabic*)]
\item $a\ge3$, $p\ge13$, $2r\le n_0$, $N+r\le K$, $k/2\le K-2N$ and $N\le a$;
\item $k\le B_i\le2a$ for $i\le n_0$, and $\frac{p^2}5B_i\le B_{i'}$ for $i<i'\le n_0$ of the same type;
\item the letters $i\le r$ and the letters $i>n_0-r$ are vector letters;
\item for $r<i\le N$, $o(i)$ is a vector letter and exactly $r$ vector letters $l$ satisfy $o(i)\le l<i$;
\item for $r<i\le N$, $B_l\le4b_i/p$ for all $l<i$ with $l\le n_0$; and $\Lambda>0$.
\end{enumerate}
Put $\ell_0=\log(1+2aN)$ and $k'=K-2N$.
\formal{Formalism.MainNotions.SelectionLemmas.Pf.T.HypT}
\formal{Formalism.MainNotions.SelectionLemmas.Pf.T.SAll}
\formal{Formalism.MainNotions.SelectionLemmas.Pf.T.SGood}
\formal{Formalism.InvariantFactorPaths.BranchingDimensions.Letters.Typical}
\end{definition}

\begin{lemma}[Paths are few]\label{pf:linkG:pathcount}
Let the data be admissible and $\mu$ a path. Then $\mu_i(m)\ge0$ for $m\le k-N$, $\mu_i(m)\le0$ for $m>N$, and
$|\mu_i(m)|\le\beta\le2aN$ for all $i,m$. The set of paths is finite, with at most $e^{2N^2\ell_0}$ elements.
\formallabel{Formalism.MainNotions.SelectionLemmas.Pf.T.path_count_finite}
\formallabel{Formalism.MainNotions.SelectionLemmas.Pf.T.path_le_beA}
\formallabel{Formalism.MainNotions.SelectionLemmas.Pf.T.beA_le}
\end{lemma}

\begin{proof}
If $\nu/\lambda$ is a horizontal strip with $\lambda$ dominant, then $\nu$ is dominant, has at most one more positive entry than
$\lambda$ (as $\nu(m+1)\le\lambda(m)$) and no more negative entries (as $\nu\ge\lambda$), and $\nu(1)\le\lambda(1)+d$; symmetrically
for $\lambda/\nu$. Starting from $\mu_0=0$, after $i\le N$ steps there are at most $N$ positive and at most $N$ negative entries,
which by dominance gives the sign conditions. For $i\le n_0$, $\mu_i(1)\le\sum_{l\le i\text{ vector}}B_l\le\beta$ and
$\mu_i(k)\ge-\sum_{l\le i\text{ covector}}B_l\ge-\beta$. For $i>n_0$ the letter $i$ is a vector letter by (A3), $\mu_i(1)=\beta$ by
the head condition, and $\mu_i(m)\le\mu_{i-1}(m-1)\le\beta$ for $m\ge2$ by induction; and $\mu_i(k)\ge\mu_{i-1}(k)\ge-\beta$. By (A2),
$\beta\le2an_0\le2aN$. As $K>2N$ by (A1), the entries $\mu_i(m)$ with $N<m\le k-N$ vanish, so a path is determined by the
$2N\cdot N$ entries $\mu_i(m)$ with $i\ge1$ and $m\le N$ or $m>k-N$, each of absolute value at most $\beta$. Hence there are at
most $(\beta+1)^{2N^2}\le(1+2aN)^{2N^2}=e^{2N^2\ell_0}$ paths.
\end{proof}

\begin{lemma}[Total weight]\label{pf:linkG:totalweight}
For admissible data, $S_{\mathrm{all}}$ is finite and $|S_{\mathrm{all}}|=\sum_\mu w(\mu)=\rho$.
\formallabel{Formalism.MainNotions.SelectionLemmas.Pf.T.path_weight_eq_rank}
\formallabel{Formalism.MainNotions.SelectionLemmas.Pf.T.pre_count}
\formallabel{Formalism.MainNotions.SelectionLemmas.Pf.T.tau_count}
\end{lemma}

\begin{proof}
Finiteness follows from \cref{pf:linkG:pathcount} and the finiteness of the standard families of a given shape; by (i),
$|S_{\mathrm{all}}|=\sum_\mu w(\mu)$. For $s\le N$ and an initial segment $(\mu_0,\dots,\mu_s)$ of a path let $A(\mu_0,\dots,\mu_s)$ be the
number of pairs $(\mu',I)\in S_{\mathrm{all}}$ with $\mu'_i=\mu_i$ for $i\le s$.

\emph{The tail.} Let $s=N-t\ge n_0$ ($t\le r$). By the head condition the first $r-t$ entries of $\mu_s$ equal $\beta$; let $\nu_s$ be
the weight of the last $K+t$ entries. We show $A(\mu_0,\dots,\mu_s)=w_{K+t}(\nu_s)$ by induction on $t$. For $t=0$ the path is
complete and $A=w_K(\tau(\mu_N))$ by (i). For $t\ge1$ the letter $s+1$ is a vector letter; the possible $\mu_{s+1}$ consist of
$r-t+1$ entries $\beta$ followed by a weight $\nu_{s+1}\in\Z^{K+t-1}$, and since all entries are at most $\beta$
(\cref{pf:linkG:pathcount}), $\mu_{s+1}/\mu_s$ is a horizontal strip exactly when $\nu_{s+1}$ interlaces $\nu_s$. By induction and
the branching rule (iii), $A(\mu_0,\dots,\mu_s)=\sum_{\nu_{s+1}}w_{K+t-1}(\nu_{s+1})=w_{K+t}(\nu_s)$.

\emph{The head.} For $s\le n_0$ we show $A(\mu_0,\dots,\mu_s)=w_k(\mu_s)\prod_{s<i\le n_0}\binom{B_i+k-1}{k-1}$ by descending
induction. For $s=n_0$ this is the tail statement with $t=r$ ($\nu_{n_0}=\mu_{n_0}$). For $s<n_0$, the possible $\mu_{s+1}$ are
the dominant weights with $\mu_{s+1}/\mu_s$ (or $\mu_s/\mu_{s+1}$) a horizontal strip of size $B_{s+1}$; by induction and Pieri's
rule (ii), $A(\mu_0,\dots,\mu_s)=\sum_{\mu_{s+1}}w_k(\mu_{s+1})\prod_{s+1<i\le n_0}\binom{B_i+k-1}{k-1}
=w_k(\mu_s)\prod_{s<i\le n_0}\binom{B_i+k-1}{k-1}$. For $s=0$, $w_k(0)=1$ gives $|S_{\mathrm{all}}|=A(\mu_0)=\rho$.
\end{proof}

\begin{lemma}[Typical paths are good]\label{pf:linkG:typicalgood}
For admissible data, every typical path is good.
\formallabel{Formalism.MainNotions.SelectionLemmas.Pf.T.typical_path_orders_small}
\formallabel{Formalism.MainNotions.SelectionLemmas.Pf.T.nu_last_le}
\formallabel{Formalism.MainNotions.SelectionLemmas.Pf.T.dplus_le_geom}
\formallabel{Formalism.MainNotions.SelectionLemmas.Pf.T.typ_of_Typical}
\end{lemma}

\begin{proof}
Let $\mu$ be a typical path and $r<i\le N$. By \cref{pf:linkG:pathcount} and $N+r\le K$, the entries $\mu_s(m)$ with $m\le r+1$
are non-negative, so the plane part contributes nothing to $\lambda^-$. Covector letters satisfy $i\le n_0-r$ by (A3), so
(T1) is the inequality $\mu_{i-1}^+-\mu_i^+\le X_i$; and as the omitted head entries are positive, (T2) is the inequality
$\mu_{i-1}^--\mu_i^-\le X_i$ computed with all entries. (G1) is (T3).

\emph{A geometric bound.} For $m\le l\le n_0$ with $l$ a vector letter,
$\sum_{l'<m,\ l'\text{ vector}}B_{l'}\le\frac{10}{p^2}B_l$. Indeed, by induction on $m$: if $m-1$ is a vector letter, the
induction hypothesis for $l=m-1$ and (A2) give $\sum_{l'<m}B_{l'}\le(1+\frac{10}{p^2})B_{m-1}\le(1+\frac{10}{p^2})\frac5{p^2}B_l\le
\frac{10}{p^2}B_l$, since $p^2\ge169$ gives $\frac{10}{p^2}\le1$.

\emph{The last plane entry.} We show $\mu_{i-1}(r)\le2B_{o(i)}$. Between $o(i)+1$ and $i-1$ there are $r-1$ vector letters by (A4),
and along a vector step an entry can only move up by one position under interlacing, while covector steps decrease
entries; hence $\mu_{i-1}(r)\le\mu_{o(i)}(1)$. Moreover $\mu_{o(i)}(1)\le\sum_{l\le o(i)\text{ vector}}B_l=
\sum_{l<o(i)\text{ vector}}B_l+B_{o(i)}\le(\frac{10}{p^2}+1)B_{o(i)}\le2B_{o(i)}$, using $o(i)\le n_0$ (which follows from (A4) and
$N=n_0+r$) and the geometric bound.

\emph{(G2).} Let $i$ be a vector letter. Writing $x=\max(x,0)-\max(-x,0)$ and using that the plane entries are
non-negative,
\[
\sum_c\bigl(\tau(\mu_i)(c)-\tau(\mu_{i-1})(c)\bigr)=\Bigl(\sum_c\tau(\mu_i)(c)^+-\sum_c\tau(\mu_{i-1})(c)^+\Bigr)+(\mu_{i-1}^--\mu_i^-).
\]
By interlacing, $\tau(\mu_i)(c+1)\le\tau(\mu_{i-1})(c)$, so $\sum_{c\le K}\tau(\mu_i)(c)^+\le\tau(\mu_i)(1)^++\sum_{c\le K}\tau(\mu_{i-1})(c)^+$;
and $\tau(\mu_i)(1)=\mu_i(r+1)\le\mu_{i-1}(r)$ with $\mu_{i-1}(r)\ge0$. Hence the sum is at most
$\mu_{i-1}(r)+(\mu_{i-1}^--\mu_i^-)\le2B_{o(i)}+X_i$ by the previous paragraph and (T2).

\emph{(G3).} Let $i$ be a covector letter. As the plane entries are non-negative, $\lambda^+=\sum_c\pi(\lambda)(c)+\sum_c\tau(\lambda)(c)^+$,
and $\tau(\mu_i)\le\tau(\mu_{i-1})$ entrywise. Hence $\sum_c(\pi(\mu_{i-1})(c)-\pi(\mu_i)(c))\le\mu_{i-1}^+-\mu_i^+\le X_i$ by (T1).
\end{proof}

\paragraph{Binomial sums.} For $k'\ge1$ and $y\ge0$ let
\[
f(y)=\sum_{j=1}^{k'}\log\frac{y+j}{j},\qquad f'(y)=\sum_{j=1}^{k'}\frac1{y+j},
\]
so that $f(n)=\log\binom{n+k'}{k'}$ for $n\in\N$, $f(0)=0$, and $f'$ is the derivative of $f$.

\begin{lemma}[Properties of $f$]\label{pf:linkG:fT}
For $y,y'\ge0$: $f'$ is positive and decreasing; $\frac{k'}{y+k'}\le f'(y)\le\frac{k'}{y+1}$; if $y\ge k'$ then $f'(y)\ge\frac{k'}{2y}$;
if $y\ge k'$ and $y'\ge4y$ then $f'(y')\le\frac12f'(y)$; and $f(x)-f(y)\le f'(y)(x-y)$ for $x\ge0$.
\formallabel{Formalism.InvariantFactorPaths.BranchingDimensions.logBinomialProperties}
\end{lemma}

\begin{proof}
All but the last two claims are termwise ($j\le k'$ gives $\frac1{y+k'}\le\frac1{y+j}\le\frac1{y+1}$, and $y+k'\le2y$ for $y\ge k'$).
If $y'\ge4y$ and $j\le k'\le y$, then $y'+j\ge4y\ge2(y+j)$, so $\frac1{y'+j}\le\frac12\cdot\frac1{y+j}$. Concavity: termwise,
$\log\frac{x+j}{y+j}\le\frac{x+j}{y+j}-1=\frac{x-y}{y+j}$.
\end{proof}

\begin{lemma}[Dominated sums]\label{pf:linkG:deficit}
Let $k'\ge1$, $L\in\N$, $\alpha_0,\dots,\alpha_{L-1}\ge0$ and $\beta_0,\dots,\beta_{L-1}$ with $\beta_l\ge k'$ and $4\beta_{l+1}\le\beta_l$.
Assume the tail sums satisfy $\delta_u:=\sum_{l\ge u}(\alpha_l-\beta_l)\le0$ for all $u<L$. Then
$\sum_lf(\alpha_l)-\sum_lf(\beta_l)\le0$; and if $\delta_{u_0}\le-X$ for some $u_0<L$ and $X\ge0$, then
\[
\sum_lf(\alpha_l)-\sum_lf(\beta_l)\le-\frac{Xf'(\beta_{u_0})}2\le-\frac{Xk'}{4\beta_{u_0}} .
\]
\formallabel{Formalism.InvariantFactorPaths.BranchingDimensions.logBinomialDeficit}
\end{lemma}

\begin{proof}
Put $c_u=f'(\beta_u)-f'(\beta_{u-1})$ for $u\ge1$ and $c_0=f'(\beta_0)$. As $\beta_{u-1}\ge4\beta_u$ and $\beta_u\ge k'$,
\cref{pf:linkG:fT} gives $f'(\beta_{u-1})\le\frac12f'(\beta_u)$, so $c_u\ge\frac12f'(\beta_u)\ge0$. By concavity and Abel
summation,
\[
\sum_l\bigl(f(\alpha_l)-f(\beta_l)\bigr)\le\sum_lf'(\beta_l)(\alpha_l-\beta_l)=\sum_{u<L}c_u\,\delta_u .
\]
Every term is $\le0$, which gives the first claim, and the term $u_0$ alone is at most $-c_{u_0}X\le-\frac12f'(\beta_{u_0})X$.
Finally $f'(\beta_{u_0})\ge\frac{k'}{2\beta_{u_0}}$ as $\beta_{u_0}\ge k'$.
\end{proof}

\begin{lemma}[Logarithm of the rank]\label{pf:linkG:logrank}
For admissible data and $k'=K-2N$, $0\le\log\rho-\sum_{i\le n_0}f(B_i)\le N(r+2N)\log(1+2a)$.
\formallabel{Formalism.InvariantFactorPaths.BranchingDimensions.logRankEstimate}
\end{lemma}

\begin{proof}
Write $k-1=k'+d$ with $d=r+2N-1$. Then $\log\binom{B_i+k-1}{k-1}=\sum_{j=1}^{k'+d}\log\frac{B_i+j}{j}$, and the terms with
$k'<j\le k'+d$ lie between $0$ and $\log(1+B_i)\le\log(1+2a)$. Sum over the at most $N$ indices $i\le n_0$.
\end{proof}

\begin{lemma}[Logarithm of Weyl's dimension]\label{pf:linkG:logweyl}
Let $2N\le K\le n\le r+K$, $k'=K-2N$, and let $\lambda\in\Z^n$ be dominant with at most $N$ positive and at most $N$ negative
entries, and $\lambda(1)^++\lambda(n)^-\le A$. Then $w_n(\lambda)>0$ and, with $L=\log(1+A)$,
\[
-2N^2L\le\log w_n(\lambda)-\sum_mf(\lambda(m)^+)-\sum_mf(\lambda(m)^-)\le(2Nr+8N^2)L .
\]
\formallabel{Formalism.InvariantFactorPaths.BranchingDimensions.logWeylDimensionEstimate}
\end{lemma}

\begin{proof}
Write $h(z,d)=\log\frac{z+d}d$, so that $\log w_n(\lambda)=\sum_{i<j}h(\lambda(i)-\lambda(j),j-i)$, every term in $[0,L]$. The positive
entries are $\lambda(1),\dots,\lambda(P)$ with $P\le N$ and the negative ones are among the last $N$. Split the pairs $i<j$: (a) both
entries non-zero; (b) $\lambda(j)=0$, contributing $h(\lambda(i)^+,j-i)$; (c) $\lambda(i)=0$, contributing $h(\lambda(j)^-,j-i)$ (pairs of
zeros contribute $0$). There are at most $(2N)^2$ pairs of type (a), so their sum lies in $[0,4N^2L]$. For a positive entry
$\lambda(i)$, the distances $D_i=\{j-i:\lambda(j)=0\}$ contain every $d\le k'$ except those with $\lambda(i+d)>0$ (at most $N$ of them;
$i+d\le N+k'<n-N$ excludes negative entries), and the elements of $D_i$ larger than $k'$ are at most $n-1-k'\le r+2N$ in number.
Hence $f(\lambda(i))-NL\le\sum_{d\in D_i}h(\lambda(i),d)\le f(\lambda(i))+(r+2N)L$. Summing over the at most $N$ positive entries, the
type (b) sum lies within $[-N^2L,N(r+2N)L]$ of $\sum_mf(\lambda(m)^+)$; symmetrically for (c). Adding the three types gives the
claim.
\end{proof}

\paragraph{Stages and tails.} For a path $\mu$ and $i\le N$, the \emph{stage} $\sigma_i$ is $\mu_i$ for $i\le n_0$, and $\mu_i$ without its
first $i-n_0$ entries (which equal $\beta$) for $i>n_0$; thus $\sigma_N=\tau(\mu_N)$. For a weight $\lambda$ and $u\ge0$ let
$T^+(\lambda,u)=\sum_{m>u}\lambda(m)^+$ and $T^-(\lambda,u)=\sum_{m>u}\lambda(n+1-m)^-$ (tails of the positive and negative parts,
counted from the largest, resp.\ the most negative, entry). For $i\le N$ let $V(i)=(V_0(i),V_1(i),\dots)$ be the sizes $B_l$ of
the vector letters $l\le\min(i,n_0)$, latest first, and $C(i)$ the same for covector letters; $S^\pm(u,i)$ is the sum of the
entries of $V(i)$, resp.\ $C(i)$, from position $u$ on.

\begin{lemma}[Tails are dominated]\label{pf:linkG:tails}
For every path and $i\le N$, $u\ge0$: $T^+(\sigma_i,u)\le S^+(u,i)$ and $T^-(\sigma_i,u)\le S^-(u,i)$. If
$T^+(\sigma_{i_0},u)\le S^+(u,i_0)-X$ for some $i_0,u$ and $X\ge0$, then $T^+(\sigma_N,u+v)\le S^+(u+v,n_0)-X$, where $v$ is the
number of vector letters in $(i_0,n_0]$; symmetrically for $T^-$ with covector letters.
\formallabel{Formalism.InvariantFactorPaths.BranchingDimensions.pathTailBounds}
\end{lemma}

\begin{proof}
Consider a step $i-1\to i$. If $i\le n_0$ is a vector letter, $\sigma_i/\sigma_{i-1}$ is a horizontal strip of size $B_i$; interlacing
$\sigma_i(m+1)\le\sigma_{i-1}(m)$ gives $T^+(\sigma_i,u+1)\le T^+(\sigma_{i-1},u)$, and $T^+(\sigma_i,0)\le T^+(\sigma_{i-1},0)+B_i$ since the
total size grows by $B_i$ and the negative parts do not grow. As $V(i)$ is $V(i-1)$ with $B_i$ put in front,
$S^+(u+1,i)=S^+(u,i-1)$ and $S^+(0,i)=S^+(0,i-1)+B_i$. If $i\le n_0$ is a covector letter, $\sigma_i\le\sigma_{i-1}$ entrywise and
$V(i)=V(i-1)$. If $i>n_0$, $\sigma_i$ interlaces $\sigma_{i-1}$ (the first entry of $\sigma_{i-1}$ is replaced by the head entry
$\beta$), so $T^+(\sigma_i,u)\le T^+(\sigma_{i-1},u)$, and $V(i)=V(i-1)$. Starting from $\sigma_0=0$, induction on $i$ gives the first
claim, and the same induction started at $i_0$ (each vector letter shifts $u$ by one) gives the propagation of a deficit.
The negative parts are treated by reversing the order of the entries and exchanging vector and covector letters.
\end{proof}

\begin{theorem}[Atypical paths have small weight]\label{pf:linkG:lemmaT}
For admissible data and every path $\mu$, $w(\mu)\le e^{9N^2\ell_0}\rho$, and if $\mu$ is not typical,
\[
w(\mu)\le\exp\Bigl(9N^2\ell_0-\min\Bigl(\frac{\Lambda k'}{16k},\frac{k'}{12}\Bigr)\Bigr)\rho .
\]
\formallabel{Formalism.InvariantFactorPaths.BranchingDimensions.Top.lemma_t_one}
\formallabel{Formalism.InvariantFactorPaths.BranchingDimensions.pathWeightBound_of}
\formallabel{Formalism.MainNotions.SelectionLemmas.Pf.T.atypical_path_small_weight}
\end{theorem}

\begin{proof}
Let $\kappa=\sigma_N=\tau(\mu_N)\in\Z^K$, so $w(\mu)=w_K(\kappa)$. By \cref{pf:linkG:pathcount}, $\kappa$ has at most $N$ positive and at most
$N$ negative entries; its entries lie between $\mu_{n_0}(k)\ge-\sum_{l\text{ covector}}B_l$ and $\mu_{n_0}(1)\le\sum_{l\text{ vector}}B_l$
(the tail steps interlace), so $\kappa(1)^++\kappa(K)^-\le\beta\le2aN$. \Cref{pf:linkG:logweyl} with $n=K$, $A=2aN$ gives
\[
\log w(\mu)\le\sum_mf(\kappa(m)^+)+\sum_mf(\kappa(m)^-)+(2Nr+8N^2)\ell_0 .
\]
By \cref{pf:linkG:tails} at $i=N$, the sequences $\alpha_l=\kappa(l+1)^+$ and $\beta_l=V_l(n_0)$ satisfy the hypothesis of
\cref{pf:linkG:deficit} (the entries $V_l$ are at least $k\ge k'$ by (A2), and $4V_{l+1}\le V_l$ by (A2) as $p^2\ge169$; entries of
$\kappa$ beyond the length of $V(n_0)$ are $\le0$ since their tail is bounded by $0$); likewise for the negative parts and
$C(n_0)$. Hence, writing $\Delta^\pm\ge0$ for the deficits produced by \cref{pf:linkG:deficit},
\begin{align*}
\sum_mf(\kappa(m)^+)+\sum_mf(\kappa(m)^-)&\le\sum_lf(V_l(n_0))+\sum_lf(C_l(n_0))-\Delta^+-\Delta^-\\
&=\sum_{i\le n_0}f(B_i)-\Delta^+-\Delta^-\le\log\rho-\Delta^+-\Delta^-,
\end{align*}
by \cref{pf:linkG:logrank}. Since $2r\le n_0\le N$, $2Nr+8N^2\le9N^2$, and therefore $w(\mu)\le\exp(9N^2\ell_0-\Delta^+-\Delta^-)\rho$.
It remains to show $\Delta^++\Delta^-\ge\min(\frac{\Lambda k'}{16k},\frac{k'}{12})$ if $\mu$ is not typical; there are three cases.
\begin{enumerate}[label=(\roman*)]
\item (T1) fails at a covector letter $i$ (so $r<i\le n_0$ by (A3)). As $V(i)=V(i-1)$,
$T^+(\sigma_i,0)<T^+(\sigma_{i-1},0)-X_i\le S^+(0,i)-X_i$. By \cref{pf:linkG:tails} the deficit $X_i$ propagates to $\kappa$ at the
position $v$, and $V_v(n_0)=B_l$ for the latest vector letter $l<i$. By \cref{pf:linkG:deficit}, $\Delta^+\ge X_ik'/(4B_l)$, and by
(A5), $B_l\le4b_i/p$, so $\Delta^+\ge\frac{\Lambda b_i}{pk}\cdot\frac{k'p}{16b_i}=\frac{\Lambda k'}{16k}$.
\item (T2) fails at a vector letter $i>r$: the same argument with the negative parts and covector letters gives
$\Delta^-\ge\frac{\Lambda k'}{16k}$.
\item (T3) fails at $i>r$: $\mu_{i-1}(r)<\frac23B_{o(i)}$. Let $u\ge0$ be such that $\mu_{i-1}(r)$ is the entry of $\sigma_{i-1}$ at position
$u+1$. By (A4) and the definition of the lists, $B_{o(i)}=V_u(i-1)$, so, as $\mu_{i-1}(r)\ge0$,
$T^+(\sigma_{i-1},u)=\mu_{i-1}(r)+T^+(\sigma_{i-1},u+1)\le\frac23V_u(i-1)+S^+(u+1,i-1)=S^+(u,i-1)-\frac13B_{o(i)}$:
a deficit $X=\frac13B_{o(i)}$, which propagates to $\kappa$ at an entry of $V(n_0)$ equal to $B_{o(i)}$. \Cref{pf:linkG:deficit} gives
$\Delta^+\ge\frac13B_{o(i)}\cdot\frac{k'}{4B_{o(i)}}=\frac{k'}{12}$.
\end{enumerate}
\end{proof}

\begin{proposition}[Lemma T]\label{pf:linkG:lemT}
For admissible data, $S_{\mathrm{good}}$ is finite and
\[
\Bigl(1-\exp\Bigl(11N^2\ell_0-\frac{\min(\Lambda,k)}{32}\Bigr)\Bigr)\rho\le|S_{\mathrm{good}}| .
\]
\formallabel{Formalism.MainNotions.SelectionLemmas.Pf.lemT_of_adm}
\formallabel{Formalism.MainNotions.SelectionLemmas.Pf.T.lemT_of}
\end{proposition}

\begin{proof}
$S_{\mathrm{good}}\subseteq S_{\mathrm{all}}$ is finite by \cref{pf:linkG:totalweight}. By \cref{pf:linkG:typicalgood}, every pair of
$S_{\mathrm{all}}$ whose path is typical lies in $S_{\mathrm{good}}$, so $\rho=|S_{\mathrm{all}}|\le|S_{\mathrm{good}}|+|S_{\mathrm{bad}}|$, where
$S_{\mathrm{bad}}$ is the set of pairs with an atypical path. The pairs with a given path $\mu$ are $w(\mu)$ in number; by
\cref{pf:linkG:lemmaT,pf:linkG:pathcount},
\[
|S_{\mathrm{bad}}|\le e^{2N^2\ell_0}\exp\bigl(9N^2\ell_0-M\bigr)\rho,\qquad M=\min\Bigl(\frac{\Lambda k'}{16k},\frac{k'}{12}\Bigr).
\]
By (A1), $k'=K-2N\ge k/2$, so $\frac{\Lambda k'}{16k}\ge\frac\Lambda{32}$ and $\frac{k'}{12}\ge\frac k{24}\ge\frac k{32}$; hence
$M\ge\min(\Lambda,k)/32$ and the claim follows.
\end{proof}

\paragraph{The letter dictionary.} We now apply this to a segment $j$ of a valid frame, with the data of clause 27 and degree
parameters $b$ admissible for $j$ (in the sense of Step 4 of the proof of \cref{pf:linkG:secondpart}). For a letter $e$ of the
segment put $P(e)=p^{\mathrm{dist}(j,e)+R_j}$, where $R_j$ is the logarithmic offset of the segment.

\begin{lemma}[The letter dictionary]\label{pf:linkG:dictionary}
Let $C_\Lambda,C_k\ge1088$. Let $e_1,\dots,e_N$ be the letters of the segment $j$ in their order, $n_0=N-r$, with the type of
each letter given by the frame, sizes $B_i$ equal to the segment degrees of $e_i$ (for $i\le n_0$ these are $b_x$ for the late
letters $x\in LL(j)$ and $(B_w(j,i)-m_0)^+$ for the $r$ letters of the frame of $j-1$), $b_i=a/P(e_i)$, $\Lambda=\Lambda_a$, and $o(i)$
the oldest vector letter with exactly $r$ vector letters in $[o(i),i)$. These data are admissible; the paths and good paths of the
segment are the paths and good paths of these data; and $\rho$ is the product of binomial coefficients in the segment bound.
\formallabel{Formalism.MainNotions.SelectionLemmas.Pf.TE.letter_dictionary_bounds}
\end{lemma}

\begin{proof}
The hypotheses of clause 27 give $a\ge416N_s$, $N_s\ge6$, $p\ge13$, $N\le N_s$, $r\le N_s$, $k\ge1088N_s^2$, and, with the top
count lower bounds (valid for constants $\ge32$), $\frac9{10}a\le a'_t\le2a$ for the reduced degrees $a'_t$ at the positions of
$TT$, and $a'_{t(e)}/P(e)\ge36kN_s$ for every letter $e$ of the segment. For $i\le n_0$ the segment degree satisfies
$a'_{t(e_i)}/P(e_i)-2kN_s\le B_i\le a'_{t(e_i)}/P(e_i)$. We check (A1)--(A5).
\emph{(A2), first part:} $B_i\ge36kN_s-2kN_s\ge k$, and $B_i\le2a/P(e_i)\le2a$.
\emph{(A2), growth:} if $i<i'\le n_0$ have the same type, then $\mathrm{dist}(j,e_i)\ge\mathrm{dist}(j,e_{i'})+2$, so
$P(e_i)\ge p^2P(e_{i'})$ and $\frac{p^2}5B_i\le\frac25\cdot\frac a{P(e_{i'})}$; on the other hand $kN_s\le
\frac{a'_t}{36P(e_{i'})}\le\frac a{18P(e_{i'})}$, so $B_{i'}\ge(\frac9{10}-\frac19)\frac a{P(e_{i'})}\ge\frac25\cdot\frac a{P(e_{i'})}\ge
\frac{p^2}5B_i$.
\emph{(A1):} $a\ge416\ge3$, $N\le N_s\le a$, $2r\le n_0$ and $N+r\le K$ by the validity of the frame, and
$K-2N=k-r-2N\ge k-3N_s\ge k/2$ as $k\ge1088N_s^2\ge6N_s$.
\emph{(A5):} for $l<i$, $\mathrm{dist}(j,e_l)>\mathrm{dist}(j,e_i)$, so $B_l\le2a/P(e_l)\le2a/(pP(e_i))\le4b_i/p$; and
$\Lambda_a=C_\Lambda N_s^2\log a>0$.
\emph{(A3), (A4):} the early letters and the letters of the frame are of vector type, and the frame provides, before every letter
$i>r$, at least $r$ vector letters, whence $o(i)$.
The remaining assertions are comparisons of definitions: the path condition of the segment is the path condition above
(the head condition is the condition on the frame entries), the threshold of the segment is
$X_i=\Lambda_ab_i/(pk)=\Lambda_aa/(pkP(e_i))$, and the conditions (G1)--(G3) imply the goodness conditions of the segment;
$\rho$ is the product over the $n_0$ early letters of $\binom{B_i+k-1}{k-1}$.
\end{proof}
\begin{lemma}[The final estimate]\label{pf:linkG:finalnum}
Under the hypotheses of clause 27 with $C_\Lambda,C_k\ge1088$, and for $N\le N_s$,
$\exp(11N^2\ell_0-\min(\Lambda_a,k)/32)\le a^{-N_s^2}$.
\formallabel{Formalism.MainNotions.SelectionLemmas.Pf.TE.final_num}
\end{lemma}

\begin{proof}
The size conditions of clause 27 give $a\ge416N_s$ and $N_s\ge6$, so $a\ge416$ and $\log a\ge1$. Hence $1+2aN\le1+2a\cdot
\frac a{416}\le a^2$ and $\ell_0\le2\log a$. Moreover $\Lambda_a=C_\Lambda N_s^2\log a\ge1088N_s^2\log a$ and $k\ge C_kN_s^2\log a\ge
1088N_s^2\log a$. Therefore $11N^2\ell_0-\min(\Lambda_a,k)/32\le22N_s^2\log a-34N_s^2\log a\le-N_s^2\log a$.
\end{proof}

\begin{theorem}[Good path weight]\label{pf:linkG:gpw}
Under the thirteen statements of the second-part list, for the data of clause 27 with $C_\Lambda,C_k\ge1088$, every segment $j$
and all admissible degree parameters $b$, the set of pairs of a good path of the segment and a standard family is finite,
and its cardinality is at least
\[
\bigl(1-a^{-N_s^2}\bigr)\prod_{x\in LL(j)}\binom{b_x+k-1}{k-1}\prod_{i<r}\binom{(B_w(j,i)-m_0)^++k-1}{k-1}.
\]
\formallabel{Formalism.MainNotions.SelectionLemmas.Pf.good_path_weight}
\formallabel{Formalism.MainNotions.SelectionLemmas.Pf.good_path_weight_of_top_count}
\formallabel{Formalism.MainNotions.SelectionLemmas.Pf.good_path_weight_of_parts}
\end{theorem}

\begin{proof}
By \cref{pf:linkG:dictionary}, the set in question contains the set $S_{\mathrm{good}}$ of the corresponding admissible data
and is contained in its $S_{\mathrm{all}}$, hence is finite. By \cref{pf:linkG:lemT,pf:linkG:finalnum} (with $N\le N_s$),
$|S_{\mathrm{good}}|\ge(1-a^{-N_s^2})\rho$, and $\rho$ is the stated product of binomial coefficients.
\end{proof}

% ---------------------------------------------------------------------------
\subsection{Choice and membership for Proposition 5}\label{sec:pf:linkG:pick}

We prove the choice part (i) of the proof of \cref{pf:linkG:prop5} (\cref{pf:linkG:pick} below) and give the main steps of the
membership part (ii) (\cref{pf:linkG:memb}). Throughout we fix test-side data with constant $C$ (\cref{pf:linkG:testdata}),
with its valid frame of $J$ segments, $k=m+1=r+K$, the multidegree $d=d_{TT}$ ($d_s$ the flat degree at $s\in TT$, $d_s=0$
otherwise), the integer $c'$ of $\linkGES$, the uniformiser $\varpi$ of $\OO$ with $\varpi^2=p$, the fraction field $K_f$ of
$\OO$, and the valuation $v_p$ on $K_f$ with $v_p(p)=1$. We put $L=|TT|k^2\log_pa$ and $\nu_2(d)=v_p(\prod_sd_s!)$.

\paragraph{Letters, levels and assembly.} For $j\in\Z/J$ the \emph{letters of level $j$} are the $n_L(j)$ late letters of the
segment $j$, in the increasing order of their positions, followed by the $r$ frame letters $\bar m$ of the segment, taken in
the order $\bar m=r-1,\dots,0$ (the frame letter $\bar m$ sits at the frame position $\mathrm{fr}(j,r-1-m)$). Let $\mathcal L$
be the set of all letters, $\mathrm{pos}(x)\in TT$ the position of a letter $x$ of level $j$, and
$P_x=p^{\mathrm{dist}(j,\mathrm{pos}(x))}$ its \emph{low power}; $N(j,u)$ denotes the set of the first $u$ letters of level
$j$. To every letter $x$ we attach $r$ \emph{plane variables} $X_{x,a}$ and $K$ \emph{normal variables} $Y_{x,b}$; for
$e:\mathcal L\to\N$ let $\mathcal P_r(e)$, resp.\ $\mathcal P_K(e)$, be the polynomials in the plane, resp.\ normal,
variables which are homogeneous of degree $e(x)$ in the variables of each letter $x$, and $\nu(e)=v_p(\prod_xe(x)!)$. For a
monomial $\beta$ we write $\beta!=\prod_v\beta(v)!$; the Fock pairing is $\langle F,G\rangle=\sum_\beta F_\beta G_\beta\beta!$ (in every
polynomial ring), and the \emph{$B$-valuation} is $v_B(F)=\min_{\beta}\bigl(v_p(F_\beta)+\frac12v_p(\beta!)\bigr)$, the minimum over the
monomials of $F$ ($v_B(0)=\infty$); $F$ is \emph{$B$-integral} if $v_B(F)\ge0$.

\begin{definition}[Block data, assembly, level images and chain data]\label{pf:linkG:blockdefs}
\emph{Block data} are maps $\rho:\mathcal L\times\{0,\dots,k-1\}\to\N$ and $d_{\mathrm{Pl}},d_{\mathrm{No}}:\mathcal L\to\N$ with
$\rho(x,c)<P_x$ and $\sum_c\rho(x,c)+P_x(d_{\mathrm{Pl}}(x)+d_{\mathrm{No}}(x))=d_{\mathrm{pos}(x)}$ for every letter $x$. For such
data, $A\in\mathcal P_r(d_{\mathrm{Pl}})$ and $N\in\mathcal P_K(d_{\mathrm{No}})$, the \emph{assembled polynomial} is
\[
\mathrm{asm}_\rho(A,N)=\prod_{x,c}X_{\mathrm{pos}(x),c}^{\rho(x,c)}\cdot A\bigl(X_{x,a}\mapsto X_{\mathrm{pos}(x),a}^{P_x}\bigr)\cdot
N\bigl(Y_{x,b}\mapsto X_{\mathrm{pos}(x),r+b}^{P_x}\bigr).
\]
A monomial $\mathbf m$ of the tangent ring has \emph{digits} $\rho$ if $\mathbf m(\mathrm{pos}(x),c)\equiv\rho(x,c)\bmod P_x$ for all
$x,c$; its \emph{quotients} are the plane and normal monomials $\beta(x,a)=\lfloor\mathbf m(\mathrm{pos}(x),a)/P_x\rfloor$ and
$\beta'(x,b)=\lfloor\mathbf m(\mathrm{pos}(x),r+b)/P_x\rfloor$, and its \emph{blocks} are their degrees at the letters. The \emph{level
image} $\mathrm{lev}(F)=\mathrm{lev}_{\rho,d_{\mathrm{Pl}},d_{\mathrm{No}}}(F)$ of $F$ is $\sum_{\mathbf m}F_{\mathbf m}X^\beta Y^{\beta'}$, the sum over the
monomials $\mathbf m$ of $F$ with digits $\rho$ and blocks $(d_{\mathrm{Pl}},d_{\mathrm{No}})$, viewed as a polynomial in the plane
variables with coefficients $\mathrm{lev}(F)_\beta$ in the normal variables. Put
\[
\gamma(P,\rho,\beta)=\frac{(\rho+P\beta)!}{\rho!\,\beta!\,(P!)^\beta},\qquad
c(\rho,d_{\mathrm{Pl}},d_{\mathrm{No}})=\prod_{x,c}\rho(x,c)!\prod_x(P_x!)^{d_{\mathrm{Pl}}(x)+d_{\mathrm{No}}(x)},
\]
$\gamma_{\mathrm{Pl}}(\beta)=\prod_{x,a}\gamma(P_x,\rho(x,a),\beta(x,a))$ and $\gamma_{\mathrm{No}}(\beta')=\prod_{x,b}\gamma(P_x,\rho(x,r+b),\beta'(x,b))$;
for a function $g$ on monomials, $g\star F=\sum_\beta g(\beta)F_\beta X^\beta$. Finally
\[
D_{\mathrm{No}}(d_{\mathrm{No}})=\sum_{j\in\Z/J}\ \sum_{u=0}^{n_L(j)+r}\binom K2\Bigl\lfloor\log_p\bigl(2\deg_{N(j,u)}d_{\mathrm{No}}+K-1\bigr)\Bigr\rfloor,
\]
where $\deg_Sd_{\mathrm{No}}=\sum_{x\in S}d_{\mathrm{No}}(x)$.

\emph{Signed chain data} with $J$ levels consist of the numbers $n_L(j)$ (which determine the letters as above), a type
(vector or covector) for every late letter, plane degrees $d(x)$, exponents $g(j,m)$, increasing chains of sets of letters of
level $j$ containing all frame letters, thresholds $\theta(j,i)$ and slopes $s(j,i)$ along the chains, link thresholds
$\theta_L(j,m)>0$ ($m+2\le r$), a unit-valued function $\gamma$ on plane monomials such that $\gamma^{\pm1}$ is congruent modulo $p^t$
to an $\OO$-combination of products of binomial coefficients $\prod_v\binom{\beta(v)}{\mu(v)}$ with $\sum_v\mu(v)\le6t$ (for every
$t\ge1$), and a constant $c\ge1$. They define the \emph{$B$-form} $B(F,A)=\langle F,\gamma\star A\rangle$, the number $n_0=\frac12\nu(d)$,
the \emph{$X$-space} and the \emph{$Y$-spaces} (valuation conditions along the chains and links, of
\cref{pf:linkG:teststeps}); the $Y$-spaces consist of polynomials in the \emph{model variables}: for every level $j$, the plane
variables of its late letters and three copies (label, wrapped and top variables) of those of its frame letters, and the substitution
$\varphi$ maps them to the plane variables of the letters, a wrapped variable of the frame letter $\bar m$ of level $j$ to the
$p^{g(j,m)}$-th power of the corresponding variable. They are \emph{tied to the frame with weights} $\kappa=(\kappa(j,u))$, $\kappa(j,u)\in\Z^K$, and plane
degrees $d_{\mathrm{Pl}}$ if $n_L(j)$ is the number of late letters of the segment $j$, the types are the parities of the positions,
$d=d_{\mathrm{Pl}}$, $g(j,m)$ is the shift of the frame from the segment $j$ to $j+1$, the chains are the sets
$\{\text{frame letters}\}\cup\{\text{late letters }l<i\}$ ($i\le n_L(j)$), and
$\theta(j,i)=\theta^*_\kappa(j,i)-\frac18s_U(j,i)$, $\theta_L(j,m)=\theta^*_\kappa(j,u_m)-\frac38s_U(j,u_m)$ with $u_m=n_L(j)+r-1-m$. Here
$s_U(j,u)=\lfloor B_j(x)/(Yk^2)\rfloor$ is the slope of the $u$-th letter $x$ of level $j$, $B_j(x)>0$ its block bound, and the
\emph{tie threshold} $\theta^*_\kappa(j,u)$ is the maximum of $B_j(x)/2$, of $\kappa(j,u+1)_0+m_*+s_U(j,u)$ if the position of $x$ is
even, and of $\kappa(j,u-1)_0+m_*+s_U(j,u)$ if $u\ge1$ and the position of the $(u-1)$-th letter is odd ($m_*$ is the tie constant of
the frame). They satisfy the \emph{degree condition of the frame} if moreover $c=c'$ and
$s(j,i)=\frac38s_U(j,i)$. A \emph{norm test} for such data is an integral $N_0\in\mathcal P_K(d_{\mathrm{No}})$ whose lowest weight with
respect to the letters of $N(j,u)$ is at most $\kappa(j,u)_0+4c'$, for all $j$ and $u\le n_L(j)+r$. A polynomial $F$ of the tangent
ring satisfies the \emph{membership condition} $\mathsf{Mem}$ if for all chain data tied to the frame (with any weights $\kappa$ and
plane degrees $d_{\mathrm{Pl}}$), all $\rho,d_{\mathrm{No}}$, every joint projector $e$ for $\kappa$ (\cref{pf:linkG:jointproj}) and every
normal monomial $\nu$, the plane polynomial $\sum_\beta\bigl(e(\mathrm{lev}(F)_\beta)\bigr)_\nu X^\beta$ lies in the $X$-space.
\formal{Formalism.MainNotions.Proposition5Numerics.BlockOK}
\formal{Formalism.MainNotions.Proposition5Numerics.asm}
\formal{Formalism.MainNotions.TestSideWords.levImg}
\formal{Formalism.MainNotions.Proposition5Numerics.DnoOf}
\formal{Formalism.MainNotions.PlaneGramData.SignedChainData}
\formal{Formalism.MainNotions.TestSideWords.Tied}
\formal{Formalism.MainNotions.Proposition5Numerics.SignedChainDegreeCondition}
\formal{Formalism.MainNotions.Proposition5Numerics.NormTest}
\formal{Formalism.MainNotions.TestSideWords.MemAt}
\end{definition}

\begin{lemma}[Digit splitting]\label{pf:linkG:digits}
Let $F$ be a polynomial of multidegree $d$.
\begin{enumerate}[label=(\alph*)]
\item For maps $\rho,d_{\mathrm{Pl}},d_{\mathrm{No}}$ with $\rho(x,c)<P_x$, and $A\in\mathcal P_r(d_{\mathrm{Pl}})$, $N\in\mathcal P_K(d_{\mathrm{No}})$,
\[
\langle F,\mathrm{asm}_\rho(A,N)\rangle=c(\rho,d_{\mathrm{Pl}},d_{\mathrm{No}})\sum_{\beta,\beta'}\mathrm{lev}(F)_{\beta,\beta'}\,
\gamma_{\mathrm{Pl}}(\beta)A_\beta\beta!\;\gamma_{\mathrm{No}}(\beta')N_{\beta'}\beta'! .
\]
\item The coefficients $\mathrm{lev}(F)_\beta$ lie in $\mathcal P_K(d_{\mathrm{No}})$.
\item Every monomial $\mathbf m$ of $F$, with digits $\rho$, quotients $\beta,\beta'$ and blocks $(d_{\mathrm{Pl}},d_{\mathrm{No}})$, defines block
data, $\mathrm{lev}_{\rho,d_{\mathrm{Pl}},d_{\mathrm{No}}}(F)_{\beta,\beta'}=F_{\mathbf m}$, and
$\mathbf m!=c(\rho,d_{\mathrm{Pl}},d_{\mathrm{No}})\,\beta!\,\beta'!\,\gamma_{\mathrm{Pl}}(\beta)\gamma_{\mathrm{No}}(\beta')$.
\end{enumerate}
\formallabel{Formalism.MainNotions.Proposition5Numerics.digFacts}
\formallabel{Formalism.MainNotions.Proposition5Numerics.digFacts_of}
\end{lemma}

\begin{proof}
For a valid frame, $x\mapsto\mathrm{pos}(x)$ is a bijection from the letters onto $TT$
(\lref{Formalism.MainNotions.Proposition5Numerics.dg_surj}: every position of $TT$ is the position of a letter;
\lref{Formalism.MainNotions.Proposition5Numerics.posTL_inj}: distinct letters have distinct positions).
As $d_s=0$ for $s\notin TT$, every monomial $\mathbf m$ of $F$ involves only variables at positions of $TT$, and division with
remainder by $P_x$ writes it uniquely as $\mathbf m=\ell(\rho)+P\cdot(\beta,\beta')$, where $\ell(\rho)$ is the exponent of the first factor
of $\mathrm{asm}_\rho$. For each variable $v$ at a letter $x$, with remainder $\rho_v<P_x$ and quotient $\beta_v$,
$(\rho_v+P_x\beta_v)!=\rho_v!\,\beta_v!\,(P_x!)^{\beta_v}\gamma(P_x,\rho_v,\beta_v)$ by the definition of $\gamma$; the product over all variables
gives (c), since the exponents of $P_x!$ at $x$ add up to $d_{\mathrm{Pl}}(x)+d_{\mathrm{No}}(x)$, and $\sum_c\rho(x,c)+P_x(d_{\mathrm{Pl}}(x)+
d_{\mathrm{No}}(x))=d_{\mathrm{pos}(x)}$ because $\mathbf m$ has multidegree $d$. The monomials of $\mathrm{asm}_\rho(A,N)$ are the
$\ell(\rho)+P\cdot(\beta,\beta')$, with coefficients $A_\beta N_{\beta'}$; by uniqueness, the monomials of $F$ which also occur there are
those with digits $\rho$ and quotients $(\beta,\beta')$, hence with blocks $(d_{\mathrm{Pl}},d_{\mathrm{No}})$ by homogeneity of $A$ and $N$.
Inserting the factorial identity into $\langle F,\mathrm{asm}_\rho(A,N)\rangle=\sum_{\mathbf m}F_{\mathbf m}\,\mathrm{asm}_\rho(A,N)_{\mathbf m}\,\mathbf m!$
gives (a), and (b) holds as the blocks of the monomials in $\mathrm{lev}(F)$ are $(d_{\mathrm{Pl}},d_{\mathrm{No}})$.
\end{proof}

\begin{lemma}[The digit factors]\label{pf:linkG:gammaunit}
Let $p\ge5$, $P=p^g$ and $0\le\rho<P$. Then $v_p(\gamma(P,\rho,\beta))=0$ for every $\beta\in\N$. Moreover, for maps $g,\rho$ on a finite
set with $\rho(v)<p^{g(v)}$, every $t\in\N$ and both signs, $\bigl(\prod_v\gamma(p^{g(v)},\rho(v),\beta(v))\bigr)^{\pm1}$ is congruent
modulo $p^t$, uniformly in $\beta$, to an $\OO$-combination of products $\prod_v\binom{\beta(v)}{\mu(v)}$ with $\sum_v\mu(v)\le4t$. In
particular $\gamma_{\mathrm{Pl}}$ and $\gamma_{\mathrm{No}}$ take unit values, and $\gamma_{\mathrm{Pl}}$ is a unit-valued function as in the
definition of chain data.
\formallabel{Formalism.MainNotions.Proposition5Numerics.gaAt}
\formallabel{Formalism.MainNotions.Proposition5Numerics.digAbs}
\formallabel{Formalism.MainNotions.Proposition5Numerics.digQ}
\end{lemma}

\begin{proof}
We use Legendre's formula $v_p(n!)=\sum_{i\ge1}\lfloor n/p^i\rfloor$, which follows from
$v_p(n!)=\sum_{l\le n}v_p(l)=\sum_{i\ge1}\#\{l\le n:p^i\mid l\}$. For $i\le g$,
$\lfloor(\rho+P\beta)/p^i\rfloor=\lfloor\rho/p^i\rfloor+p^{g-i}\beta$, and for $i>g$, $\lfloor(\rho+P\beta)/p^i\rfloor=\lfloor\beta/p^{i-g}\rfloor$ as
$\rho<P$. Hence $v_p((\rho+P\beta)!)=v_p(\rho!)+\beta\frac{p^g-1}{p-1}+v_p(\beta!)=v_p(\rho!)+v_p(\beta!)+\beta\,v_p(P!)$, which is the first
claim. The approximation by binomial combinations (with rational coefficients integral at $p$, which lie in $\OO$) is the first
part of the factorial digit lemma of the development (\lref{Formalism.MainNotions.G.GLWeightTheory.lemma_five_two_first_holds}:
for $p\ge5$ these products and their inverses are congruent modulo $p^t$ to $p$-integral combinations of products of binomial
coefficients of total degree at most $4t$). Since $4t\le6t$, $\gamma_{\mathrm{Pl}}$ satisfies the condition on $\gamma$.
\end{proof}

\begin{lemma}[Multinomial valuations]\label{pf:linkG:multinomial}
Let $e:I\to\N$ on a finite set $I$, $n\ge1$, and let $\beta$ be a monomial in variables $v_{x,c}$ ($x\in I$, $c<n$) with
$\sum_c\beta(x,c)=e(x)$. Then $0\le\nu(e)-v_p(\beta!)\le\sum_xn\lfloor\log_pe(x)\rfloor$ (with $\lfloor\log_p0\rfloor=0$). Consequently, if $G$ has multidegree $d$ and integral
coefficients and all coefficients of $T$ have valuation $\ge t$, then $v_p(\langle G,T\rangle)\ge t+\nu_2(d)-\kappa_{\mathrm{fl}}$.
\formallabel{Formalism.MainNotions.Proposition5Numerics.mult_level}
\formallabel{Formalism.MainNotions.Proposition5Numerics.padicVal_expFact_le}
\formallabel{Formalism.MainNotions.Proposition5Numerics.geQ_expFact}
\formallabel{Formalism.MainNotions.Proposition5Numerics.geQ_fock}
\end{lemma}

\begin{proof}
For one letter, $e!/\prod_c\beta_c!$ is a multinomial coefficient, so the difference is non-negative, and by Legendre's formula it equals
$\sum_{i\ge1}\bigl(\lfloor e/p^i\rfloor-\sum_c\lfloor\beta_c/p^i\rfloor\bigr)$. Each term lies in $[0,n-1]$, since the fractional parts of the
$\beta_c/p^i$ add up to less than $n$, and vanishes for $p^i>e$; so the sum is at most $(n-1)\lfloor\log_pe\rfloor$ (and $0$ if $e=0$). Sum over
the letters. For the consequence, $\langle G,T\rangle=\sum_{\mathbf m}G_{\mathbf m}T_{\mathbf m}\mathbf m!$ over monomials of multidegree $d$ in $k$
variables per position, and $v_p(\mathbf m!)\ge\nu_2(d)-\sum_sk\lfloor\log_pd_s\rfloor=\nu_2(d)-\kappa_{\mathrm{fl}}$.
\end{proof}

\begin{lemma}[Sizes]\label{pf:linkG:sizes}
Let $C\ge4$. (a) For block data and monomials $\beta\in\mathcal P_r(d_{\mathrm{Pl}})$, $\beta'\in\mathcal P_K(d_{\mathrm{No}})$,
\[
D_{\mathrm{No}}(d_{\mathrm{No}})+\tfrac12\bigl(\nu(d_{\mathrm{Pl}})-v_p(\beta!)\bigr)+\tfrac12\bigl(\nu(d_{\mathrm{No}})-v_p(\beta'!)\bigr)\le9L .
\]
(b) For every level $j$ and $u<n_L(j)+r$, with $x$ the $u$-th letter, $\frac38s_U(j,u)<\frac12B_j(x)$. (c) $c'\ge1$ and $J\ge1$.
\formallabel{Formalism.MainNotions.Proposition5Numerics.sizeStep}
\formallabel{Formalism.MainNotions.Proposition5Numerics.sizeStep'}
\formallabel{Formalism.MainNotions.Proposition5Numerics.sizeAt_of}
\formallabel{Formalism.MainNotions.Proposition5Numerics.DnoOf_le}
\formallabel{Formalism.MainNotions.Proposition5Numerics.mult_level_real}
\formallabel{Formalism.MainNotions.Proposition5Numerics.natlog_le_three}
\formallabel{Formalism.MainNotions.Proposition5Numerics.size_fin}
\formallabel{Formalism.MainNotions.Proposition5Numerics.sl_lt_Bo}
\formallabel{Formalism.MainNotions.Proposition5Numerics.sl_lt_real}
\formallabel{Formalism.MainNotions.Proposition5Numerics.card_TL_le}
\formallabel{Formalism.MainNotions.Proposition5Numerics.card_idx}
\formallabel{Formalism.MainNotions.Proposition5Numerics.J_le_card}
\formallabel{Formalism.MainNotions.Proposition5Numerics.triangular_count_le_sq}
\end{lemma}

\begin{proof}
By the hypotheses of $\linkGES$ with $C\ge4$ (\lref{Formalism.MainNotions.Proposition5Numerics.k_le_a} and \lref{Formalism.MainNotions.Proposition5Numerics.nb_of}:
the elementary size relations between $a$, $p$, $k$, $r$, $|TT|$ and $C$) we have $a\ge2$, $p\le a$, $\log_pa\ge1$,
$4|TT|\le C|TT|\le a$, $r\le k\le a$, and $a\le d_s\le2a$ for $s\in TT$, by the weight hypotheses; moreover $Y\ge2$, $c'\ge1$ (as
$c_*\ge1$), $J\ge1$ by the validity of the frame, and the number of letters is at most $|TT|$, as distinct letters have distinct
positions in $TT$ (\lref{Formalism.MainNotions.Proposition5Numerics.posTL_inj}).
Write $N=|TT|$ and $\ell=\log_pa$.

(a) By the block condition, $P_x(d_{\mathrm{Pl}}(x)+d_{\mathrm{No}}(x))\le d_{\mathrm{pos}(x)}\le2a$, so $d_{\mathrm{Pl}}(x),d_{\mathrm{No}}(x)\le2a$. If
$1\le M\le a^3$, then $\lfloor\log_pM\rfloor\le3\ell$. \Cref{pf:linkG:multinomial} gives
$\nu(d_{\mathrm{Pl}})-v_p(\beta!)\le N\cdot r\cdot3\ell$ and $\nu(d_{\mathrm{No}})-v_p(\beta'!)\le N\cdot K\cdot3\ell$, as $2a\le a^3$. In
$D_{\mathrm{No}}$, $\binom K2\le k^2$ and $2\deg_{N(j,u)}d_{\mathrm{No}}+K-1\le2\cdot N\cdot2a+k\le a^2+a\le a^3$ (using $4N\le a$, $k\le a$,
$a\ge2$), so each term is at most $3k^2\ell$; the number of terms is $\sum_j(n_L(j)+r+1)=|\mathcal L|+J\le2|\mathcal L|\le2N$ (each
level has $r\ge1$ letters). Hence $D_{\mathrm{No}}\le6Nk^2\ell$, and the left side of (a) is at most
$6Nk^2\ell+\frac32N(r+K)\ell\le6Nk^2\ell+\frac32Nk^2\ell\le9Nk^2\ell$.

(b) The block bound is positive, $B_j(x)>0$, as the frame is valid and $d_s\ge a>0$ on $TT$. Since $Yk^2\ge2$,
$s_U(j,u)\le B_j(x)/(Yk^2)\le B_j(x)/2$, so $\frac38s_U(j,u)\le\frac3{16}B_j(x)<\frac12B_j(x)$.
\end{proof}

\begin{lemma}[Joint projectors]\label{pf:linkG:jointproj}
Let chain data with $J$ levels be given, $d_{\mathrm{No}}:\mathcal L\to\N$ and $y\in\mathcal P_K(d_{\mathrm{No}})$. There are weights $\kappa(j,u)\in
\Z^K$ ($u\le n_L(j)+r$) and a linear operator $e$ on the normal polynomials (a \emph{joint projector for $\kappa$}) such that:
\begin{enumerate}[label=(\alph*)]
\item for every $(j,u)$, $e$ maps into the isotypic component of lowest weight $\kappa(j,u)$ for the action of $\GL_K$ on the
variables of the letters of $N(j,u)$; $e$ is the identity on the joint isotypic component with the weights $\kappa$ and vanishes
on the joint components with other weights;
\item $\langle eF,G\rangle=\langle F,eG\rangle$ for all $F,G$;
\item $e$ preserves $\mathcal P_K(d_{\mathrm{No}})$, and $p^{D_{\mathrm{No}}(d_{\mathrm{No}})}eF$ is $B$-integral for every $B$-integral
$F\in\mathcal P_K(d_{\mathrm{No}})$;
\item $v_B(ey)\le v_B(y)$.
\end{enumerate}
\formallabel{Formalism.MainNotions.TestSideWords.restJP}
\formallabel{Formalism.MainNotions.TestSideWords.greedy_list}
\formallabel{Formalism.MainNotions.TestSideWords.greedy_one}
\formallabel{Formalism.MainNotions.TestSideWords.vBB_sum}
\formallabel{Formalism.MainNotions.TestSideWords.jointE_int}
\formallabel{Formalism.MainNotions.TestSideWords.jointE_int_list}
\formallabel{Formalism.MainNotions.TestSideWords.eeN_int}
\end{lemma}

\begin{proof}
For an index $i=(j,u)$ and $\mu\in\Z^K$ let $e_{i,\mu}$ be the projector onto the isotypic component of lowest weight $\mu$ for the
action of $\GL_K$ on the variables of $N(j,u)$, so that $\sum_\mu e_{i,\mu}=\mathrm{id}$ (finitely many non-zero terms on each
polynomial). By the filtered projector corollary
(\lref{Formalism.MainNotions.G.GLWeightTheory.filtered_projector_corollary_field}: the isotypic projectors preserve the
degrees, and on homogeneous polynomials of degree $\delta$ in the variables of a set of letters their denominators on the
$B$-integral lattice are at most $p^{\binom K2\lfloor\log_p(2\delta+K-1)\rfloor}$), $p^{\mathrm{cost}(i)}e_{i,\mu}$ maps $B$-integral elements of
$\mathcal P_K(d_{\mathrm{No}})$ to $B$-integral elements, with $\mathrm{cost}(i)=\binom K2\lfloor\log_p(2\deg_{N(j,u)}d_{\mathrm{No}}+K-1)\rfloor$. The
projectors for different indices commute, and the composite $e=\prod_ie_{i,\kappa(i)}$ has the properties (a), (b) and the degree
part of (c) (\lref{Formalism.MainNotions.TestSideWords.jointN_props}: each $e_{i,\mu}$ is self-adjoint for the Fock pairing,
preserves $\mathcal P_K(d_{\mathrm{No}})$ and commutes with the others). Applying the factors one after the other, the denominators add
up: $p^{\sum_i\mathrm{cost}(i)}e$ maps $B$-integral elements of $\mathcal P_K(d_{\mathrm{No}})$ to $B$-integral elements, and
$\sum_i\mathrm{cost}(i)=D_{\mathrm{No}}(d_{\mathrm{No}})$. This is (c).

It remains to choose $\kappa$ with (d). The $B$-valuation satisfies $v_B(\sum_tF_t)\ge\min_tv_B(F_t)$, since $v_p$ is non-archimedean
and the coefficient of $X^\beta$ in the sum is the sum of the coefficients. Enumerate the indices $i_1,\dots,i_n$ and put $y_0=y$.
Given $y_{s-1}$, $y_{s-1}=\sum_\mu e_{i_s,\mu}y_{s-1}$ implies $\min_\mu v_B(e_{i_s,\mu}y_{s-1})\le v_B(y_{s-1})$ (if $y_{s-1}=0$, any $\mu$ will
do); choose $\kappa(i_s)$ attaining the minimum and put $y_s=e_{i_s,\kappa(i_s)}y_{s-1}$. Then $ey=y_n$ and
$v_B(y_n)\le\dots\le v_B(y_0)$.
\end{proof}

\begin{lemma}[Norm tests]\label{pf:linkG:normtest}
Let $\rho$ satisfy $\rho(x,c)<P_x$, let $e$ be a joint projector for $\kappa$ as in \cref{pf:linkG:jointproj}, and let $\nu$ be a normal
monomial of multidegree $d_{\mathrm{No}}$. There are a norm test $N_0$ for $\kappa$ and $c_Y\in K_f$ with
$v_p(c_Y)=\frac12\nu(d_{\mathrm{No}})+D_{\mathrm{No}}(d_{\mathrm{No}})-\frac12v_p(\nu!)$ such that all coefficients of
$N_0-\gamma_{\mathrm{No}}^{-1}\star\bigl(c_Y\,e(Y^\nu)\bigr)$ have valuation at least $c'$.
\formallabel{Formalism.MainNotions.Proposition5Numerics.noAt}
\formallabel{Formalism.MainNotions.Proposition5Numerics.coefGe_of_inBB}
\end{lemma}

\begin{proof}
Put $n_2=\nu(d_{\mathrm{No}})$, $v_\nu=v_p(\nu!)$, $D=D_{\mathrm{No}}(d_{\mathrm{No}})$ and $c_Y=\varpi^{n_2}p^D\varpi^{-v_\nu}$; then
$v_p(c_Y)=\frac{n_2}2+D-\frac{v_\nu}2$. The polynomial $\varpi^{-v_\nu}Y^\nu\in\mathcal P_K(d_{\mathrm{No}})$ has $B$-valuation
$-\frac{v_\nu}2+\frac{v_\nu}2=0$, so $p^De(\varpi^{-v_\nu}Y^\nu)$ is $B$-integral by \cref{pf:linkG:jointproj}(c). A $B$-integral
$F\in\mathcal P_K(d_{\mathrm{No}})$ satisfies $v_p(F_{\beta'})\ge-\frac12v_p(\beta'!)\ge-\frac12n_2$ for all $\beta'$ (\cref{pf:linkG:multinomial}), so
$\varpi^{n_2}F$ is integral. Hence $Z=c_Y\,e(Y^\nu)=\varpi^{n_2}\,p^De(\varpi^{-v_\nu}Y^\nu)$ is integral, and it lies in
$\mathcal P_K(d_{\mathrm{No}})$. By \cref{pf:linkG:gammaunit} there is an $\OO$-combination $\Lambda(\beta')=\sum_\mu\lambda_\mu\prod_v\binom{\beta'(v)}{\mu(v)}$
with $\sum_v\mu(v)\le4c'$ for $\lambda_\mu\ne0$ and $v_p(\Lambda(\beta')-\gamma_{\mathrm{No}}(\beta')^{-1})\ge c'$ for all $\beta'$. Put
$N_0=\Lambda\star Z$. It is integral (the values of $\Lambda$ are integral), lies in $\mathcal P_K(d_{\mathrm{No}})$, and
$N_0-\gamma_{\mathrm{No}}^{-1}\star Z=(\Lambda-\gamma_{\mathrm{No}}^{-1})\star Z$ has coefficients of valuation $\ge c'$. For the lowest weights: by
\cref{pf:linkG:jointproj}(a), $Z$ lies in the isotypic component of lowest weight $\kappa(j,u)$ for $N(j,u)$, so its lowest weight is at most
$\kappa(j,u)_0$ (\lref{Formalism.MainNotions.Proposition5Numerics.lowest_weight_le_of_isotypic}); and the operator $\Lambda\star$, a
combination of the operators $X^\beta\mapsto\prod_v\binom{\beta(v)}{\mu(v)}X^\beta$ of order $\sum_v\mu(v)\le4c'$, raises the lowest weight by at most
$4c'$ (\lref{Formalism.MainNotions.G.GLWeightTheory.degree_truncation}: the degree truncation lemma of the weight theory). Hence
$N_0$ is a norm test for $\kappa$.
\end{proof}

\begin{theorem}[Choice of the chain data]\label{pf:linkG:pick}
For $C\ge4$ the choice part of \cref{pf:linkG:prop5} holds with $C_{Dn}=9$: let $G$ be a polynomial of multidegree $d$ with integral
coefficients, not all in $\varpi\OO$, which satisfies $\mathsf{Mem}$. Then there are signed chain data $\mathcal D$ tied to the frame
with weights $\kappa$ and satisfying the degree condition of the frame, block data $(\rho,d_{\mathrm{Pl}},d_{\mathrm{No}})$, an element
$F_2$ of the $X$-space of $\mathcal D$ with $v_B(F_2)=0$, a norm test $N_0$ for $\kappa$ and $k_0\in K_f$ with
$v_p(k_0)\le\nu_2(d)-n_0+9L$, such that for every integral $A\in\mathcal P_r(d_{\mathrm{Pl}})$
\[
v_p\bigl(\langle G,\mathrm{asm}_\rho(A,N_0)\rangle-k_0\,B(F_2,A)\bigr)\ge c'+\nu_2(d)-\kappa_{\mathrm{fl}} .
\]
\formallabel{Formalism.MainNotions.Proposition5Numerics.restPick}
\formallabel{Formalism.MainNotions.Proposition5Numerics.pickStep_all}
\formallabel{Formalism.MainNotions.Proposition5Numerics.pickAt_of}
\formallabel{Formalism.MainNotions.Proposition5Numerics.pick_core}
\formallabel{Formalism.MainNotions.Proposition5Numerics.value_exact}
\formallabel{Formalism.MainNotions.Proposition5Numerics.gram_exponent_arithmetic}
\formallabel{Formalism.MainNotions.Proposition5Numerics.exists_unit_coeff}
\formallabel{Formalism.MainNotions.Proposition5Numerics.jpok_upd}
\end{theorem}

\begin{proof}
\emph{Step 1: a unit coefficient.} The values of $v_p$ on $\OO\smallsetminus\{0\}$ lie in $\frac12\N$; as the coefficients of $G$ are
integral and not all in $\varpi\OO$, some coefficient $G_{\mathbf m_0}$ is a unit. Let $(\rho,d_{\mathrm{Pl}},d_{\mathrm{No}})$ be the digits and
blocks of $\mathbf m_0$ and $\beta_0,\beta'_0$ its quotients. By \cref{pf:linkG:digits}(c) these are block data, the coefficient
$y=\mathrm{lev}(G)_{\beta_0}\in\mathcal P_K(d_{\mathrm{No}})$ satisfies $y_{\beta'_0}=G_{\mathbf m_0}$, and
$v_p(\mathbf m_0!)=v_p(c(\rho,d_{\mathrm{Pl}},d_{\mathrm{No}}))+v_p(\beta_0!)+v_p(\beta'_0!)$, as the digit factors are units
(\cref{pf:linkG:gammaunit}).

\emph{Step 2: the chain data.} For weights $\kappa$ let $\mathcal D_\kappa$ be the chain data tied to the frame with weights $\kappa$ and
plane degrees $d_{\mathrm{Pl}}$, with $\gamma=\gamma_{\mathrm{Pl}}$ (admissible by \cref{pf:linkG:gammaunit}), $c=c'\ge1$ and slopes
$\frac38s_U$; they satisfy the degree condition of the frame by construction. The link thresholds are positive:
$\theta_L(j,m)\ge\frac12B_j(x)-\frac38s_U(j,u_m)>0$ by \cref{pf:linkG:sizes}(b), as the tie threshold is at least $\frac12B_j(x)$. The
properties of joint projectors involve only the levels and the types of the letters, which do not depend on $\kappa$; so we may first
choose $e$ and $\kappa$ and then use $\mathcal D=\mathcal D_\kappa$.

\emph{Step 3: the projector.} \Cref{pf:linkG:jointproj} applied to $y$ gives $\kappa$ and $e$ with $v_B(ey)\le v_B(y)\le
v_p(y_{\beta'_0})+\frac12v_p(\beta'_0!)=\frac12v_p(\beta'_0!)$. Thus $q_1:=v_B(ey)$ is finite; let $\nu$ be a monomial with
$q_1=q_2+\frac12v_p(\nu!)$, $q_2=v_p((ey)_\nu)$. It has multidegree $d_{\mathrm{No}}$, as $ey\in\mathcal P_K(d_{\mathrm{No}})$.

\emph{Step 4: the dual element.} Let $\Pi=\sum_\beta\bigl(e(\mathrm{lev}(G)_\beta)\bigr)_\nu X^\beta$. Its coefficient at $\beta_0$ is $(ey)_\nu$, so
$q_3:=v_B(\Pi)\le q_2+\frac12v_p(\beta_0!)$; let $\beta_3$ attain $q_3$ and $z=\Pi_{\beta_3}\varpi^{v_p(\beta_3!)}$, so that $v_p(z)=q_3$ and
$F_2=z^{-1}\Pi$ has $v_B(F_2)=0$. By $\mathsf{Mem}$ for $G$ (applied to $\mathcal D$, $\rho$, $d_{\mathrm{No}}$, $e$ and $\nu$), $\Pi$ lies in the
$X$-space of $\mathcal D$, and so does its multiple $F_2$ (the conditions defining the $X$-space are invariant under non-zero scalars).

\emph{Step 5: the norm test.} \Cref{pf:linkG:normtest} gives a norm test $N_0$ for $\kappa$ and $c_Y$ with
$v_p(c_Y)=\frac12\nu(d_{\mathrm{No}})+D_{\mathrm{No}}-\frac12v_p(\nu!)$, and $N_0\equiv N_{\mathrm{ex}}:=\gamma_{\mathrm{No}}^{-1}\star(c_Ye(Y^\nu))$
modulo $p^{c'}$, coefficientwise. Put $k_0=c(\rho,d_{\mathrm{Pl}},d_{\mathrm{No}})\,c_Y\,\nu!\,z$.

\emph{Step 6: the valuation of $k_0$.} Write $c=c(\rho,d_{\mathrm{Pl}},d_{\mathrm{No}})$. By Steps 3--5,
\begin{align*}
v_p(k_0)&=v_p(c)+\tfrac12\nu(d_{\mathrm{No}})+D_{\mathrm{No}}-\tfrac12v_p(\nu!)+v_p(\nu!)+q_3\\
&\le v_p(c)+\tfrac12\nu(d_{\mathrm{No}})+D_{\mathrm{No}}+q_1+\tfrac12v_p(\beta_0!)\\
&\le v_p(c)+\tfrac12\nu(d_{\mathrm{No}})+D_{\mathrm{No}}+\tfrac12v_p(\beta'_0!)+\tfrac12v_p(\beta_0!) .
\end{align*}
By Step 1 and \cref{pf:linkG:multinomial}, $v_p(c)=v_p(\mathbf m_0!)-v_p(\beta_0!)-v_p(\beta'_0!)\le\nu_2(d)-v_p(\beta_0!)-v_p(\beta'_0!)$. Hence,
with $n_0=\frac12\nu(d_{\mathrm{Pl}})$ (the plane degrees of $\mathcal D$ are $d_{\mathrm{Pl}}$),
\[
v_p(k_0)\le\nu_2(d)-n_0+\Bigl(D_{\mathrm{No}}+\tfrac12\bigl(\nu(d_{\mathrm{Pl}})-v_p(\beta_0!)\bigr)+\tfrac12\bigl(\nu(d_{\mathrm{No}})-v_p(\beta'_0!)\bigr)\Bigr)
\le\nu_2(d)-n_0+9L
\]
by \cref{pf:linkG:sizes}(a).

\emph{Step 7: the value.} Let $A\in\mathcal P_r(d_{\mathrm{Pl}})$ be integral. Since $e(Y^\nu)\in\mathcal P_K(d_{\mathrm{No}})$, \cref{pf:linkG:digits}(a),
the definition of $N_{\mathrm{ex}}$ and the self-adjointness of $e$ give
\begin{align*}
\langle G,\mathrm{asm}_\rho(A,N_{\mathrm{ex}})\rangle&=c\,c_Y\sum_\beta\gamma_{\mathrm{Pl}}(\beta)A_\beta\beta!\sum_{\beta'}\mathrm{lev}(G)_{\beta,\beta'}\,
e(Y^\nu)_{\beta'}\,\beta'!\\
&=c\,c_Y\sum_\beta\gamma_{\mathrm{Pl}}(\beta)A_\beta\beta!\,\bigl\langle e(\mathrm{lev}(G)_\beta),Y^\nu\bigr\rangle
=c\,c_Y\,\nu!\,\langle\Pi,\gamma_{\mathrm{Pl}}\star A\rangle\\
&=k_0\,z^{-1}\langle\Pi,\gamma\star A\rangle=k_0\,B(F_2,A).
\end{align*}
As $\mathrm{asm}_\rho(A,\cdot)$ is linear, $\langle G,\mathrm{asm}_\rho(A,N_0)\rangle-k_0B(F_2,A)=\langle G,\mathrm{asm}_\rho(A,N_0-N_{\mathrm{ex}})\rangle$,
and the coefficients of $\mathrm{asm}_\rho(A,N_0-N_{\mathrm{ex}})$ are products of coefficients of $A$ and of $N_0-N_{\mathrm{ex}}$, of
valuation $\ge c'$. \Cref{pf:linkG:multinomial} gives the claimed bound. Finally, the data are those required in part (i) of the
proof of \cref{pf:linkG:prop5}, with $C_{Dn}=9$ and $C_1=4$.
\end{proof}

\begin{theorem}[Membership of assembled tests]\label{pf:linkG:memb}
For $C$ beyond a constant the membership part of \cref{pf:linkG:prop5} holds: let $\mathcal D$ be chain data tied to the frame with
weights $\kappa$ and satisfying the degree condition of the frame, $(\rho,d_{\mathrm{Pl}},d_{\mathrm{No}})$ block data, $N_0$ a norm test for
$\kappa$, and $G_p$ an integral polynomial in the model variables which lies in the $Y$-space of $\mathcal D$ with the constants
$(24,60)$ and whose monomials are mapped into $\mathcal P_r(d_{\mathrm{Pl}})$ by the substitution $\varphi$ of the model variables. Then
$\mathrm{asm}_\rho(\varphi(G_p),N_0)$ is a sum of one-sided elements of multidegree $d$.
\formallabel{Formalism.MainNotions.Proposition5Numerics.restMemb}
\formallabel{Formalism.MainNotions.Proposition5Numerics.membership_of_three_parts}
\formallabel{Formalism.MainNotions.Proposition5Numerics.membership_of_four_parts}
\formallabel{Formalism.MainNotions.Proposition5Numerics.restXXG_b}
\end{theorem}

\begin{proof}
Recall that a one-sided element of order $o$ is the image in the tangent ring of a polynomial
$X^{\mathbf l}\sum_i\prod_{j<J}\phi_j(\xi_{i,j})$, where $\mathbf l$ is a \emph{label} (exponents $<p^{\mathrm{dist}(j,x)}$ at the non-wrapped letters,
supported on $TT$), $\phi_j$ substitutes $X_{x,c}\mapsto X_{x,c}^{p^{\mathrm{dist}(j,x)}}$ for the letters of the segment $j$, each $\xi_{i,j}$ lies in the
segment space $X'_j$ (spanned by products of a plane and a normal polynomial on the letters of the segment, with the tie and threshold
conditions of the one-sided lemma) and is homogeneous with degrees $\mathrm{dd}_j(x)\le(1+\frac1{100})\mathrm{levdeg}_j(x)$ at the
wrapped and the late letters, and all monomials have normal degree $o$. For $C$ beyond a constant, the data satisfy the hypotheses
used below: $p$ is prime, $r\ge2$, the frame is valid for $TT$, $d$ is the flat degree on $TT$ and $0$ elsewhere, and the slopes at the
non-wrapped letters are at least $128c'$ (\lref{Formalism.MainNotions.Proposition5Numerics.mfh_of}). The proof has the
following steps.
\begin{enumerate}[label=(\arabic*)]
\item \emph{Integral models.} $G_p$ and $N_0$ have integral coefficients, so they are images of polynomials $G_o$, $N_o$ over $\OO$
(\lref{Formalism.MainNotions.Proposition5Numerics.exists_map_of_coefGe}).
\item \emph{Relabelling.} Let $\xi_{\mathrm{Pl}}(G_o)$ and $\xi_{\mathrm{No}}(N_o)$ be the polynomials obtained by renaming every model
variable, resp.\ letter variable, by the pair (segment, position) of its letter (a wrapped frame letter of level $j$ is placed in the
segment $j+1$). Then $\mathrm{asm}_\rho(\varphi(G_p),N_0)$ is the image of $X^{\mathbf l(\rho)}\Phi_{\mathrm{Pl}}(\xi_{\mathrm{Pl}}(G_o))
\Phi_{\mathrm{No}}(\xi_{\mathrm{No}}(N_o))$, where $\mathbf l(\rho)$ is the label of the digits and $\Phi_{\mathrm{Pl}},\Phi_{\mathrm{No}}$ apply
$\phi_j$ in each segment (\lref{Formalism.MainNotions.Proposition5Numerics.restAlgG}: the exponent $P_x\cdot P(j,m)$ of a wrapped
variable equals $p^{\mathrm{dist}(j+1,\cdot)}$, as $g(j,m)$ is the shift of the frame).
\item \emph{Valuation bounds.} With $\tau(j,u)=\theta^*_\kappa(j,u)-\frac12s_U(j,u)-60c'$: the valuation of $\xi_{\mathrm{Pl}}(G_o)$ on the
positions of the segment $j$ below the $u$-th letter is at least $\tau(j,u)$, and on the positions up to the $u$-th late letter at least
$\tau(j,u+1)$; the lowest weights of $\xi_{\mathrm{No}}(N_o)$ on the positions of the segment $j$ below, resp.\ up to, the $u$-th
letter are at most $\kappa(j,u)_0+4c'$, resp.\ $\kappa(j,u+1)_0+4c'$; and the supports lie on the letters of the segments, with the normal parts off the wrapped letters and of degrees
$d_{\mathrm{No}}$ (\lref{Formalism.MainNotions.Proposition5Numerics.restValBdG}). Indeed, these sets of positions are the images of
the chain sets ($u\le n_L(j)$) and of the link sets ($u=u_m$, $m+2\le r$) of $\mathcal D$, and the $Y$-space bounds give
$\theta(j,i)-s(j,i)-24c=\theta^*_\kappa-\frac18s_U-\frac38s_U-24c'\ge\tau$ on chains and $\theta_L(j,m)-60c=\theta^*_\kappa-\frac38s_U-60c'\ge\tau$ on
links (as $s_U\ge0$); the lowest weight bounds are those of the norm test.
\item \emph{Decomposition.} For every multidegree $\delta$ occurring in $\xi_{\mathrm{Pl}}(G_o)$, the homogeneous component
$\xi_\delta$ of multidegree $\delta$ and $\xi_{\mathrm{No}}(N_o)$ are sums of products over the segments,
$\xi_\delta=\sum_i\prod_jA_{i,j}$ and $\xi_{\mathrm{No}}(N_o)=\sum_{i'}\prod_jB_{i',j}$, where $A_{i,j}$ and $B_{i',j}$ live on the letters of the
segment $j$ and inherit the bounds of (3) (\lref{Formalism.MainNotions.Proposition5Numerics.decomposition_of_three_parts}: the
bounds of (3) concern sets of positions inside one segment, and slicing a polynomial along the segments, as well as taking homogeneous
components, preserves valuations and lowest weights on such sets). Expanding,
$\Phi_{\mathrm{Pl}}(\xi_\delta)\Phi_{\mathrm{No}}(\xi_{\mathrm{No}}(N_o))=\sum_{i,i'}\prod_j\phi_j(A_{i,j}B_{i',j})$.
\item \emph{Segment spaces.} Each product $A_{i,j}B_{i',j}$ lies in $X'_j$ (with the tie constant $m_*$, the thresholds
$\frac12B_j(x)-s_U$ and the order bound $\delta+d_{\mathrm{No}}$) and is homogeneous of the degrees $\delta(j,x)+d_{\mathrm{No}}(j,x)$ at the
letters $x$. Indeed, write $A=A_{i,j}$, $B=B_{i',j}$, and let $x$ be a non-wrapped letter of the segment, the $u$-th letter of level
$j$; its slope satisfies $s_U(j,u)\ge128c'$. If the position of $x$ is even, put $M_x=\kappa(j,u+1)_0+4c'$. The lowest weight of $B$ on the
positions up to $x$ is at most $M_x$ by (3); and as $\theta^*_\kappa(j,u)\ge\kappa(j,u+1)_0+m_*+s_U(j,u)$,
\[
\tau(j,u)\ge\kappa(j,u+1)_0+m_*+\tfrac12s_U(j,u)-60c'\ge\kappa(j,u+1)_0+m_*+4c'=M_x+m_*,
\]
so the valuation of $A$ on the positions below $x$ is at least $M_x+m_*$. If the position of $x$ is odd, then $x$ is a late letter (the
frame positions are even), $u<n_L(j)$, and with $M_x=\kappa(j,u)_0+4c'$ the lowest weight of $B$ on the positions below $x$ is at most
$M_x$, while $\theta^*_\kappa(j,u+1)\ge\kappa(j,u)_0+m_*+s_U(j,u+1)$ and $s_U(j,u+1)\ge128c'$ (the $(u+1)$-th letter is again a non-wrapped
letter of the segment) give, in the same way, that the valuation of $A$
on the positions up to $x$ is at least $\tau(j,u+1)\ge M_x+m_*$. Finally $\theta^*_\kappa\ge\frac12B_j(x)$ gives
$\tau(j,u)\ge\frac12B_j(x)-\frac12s_U(j,u)-60c'\ge\frac12B_j(x)-s_U(j,u)$, the threshold condition. These are the conditions defining
$X'_j$; the degrees of $A$ and $B$ give the homogeneity and the order bound.
\item \emph{Degrees.} $\delta(j,x)+d_{\mathrm{No}}(j,x)\le(1+\frac1{100})\mathrm{levdeg}_j(x)$ at the wrapped and the late letters
(\lref{Formalism.MainNotions.Proposition5Numerics.restDdG_holds}: the level degree bound of the frame for the block degrees).
\item \emph{Labels and multidegrees.} The label $\mathbf l(\rho)$ is a label supported on $TT$, since $\rho(x,c)<P_x$
(\lref{Formalism.MainNotions.Proposition5Numerics.restLabG}); and each piece
$X^{\mathbf l(\rho)}\Phi_{\mathrm{Pl}}(\xi_\delta)\Phi_{\mathrm{No}}(\xi_{\mathrm{No}}(N_o))$ has multidegree $d$, by the block condition, and all
its monomials have the same normal degree $o_\delta$ (\lref{Formalism.MainNotions.Proposition5Numerics.restDegG}).
\end{enumerate}
By (4)--(7), the image of $X^{\mathbf l(\rho)}\Phi_{\mathrm{Pl}}(\xi_\delta)\Phi_{\mathrm{No}}(\xi_{\mathrm{No}}(N_o))$ is a one-sided element of order
$o_\delta$ and multidegree $d$ (\lref{Formalism.MainNotions.Proposition5Numerics.membAt_core}), and summing over the occurring $\delta$
(the sum of the homogeneous components is $\xi_{\mathrm{Pl}}(G_o)$) gives $\mathrm{asm}_\rho(\varphi(G_p),N_0)$ by (2)
(\lref{Formalism.MainNotions.Proposition5Numerics.membAt_of}).
\end{proof}


% ---------------------------------------------------------------------------
\subsection{The components of the family statement}\label{sec:pf:linkG:famcomp}

We prove the six components used in the proof of \cref{pf:linkG:stmtEF}. Fix a family member with its scale
$(a,b,b',x,\eps,X',X''=X''(n),w''=w''(n),p_0=p_0(k))$, its parameter instance (offsets, base points of the orbits, the
exponents $e_p$ and the defects $\kappa_I$), link constants $K$, and top data $(a_t,b_t)_{t\in TT}$ in the window
($a\le a_t\le a+a/(kx^{A/2})$, $b'\le b_t$). We recall the objects of the formal statement.

\begin{definition}[Fillings, class sums and gains]\label{pf:linkG:famdefs}
Places are pairs $z=(\theta,t)$ with $\theta\in\{0,1\}$ and $t\in H$; the place is a \emph{top} if $\theta=0$ and $t\in TT$, \emph{thick}
if $\theta=1$ and $t$ is a thick position, and \emph{thin} otherwise. For a lift $j\in\{1,\dots,m\}$ the \emph{filling} is
$\mathrm{fill}_j(z)=[z\text{ top}]\,b_t+\mathrm{off}_j(z)$, where the offsets satisfy $b'/x\le\mathrm{off}_j\le b'/x+\eta$ at tops,
$w''\le\mathrm{off}_j\le w''+\eta$ at thin places and $b\le\mathrm{off}_j\le b+\eta+1$ at thick places ($\eta>0$ a size parameter);
the shifted filling is $\mathrm{fill}^G_j=\mathrm{fill}_j-\kappa_j$ with $\kappa_j=(k+(2j-1)(t+1))/2$. For a prime $p$ with
Frobenius element $h_p$, the \emph{partner} of a top $s$ at distance $e$ is $\pi_p(s,e)=([p\text{ signal and }e\text{ odd}],s+eh_p)$.
The \emph{class sums} are
\[
\mathrm{ask}^P_j(p,s)=\sum_{e\text{ odd}}\max\Bigl(0,\frac{a_s}{p^e}-\mathrm{fill}_j(\pi_p(s,e))\Bigr),\qquad
\mathrm{ask}^G_j(p,s)=\sum_{e\text{ odd}}\max\Bigl(0,\frac{a_s}{p^e}-\mathrm{fill}^G_j(\pi_p(s,e))\Bigr).
\]
Let $c_G=4S_c$ (a size constant), $N(m)$ the integer of the datum with $N(m)\le c_G$, and
\[
\mathrm{lo}(C,p)=C(c_G+\eta)\log_p(2a)+\mathrm{pr}(p)+\mathrm{or}(p),\qquad
\mathrm{fl}(C,p)=Cf\,N(m)^2\log_pa+Cf\log_p^2a,
\]
where $\mathrm{pr}(p),\mathrm{or}(p)\ge0$ are the explicit residue-class terms of the formal definition (they depend on $C$
monotonically) and $f$ is the degree of the datum. For a good orbit representative $t_0$ with set of tops $O(t_0)$ (of size
$N_v(t_0)$), the \emph{orbit gain} $g_U$, the \emph{costs} $\Pr$, $t_v$, $D_p$ and the \emph{loss}
$\mathrm{dl}(t_0)=\mathrm{dl}_1n_PN(m)^2\log_p(2a)+\mathrm{dl}_2N_v(t_0)\log_p^2(2a)$ are those of the formal statement; the
\emph{served} orbits $VV(p)$, the totals $\tau(p)=e_p+m\sum_{t_0\in VV(p)}(\lceil t_v(t_0)\rceil+D_p+1)$ (or $0$ if $VV(p)=\emptyset$)
and the gains $\mathrm{gain}_P(p,j)$ are as in \cref{pf:linkG:stmtF}. The \emph{star bound} at a candidate prime $p$, a good
representative $t_0$ and a lift $j$ is
\begin{equation}\label{eq:pf:linkG:star}
\sum_{s\in O(t_0)}\bigl(\mathrm{ask}^P_j(p,s)-\mathrm{lo}(C,p)\bigr)\le g_U(t_0)+\mathrm{dl}(t_0)-\Pr(t_0)-\lceil t_v(t_0)\rceil-D_p-1 .
\end{equation}
\formal{Formalism.MainNotions.G.askP}
\formal{Formalism.MainNotions.G.askG}
\formal{Formalism.MainNotions.G.lo}
\formal{Formalism.MainNotions.G.fl}
\formal{Formalism.MainNotions.G.gU}
\formal{Formalism.MainNotions.G.gainP}
\formal{Formalism.MainNotions.G.tauP}
\formal{Formalism.MainNotions.AuxBlock8.Pf.StarAt}
\end{definition}

\begin{lemma}[Comparison of the fillings]\label{pf:linkG:asks}
If $N(m)+1\le w''$, then for all top data with $a_t\le2a$, $b_t\ge0$, all primes $p$, lifts $j$ and tops $s$,
\[
\mathrm{ask}^P_j(p,s)\le\mathrm{ask}^G_j(p,s)\le\mathrm{ask}^P_j(p,s)+N(m)\log_p(2a).
\]
\formallabel{Formalism.MainNotions.AuxBlock8.Pf.As.asksAll_of}
\formallabel{Formalism.MainNotions.AuxBlock8.Pf.As.shift_le}
\formallabel{Formalism.MainNotions.AuxBlock8.Pf.As.shift_ge}
\formallabel{Formalism.MainNotions.AuxBlock8.Pf.As.kaJ_bounds}
\formallabel{Formalism.MainNotions.AuxBlock8.Pf.As.fill_ge}
\formallabel{Formalism.MainNotions.AuxBlock8.Pf.rest0}
\end{lemma}

\begin{proof}
Every offset is at least $w''$ (thick offsets are at least $b\ge w''$ and top offsets at least $b'/x\ge w''$), so
$\mathrm{fill}_j\ge w''$. For $1\le j\le m$, $0\le\kappa_j\le N(m)$. Put $y=a_s\in[0,2a]$ and $w_e=\mathrm{fill}^G_j(\pi_p(s,e))=
\mathrm{fill}_j-\kappa_j\ge w''-N(m)\ge1$. Then $\mathrm{ask}^P$ is the sum of $\max(0,y/p^e-(w_e+\kappa_j))$ and $\mathrm{ask}^G$ that of
$\max(0,y/p^e-w_e)$ over odd $e$; the first inequality is termwise. For the second, termwise
$\max(0,u)\le\max(0,u-\kappa_j)+\kappa_j[u>0]$, and $y/p^e-w_e>0$ forces $p^e<y\le2a$ because $w_e\ge1$. The odd $e\ge1$ with
$p^e<2a$ are at most $\log_p(2a)$ in number (none if $2a\le p$), so the second inequality holds with $\kappa_j\log_p(2a)\le
N(m)\log_p(2a)$. The size condition $N(m)+1\le w''$ is one of the admissibility conditions of the family, so the comparison
holds eventually for every member.
\end{proof}

\begin{lemma}[Raising the constant]\label{pf:linkG:addone}
Assume $N(m)\le c_G$ and the conclusions of \cref{pf:linkG:asks}. If the inequality \eqref{eq:pf:linkG:conc1} holds at all
primes $p\le X''$ of the inert or signal class with $\mathrm{ask}^P$ in place of $\mathrm{ask}^G$ and the constant $C$, then it holds
with $\mathrm{ask}^G$ and the constant $C+1$.
\formallabel{Formalism.MainNotions.AuxBlock8.Pf.As.gainBoundAll_add_one_of_gainBoundAllRequired}
\formallabel{Formalism.MainNotions.AuxBlock8.Pf.As.lo_succ}
\formallabel{Formalism.MainNotions.AuxBlock8.Pf.As.fl_succ}
\end{lemma}

\begin{proof}
The first term of $\mathrm{lo}$ grows by $(c_G+\eta)\log_p(2a)$ when $C$ is replaced by $C+1$, and the other terms of $\mathrm{lo}$
and $\mathrm{fl}$ are non-decreasing in $C$ (they are non-negative multiples of $C$ plus terms independent of $C$). Hence
$\mathrm{ask}^G_j(p,s)-\mathrm{lo}(C+1,p)\le\mathrm{ask}^P_j(p,s)+N(m)\log_p(2a)-\mathrm{lo}(C,p)-(c_G+\eta)\log_p(2a)\le
\mathrm{ask}^P_j(p,s)-\mathrm{lo}(C,p)$ and $m\,\mathrm{fl}(C+1,p)\ge m\,\mathrm{fl}(C,p)$; the left side of \eqref{eq:pf:linkG:conc1}
can only decrease.
\end{proof}

\begin{lemma}[The losses]\label{pf:linkG:losses}
Let $2\mathrm{dl}_1+C_e\le C$ and $4\mathrm{dl}_2\le C$. If $e_p\le C_efk^2\log_pa$, then
$\sum_{t_0}\mathrm{dl}(t_0)+e_p\le\mathrm{fl}(C,p)$, the sum over all orbit representatives at $p$.
\formallabel{Formalism.MainNotions.AuxBlock8.Pf.sum_dl_add_ep_le}
\formallabel{Formalism.MainNotions.AuxBlock8.Pf.sum_dl_le}
\formallabel{Formalism.MainNotions.AuxBlock8.Pf.card_reps_mul_nP_le}
\end{lemma}

\begin{proof}
The orbits of the translation by the step at $p$ partition $H$, which has $2n=f$ elements; hence (number of
representatives)$\cdot n_P\le f$ (for signal primes the step is $2h_p$ and $n_P$ is twice its order) and $\sum_{t_0}N_v(t_0)\le|TT|\le
f$. As $a\ge2$, $2a\le a^2$ and $\log_p(2a)\le2\log_pa$; hence
$\sum_{t_0}\mathrm{dl}(t_0)\le2\mathrm{dl}_1fN(m)^2\log_pa+4\mathrm{dl}_2f\log_p^2a$. Adding $e_p\le C_efk^2\log_pa\le C_efN(m)^2\log_pa$
(as $k\le N(m)$) and using $2\mathrm{dl}_1+C_e\le C$, $4\mathrm{dl}_2\le C$ gives the claim.
\end{proof}

\begin{proposition}[Gain bounds from the star bounds]\label{pf:linkG:gainstar}
Let $C\ge3$, $2\mathrm{dl}_1+C_e+1\le C$, $4\mathrm{dl}_2\le C$, and assume the unit bound $e_p\le(C_e+1)fk^2\log_pa$ for $p\le a$
(condition $\sfname{Adm9F}$). If the star bound \eqref{eq:pf:linkG:star} holds for all candidate primes, good representatives and
lifts, then for every prime $p\le X''$ of the inert or signal class and every lift $j$: $0\le\mathrm{gain}_P(p,j)\le4|TT|a/p$, and
\eqref{eq:pf:linkG:conc1} holds with $\mathrm{ask}^P$ in place of $\mathrm{ask}^G$ and the constant $C$.
\formallabel{Formalism.MainNotions.AuxBlock8.Pf.gain_bound_all_of_star}
\formallabel{Formalism.MainNotions.AuxBlock8.Pf.gain_bound_at_of_star}
\formallabel{Formalism.MainNotions.AuxBlock8.Pf.noop}
\formallabel{Formalism.MainNotions.AuxBlock8.Pf.lo_ge_cp}
\end{proposition}

\begin{proof}
The bounds for $\mathrm{gain}_P$ are \lref{Formalism.MainNotions.AuxBlock8.Pf.gain_bounds} (a lemma of the top-shape section:
the gain is a sum of non-negative terms bounded by $4|TT|a/p$). Write $\Sigma=\sum_j\sum_{s\in TT_P(p)}(\mathrm{ask}^P_j(p,s)-\mathrm{lo}(C,p))$.
Note $\mathrm{lo}(C,p)\ge3c_G\log_p(2a)$, since $C\ge3$ and the other terms are non-negative.

\emph{No served orbit.} Then $\tau(p)=0$. If $p$ is neither signal nor $\ge X'$, $TT_P(p)=\emptyset$ and the left side is
$-m\,\mathrm{fl}\le0\le\sum_j\mathrm{gain}_P(p,j)$. If $p\ge X'$ is in the inert class with $h_p=0$, the partner of a top $s$ is $s$ itself, whose filling
$b_s+\mathrm{off}\ge b'+b'/x$ exceeds $a_s/p^e$ for $e\ge1$ (as $p\ge X'$), so all class sums vanish and the left side is $\le0$. Otherwise we compare with the non-served part $g_0$ of the gain: for an odd $e$ which is not excluded,
$a_s/p^e-\mathrm{fill}_j\le(a_s-2c_G)/p^e-(\mathrm{fill}^G_j+c_G)+3c_G$, the excluded distances contribute nothing (their partner is
another top with a large filling), and at most $\log_p(2a)$ odd $e$ contribute; hence
$\sum_{s\in O(t_0)}(\mathrm{ask}^P_j(p,s)-\mathrm{lo})\le g_0(t_0)+\mathrm{dl}(t_0)$ for every orbit, and summing with \cref{pf:linkG:losses}
gives $\Sigma-m\,\mathrm{fl}\le\sum_j\sum_{t_0}\max(g_0(t_0),0)\le\sum_j\mathrm{gain}_P(p,j)$.

\emph{Served orbits.} Then $p$ is a candidate, $VV(p)$ is the set of good representatives, $TT_P(p)$ is the union of the orbits
$O(t_0)$, $t_0\in VV(p)$, and $\tau(p)=e_p+m\sum_{t_0\in VV(p)}(\lceil t_v\rceil+D_p+1)$. Summing \eqref{eq:pf:linkG:star} over
$t_0\in VV(p)$ and the $m$ lifts,
\[
\Sigma\le\sum_j\sum_{t_0\in VV(p)}(g_U-\Pr)(t_0)+m\sum_{t_0\in VV(p)}\bigl(\mathrm{dl}(t_0)-\lceil t_v(t_0)\rceil-D_p-1\bigr).
\]
By definition $\mathrm{gain}_P(p,j)\ge\sum_{t_0\in VV(p)}(g_U-\Pr)(t_0)$, and by \cref{pf:linkG:losses} (with $C_e+1$ in place of $C_e$)
and $m\ge1$, $e_p+m\sum_{t_0}\mathrm{dl}(t_0)\le m\,\mathrm{fl}(C,p)$. Therefore $\Sigma-m\,\mathrm{fl}(C,p)\le\sum_j\mathrm{gain}_P(p,j)-\tau(p)$.
\end{proof}

\begin{theorem}[Star bounds]\label{pf:linkG:starbounds}
There is $C_0\ge0$ such that for all $C\ge C_0$, eventually in $x$, the star bound \eqref{eq:pf:linkG:star} holds for every member,
all top data in the window, all candidate primes $p$, all good representatives $t_0$ and all lifts $j$.
\formallabel{Formalism.MainNotions.AuxBlock8.Pf.restStar}
\formallabel{Formalism.MainNotions.AuxBlock8.Pf.St.star_fam}
\formallabel{Formalism.MainNotions.AuxBlock8.Pf.St.starAll_of}
\formallabel{Formalism.MainNotions.AuxBlock8.Pf.St.gU_lower}
\formallabel{Formalism.MainNotions.AuxBlock8.Pf.St.kink}
\end{theorem}

\begin{proof}
Fix the data and write $\ell=\log_p(2a)\ge1$. Three estimates are combined.

\emph{(a) Lower bound for the orbit gain.} With $0\le\gamma_p\le1$,
\[
\sum_{s\in O(t_0)}\mathrm{ask}^P_j(p,s)+\sum_{s\in O(t_0)\text{ vector}}M_j(s)-N_v(t_0)\cdot3c_G\ell\le g_U(t_0)+\mathrm{dl}(t_0),
\]
where for a vector top $s$, $M_j(s)=\sum_{e\text{ odd}}\min(\mathrm{al}(s)/p^e,\mathrm{cap}_j(s,e))$ with $\mathrm{al}(s)=a_s-2c_G$ and
$\mathrm{cap}_j=\mathrm{fill}^G_j\circ\pi_p+c_G$. Indeed, for each top and odd $e$,
$\max(0,a_s/p^e-\mathrm{fill}_j)+\min(\mathrm{al}(s)/p^e,\mathrm{cap}_j)\le\mathrm{al}(s)/p^e+3c_G$ for the at most $\ell$ odd $e$ with
$p^e<2a$, while the geometric series gives $\sum_{e\text{ odd}}\mathrm{al}(s)/p^e\le\gamma_p\mathrm{al}(s)$ with
$\gamma_p=p/(p^2-1)$; summing over the orbit gives the claim from the definition of $g_U$.

\emph{(b) The kink inequality.} On an orbit of even length $n'$ with tops $TT$ split into vector and non-vector positions, for
$q\ge2$ and non-negative $a,W$:
\[
\sum_{y\text{ non-vector}}\min\Bigl(W(y),\frac{a(y-d(y))}{q^{d(y)}}\Bigr)\le\sum_{x\text{ vector}}\sum_{e\text{ odd}}\min\Bigl(\frac{a(x)}{q^e},W(x+e)\Bigr),
\]
where $d(y)$ is the distance from $y$ back to the nearest vector top (terms with no such top omitted): the map
$y\mapsto(y-d(y),d(y))$ is injective and $d(y)$ is odd, because vector and non-vector tops lie at positions of different
parity. For the orbit this gives
(\lref{Formalism.MainNotions.AuxBlock8.Pf.St.kinkHH}: the kink inequality with $q=p$, combined with the bounds for the offsets
of tops): the parts of $\Pr(t_0)$ and of $t_v(t_0)$ coming from the non-vector tops are at most
$\sum_{s\text{ vector}}M_j(s)+N_v(t_0)(b'/x+\eta)$.

\emph{(c) Arithmetic.} The remaining parts of $\Pr(t_0)+\lceil t_v(t_0)\rceil+D_p+1$ are explicit expressions in the size
parameters of the scale, in $N_v(t_0)$ and in the number of positions of the orbit, which is at most $N_v(t_0)$ times a size
parameter $G$ by the orbit counts of $\sfname{Adm1F}$ (and $N_v(t_0)\ge1$). Two inequalities between real numbers, one for
$p\ge p_0$ and one for $p<p_0$ (\lref{Formalism.MainNotions.AuxBlock8.Pf.St.caseE_arith} and \lref{Formalism.MainNotions.AuxBlock8.Pf.St.caseD_arith}:
each states that an explicit linear combination of the size parameters is at most $N_v(t_0)$ times the corresponding part of
$\mathrm{lo}(C,p)$, under explicit hypotheses on the parameters), bound them by $N_v(t_0)(\mathrm{lo}(C,p)-3c_G\ell)$, once $C\ge C_0$
and the numerical facts of the family hold
(\lref{Formalism.MainNotions.AuxBlock8.Pf.St.numStar_ev} and \lref{Formalism.MainNotions.AuxBlock8.Pf.St.B.facts_ev}: finitely many
inequalities between explicit functions of $x$, each true for large $x$ because its two sides are of different orders in $x$).
Combining (a), (b), (c):
$\sum_{s\in O(t_0)}(\mathrm{ask}^P_j-\mathrm{lo})\le g_U+\mathrm{dl}-\sum_{\text{vector}}M_j+N_v\cdot3c_G\ell-N_v\mathrm{lo}\le
g_U+\mathrm{dl}-\Pr-\lceil t_v\rceil-D_p-1$, which is \eqref{eq:pf:linkG:star}.
\end{proof}

\begin{proposition}[Mean defect and weight vectors]\label{pf:linkG:meanw}
Eventually in $x$, under the hypotheses of \cref{pf:linkG:stmtEF}, every member satisfies the mean-$\kappa$ bound of part (b)
of $\mathsf{Gain}$ and the existence of weight vectors of part (c).
\formallabel{Formalism.MainNotions.AuxBlock8.Pf.rest2}
\formallabel{Formalism.MainNotions.AuxBlock8.Pf.restW}
\formallabel{Formalism.MainNotions.AuxBlock8.Pf.second_estimate_and_mean_along_family}
\end{proposition}

\begin{proof}
Apply the admissibility conditions with $\delta=\eps/48$. The numerical facts for the mean
(\lref{Formalism.MainNotions.AuxBlock8.Pf.numMean_ev}) and the type I hypotheses
(\lref{Formalism.MainNotions.AuxBlock8.Pf.restHypI_ev}) hold eventually; these are lemmas of the section on the first
theorem. With them, the slope estimate of lemma B (\lref{Formalism.MainNotions.AuxBlock8.Pf.slopeEstimate_of_lemmaB_slopeLine})
gives the mean-$\kappa$ bound from $\sfname{Adm8F}$, and the mean excess statement
(\lref{Formalism.MainNotions.AuxBlock8.Pf.meanEx_of}) together with
\lref{Formalism.MainNotions.AuxBlock8.Pf.weight_condition_of_mean_excess} gives the weight vectors. All four are shared
lemmas of the first theorem; the threshold in $x$ is the maximum of the thresholds involved.
\end{proof}

\begin{proposition}[Points lemma and weighted prime sums]\label{pf:linkG:points}
For every $C\ge0$, eventually in $x$, every member whose gain bounds hold with the constant $C$ satisfies the points lemma of
part (c): for every set $S$ of primes of the inert class,
\begin{gather*}
\sum_{p\in S}\log p\cdot\max\Bigl(0,\sum_{j,s}\mathrm{ask}^G_j(p,s)-\Bigl(\sum_j\mathrm{gain}_P(p,j)-\tau(p)-\mathrm{io}(p)\Bigr)\Bigr)\le
\frac{5m|TT|a}x,\\
\sum_{p\in S}\log p\cdot\mathrm{io}(p)\le m|TT|ax^{-9};
\end{gather*}
and the weighted prime sum bound: $\sum_p\log p\sum_{j,s}\mathrm{ask}^G_j(p,s)\le m|TT|ax^{-16}$ over primes $p>X''$ of the inert or
signal class, and $\sum_{p\le X''}\log p\cdot C(c_G+\eta)\log_p(2a)\le bx^{-100}$, $\sum_{p\le X''}\log p\cdot\mathrm{fl}(C,p)\le
bx^{-100}$.
\formallabel{Formalism.MainNotions.AuxBlock8.Pf.rest3}
\formallabel{Formalism.MainNotions.AuxBlock8.Pf.restT}
\end{proposition}

\begin{proof}
The primes are split into ranges: $p<p_0$, $p_0\le p<X'$ (candidates and band primes), $X'\le p\le X''$, and $p>X''$. In each
range the quantity under the sum is bounded by the gain bounds and \cref{pf:linkG:asks} by $m|TT|$ times an explicit error
term ($\mathrm{lo}$, $\mathrm{fl}$, the band and small-order contributions, or, for $p>X''$, $\theta_Ww''+S_c$ for
$p<X''(1+\theta_S)$ and $0$ beyond), and Chebyshev's bound $\sum_{p\le y}\log p\le y\log4$ \cite{MontgomeryVaughan2007} bounds the
weighted number of primes in each range. The resulting inequalities are the numerical facts along the family
(\lref{Formalism.MainNotions.AuxBlock8.Pf.G5.prime_sum_facts_along_family} and \lref{Formalism.MainNotions.AuxBlock8.Pf.G5.Nu.numT_ev}:
each range contributes at most $a/(40x)$, the candidate range at most $\frac32a/x$, and $x\ge120$), which hold eventually in
$x$. The bound for $\mathrm{io}(p)=2\,\mathrm{Sig}(p)\max_{t_0}S_v(t_0)$ uses $\mathrm{Sig}(p)\le x^{-\eps}$ at served primes and the
lattice chain estimate for $S_v$ (\lref{Formalism.MainNotions.AuxBlock8.Pf.G5.latticeChain_succ_of_partial}: summing over
the residue classes of the primes and the lattice chains of the served orbits).
\end{proof}

% ---------------------------------------------------------------------------
\subsection{Remarks on the formal statements}\label{sec:pf:linkG:remarks}

\begin{remark}\label{pf:linkG:rem}
\begin{enumerate}[label=(\arabic*)]
\item The local statement $\linkGLL(C)$ is formally a proposition $\mathrm{aux}_3(D,C)$ whose first argument, a datum
$D$, does not occur in it; likewise $\mathsf{Concl}(D)$ and all $37$ clauses have the datum as a dummy argument.
\Cref{thm:families:local} is stated formally as ``there is $C_0\ge0$ with $\mathrm{aux}_3(D,C_0)$ for all $D$''; if no
datum existed this would hold vacuously with $C_0=0$, and the proof indeed treats that case separately. The
unconditional form of the theorem needs one datum.
\item The formal statement of \cref{thm:families:local} has the hypotheses $(\mathrm E)$ and $(\mathrm L6)$; they are
discharged in the proof of \cref{thm:linkG} (\cref{pf:linkG:locmodel}). The proof of $(\mathrm E)$ does not use
$\sfname{Loc}$, although the formal lemma providing it takes $\sfname{Loc}$ as an argument.
\item In $\mathsf{Crit}$ the two auxiliary existence statements are hypotheses, but the proof of the
criterion (\cref{pf:linkG:loc}) does not use them: $\sfname{Adm}$ alone implies $\mathsf{Concl}$.
\item Several clauses carry names unrelated to their content ($\sfname{ArchimedeanSignatureMatching}$ is
$(\mathrm P)$, $\sfname{HerrCocycleExtensionClass}$ is $\sfname{WittIdent}$, $\sfname{FontaineFormalProperties}$ is
$\sfname{FlatTop}$, $\sfname{FewBadOrbits}$ is Weyl's dimension formula, $\sfname{LevelPrimesProviso}$ is the
multiplicativity of the tensor lattices, etc.); see \cref{tab:pf:linkG:clauses}.
\item In the proof of \cref{thm:linkG} the Kudla family statement is applied to the lattice instance whose chain maps
are $0$. For it the valuation conditions on chain images (the unit bound at the twisted places and the third part of
$\mathsf{Len}$) hold trivially, so they play no role in the proof of $\sfname{LinkG}$; only the length part of
$\mathsf{Len}$ is used.
\item The function $\delta$ of $\sfname{LinkG}$ is the one produced by the Kudla family statement, and the constant is
$4|C|+1$; the medium-prime defects $\kappa_I$ are defined as infima of admissible defects, so the medium-prime
statement $\sfname{Adm7F}$ holds by construction, and the content at medium primes is the mean bound $\sfname{Adm8F}$.
\item The hypothesis $\sfname{Adm2F}$ of \cref{thm:linkG} enters only through the Kudla family statement (the size
condition on $|TT|$ and the bound for $B_{TT}$).
\end{enumerate}
\end{remark}
